\chapter{Max Weber and Talcott Parsons}

\section{Max Weber: Disenchantment, Authority and Bureaucracy}

\subsection{Disenchantment (Revision of a Student Doubt)}
A student asked for disenchantment to be explained once again . The meaning is very simple: disenchantment means that you are no longer enchanted .

\begin{KeyConcept}
To be no longer enchanted means that you no longer believe in the magical, spiritual, traditional world .
Put plainly: you no longer believe in the concept of God; you no longer believe in the concept of souls and spirits; you no longer believe in the concept of caste; you no longer believe in the concept of the joint family .
All those traditional elements and values are over . You have become completely rational .
\end{KeyConcept}

\paragraph{The Candlelight Dinner Analogy}
Suppose you go on a date with a girl . It is a candlelight dinner in the open, near the beach in Goa . The girl says, ``Wow, what a beautiful night, such a starry night, so many stars, they are twinkling, that is so romantic'' . And then the boy says: this is due to refraction of light .

The next morning it is rainy and a beautiful rainbow appears . She asks for a photograph with the rainbow, and the boy says there is nothing special about it---it is total internal reflection of light .

When you start giving explanations in a fully rational manner to everything around you, you lose the magic, you lose the charm, you lose the beauty of the world . That is what we call a disenchanted man, woman or person .

\paragraph{The Bureaucrat Speaks the Same Way}
This is exactly how a bureaucrat speaks . You go to a bureaucrat in a full emotional state, and he speaks in a fully rational, law-bound manner ``according to CRPC code 345'' . It is the same reaction the girl receives when she is told it is total internal reflection of light .

\begin{CriticalPoint}
There is no difference between the disenchanted person and that bureaucrat . Both are dehumanized, both are completely rational .
This is what the world is going to become . This is what Max Weber said .
\end{CriticalPoint}

\begin{QuickRevision}
\begin{itemize}
    \item Disenchantment = no longer enchanted = no longer believing in the magical, spiritual, traditional world (God, soul, caste, joint family) .
    \item Analogy: stars explained as refraction; rainbow explained as total internal reflection---the charm dies .
    \item The disenchanted person and the rule-quoting bureaucrat are the same type: dehumanized, completely rational .
\end{itemize}
\end{QuickRevision}

\subsection{Authority and the Basis of Legitimacy}
In each form of authority, you have given legitimacy . You have given legitimacy to those authorities to rule over you . The form of legitimacy is very simple .

\paragraph{Why Do You Obey the Prime Minister?}
Today you have given legitimacy to the Prime Minister to rule over you . If asked what the basis of that legitimacy is, you will say the Constitution . But let us be very honest---it is not only the Constitution . It was also the charisma of that leader---the Modi wave . This is a fact . And let us be more honest: it was also tradition . Tradition has somehow played a very big role .

\begin{KeyConcept}
Reality is highly complex and cannot be captured through concepts . That is why we form only models .
Weberians would never say that this is what absolute reality is . They would say it is a model---largely a trend, not the absolute reality .
\end{KeyConcept}

\paragraph{The American Example: Trump versus Hillary Clinton}
Consider the election in the USA in which Trump fought against Hillary Clinton . People voted for Trump . What could have been the reason? 
\begin{enumerate}
    \item \textbf{Charisma}---the way he used to speak . He connected with the masses because he used to sound very ordinary, and so he spoke in a very simple fashion . This was the first thing that attracted people .
    \item \textbf{Restoring the great American pride}---he believed past leaders of America were responsible for the current state of the economy, that the economy was in doldrums, that America had been wasting funds on defence in Iran, Iraq and Afghanistan, and that this had to stop . In this he also sounded very rational .
    \item \textbf{Legal-rational element}---he belonged to the Republican Party . A political party is a legal-rational authority; it is not a religious organisation .
    \item \textbf{Gender}---the fourth factor was that he was male, and this cannot be ignored . Hillary Clinton could have become the first woman president; for so long no woman president has been elected . Somehow the role of patriarchy plays a role in elections .
\end{enumerate}

\begin{KeyConcept}
Patriarchal bargain: in order to exercise power, Hillary Clinton bargained with patriarchy---she exchanged power by carrying a patriarchal tone and attitude into the election and into her speeches .
\end{KeyConcept}

But largely, the Prime Minister rules us on the basis of the Constitution . Largely, this has been the trend . We can no longer say that India today is governed by Brahmins, which at one point of time was the case .

\begin{ExampleBox}
We used to have Brahmanical authority . Brahmins used to decide everything; even the kings used to get appointed by Brahmins . There was no democracy---people did not elect people; it was Brahmins who elected people to the throne .
\end{ExampleBox}

That too was legitimate authority . Legitimacy is now given to legal-rational authority on the basis of rules and regulations .

\paragraph{Why We Obey Charismatic Figures}
When it comes to heroic, charismatic, supernatural-quality-possessing leaders---Gandhi, Hitler, Mother Teresa, Jesus Christ, Prophet Muhammad, Lord Buddha, Guru Nanak---we are not obeying them or giving them legitimacy to command us on the basis of the Constitution . We are simply enthralled by the heroic and supernatural qualities those people possess .
\begin{itemize}
    \item They induce hope in a hopeless situation .
    \item They somehow give us salvation and pacify our concerns .
    \item They make us believe that this crisis period will be over very soon---``trust me, yes we can'' .
    \item People fall in love with those heroic qualities and start looking at that person as someone directly sent by God to the earth .
\end{itemize}

\begin{KeyConcept}
What is legitimacy? Legitimacy is justification---what actually justified their rule .
The reasons for justification have been different throughout history, but justification is justification; only the reasons differ .
\end{KeyConcept}

\paragraph{The Legitimacy of Traditional Rulers}
The reason for traditional authority to rule upon us has largely been legitimacy on the basis of ancient traditions or texts .
\begin{itemize}
    \item Why should there be the rule of the Queen in England? Today the Queen is simply a ceremonial, nominal head---not the executive head---but at one point of time she was the executive head . What gave legitimacy was ancient tradition and text: the monarch was sent directly to earth by God from heaven . Religious texts used to give legitimacy .
    \item In ancient India, Kanishka of the Kushan empire was called \textit{Devaputra}---son of God . That is the kind of title we used to give .
    \item The first invaders of India---Greeks, Shakas, Parthians, Kushans---were all traditional authorities . None of them ruled upon us on the basis of a constitution, nor purely on the basis of heroic qualities . Largely it was the ancient text or tradition that legitimised them .
    \item In the Greek texts one finds a connection established between the Greek and Indian civilisations---that is why we always say Indo-Greek . The claim was that our culture, our language, our civilisations have been quite similar, and this justified their rule .
\end{itemize}

\begin{ExampleBox}
The British in India: they first made us believe that we are uncivilised and not capable enough to govern ourselves . It was not that they directly conquered or annexed India .
They first took permission from Jahangir to establish a port . Why did the British obey Jahangir? Because Jahangir was a traditional authority, standing in the lineage of the Mughal emperors . That is why they did not challenge him---they obeyed him and took permission .
\end{ExampleBox}

\begin{QuickRevision}
\begin{itemize}
    \item Legitimacy = justification for rule; the reasons differ across history, the need for justification does not .
    \item The PM's legitimacy mixes constitution + charisma + tradition---reality is complex, so Weber offers only models and trends .
    \item Trump's victory read through four factors: charisma, restoring national pride, party (legal-rational), and gender/patriarchy; Hillary Clinton's ``patriarchal bargain'' .
    \item Charismatic figures are obeyed because they induce hope in hopeless situations; traditional rulers are obeyed on the strength of ancient texts and traditions (Devaputra, monarchy by divine sanction, Indo-Greek claims, Jahangir and the British) .
\end{itemize}
\end{QuickRevision}

\subsection{Ideal Types of Traditional Authority}
Max Weber has done a very rigorous analysis of each and every thing in the world . The ideal types of traditional authority are the following keywords, which are going to help you be different from other candidates :
\begin{enumerate}
    \item Gerontocracy 
    \item Primary patriarchalism 
    \item Patrimonialism, also called prebendalism 
    \item Sultanism---the extreme version of patrimonialism 
\end{enumerate}

\begin{CriticalPoint}
A student asked whether traditional and charismatic authority may be illegitimate for someone . Yes, definitely .
Even the rule of the present Prime Minister is considered illegitimate by many people---they believe EVM machines have been hacked, that the Supreme Court and High Court and other institutions have been demolished, that police and army officials were involved, and that the media has become subjugated to the whims and fancies of the leader .
But that does not change the fact that he is the Prime Minister .
\end{CriticalPoint}

\paragraph{Gerontocracy---Rule of the Elderly}
Gerontocracy simply means rule of the elderly . The root word \textit{geron} describes elderly people; \textit{-cracy} means rule, while \textit{-logy} means study . Gerontology as a discipline talks about care, health and death; gerontocracy is a form of authority .
\begin{itemize}
    \item Example: Sparta . Or any village community---the village panchayat, where the elderly, anyone considered above 60, are ruling .
    \item Plato once said that anyone below 50 should not be the ruler of the nation state---anyone below 50 should not be eligible to be the ruler .
    \item That means the factor to become ruler is age . The legitimacy you give to that person to rule upon you is age . The elderly could rule because we had given them legitimacy and we gave it because they were elderly .
    \item This happens in tribal societies, in village societies, in Sparta, in the Greek and Roman empires and even in modern civilisations .
    \item Have we not had the rule of the elderly only since independence? Only the elderly are ruling . We are seeing the reality, but we must also see the concept in the same way . This is just a model .
\end{itemize}

\begin{CriticalPoint}
Students of public administration, in the name of bureaucracy, know only IAS and IPS . Their understanding of bureaucracy is un-Weberian---they read only the surface, even when reading Weber's own model .
The way they read Max Weber and the way we read Max Weber are very different .
\end{CriticalPoint}

\paragraph{Primary Patriarchalism}
Primary patriarchalism basically operates at the level of the household or family . It is the rule of the father at home; the father is the head of the household .
\begin{itemize}
    \item Even in joint families we have rule of the elderly---gerontocracy---and most of those elderly would be men .
    \item The concepts mix, because reality itself is like that . We only take the help of models to get closer to reality, and then we find that in this model too, that reality is fitting in . Reality is highly vast and complex .
\end{itemize}

\paragraph{Patrimonialism---The Most Important Concept}
Among the traditional authorities, patrimonialism is considered the most important one . Think of it simply: the rule of the father at home has been extended . Now the home has been replaced by the country, and the father is the ruler . When patriarchy is extended to a larger state, that becomes patrimonialism .

\begin{KeyConcept}
Patrimonialism is the extended version of primary patriarchalism .
The only difference: primary patriarchalism as a concept operates in the household or family; patrimonialism is beyond the family .
\end{KeyConcept}

In patrimonial authority the ruler says: all you people are like my family---the old man, the old woman, the fasting one, the poor, the rich, all of you are my family, and I take responsibility of this family, called India . This is exactly how leaders speak .
What is Max Weber studying? Social action---every time . And along with that he is moving to bigger and bigger levels, to large-scale social structures as well .

\paragraph{Features of Patrimonialism}
\begin{enumerate}
    \item The administration and military are loyalists---loyal to the leader, not to the nation . When military and administration are loyal to the leader rather than the nation, that is patrimonialism . And that leader may well have had charisma; of course, that is why he could become leader . Reality is not so segregated or layered; we are only using concepts .
    \item The administrators are appointed on the basis of contacts---kin relations or any personal affiliation . Which means not on the basis of merit . There is no merit, only contact .
    \item The administration and military are the personal instrument of the master .
    \item There is no clear-cut division between public and private ends . The so-called salary of administrators or military officers is not given from the treasury . Today the salary of an IAS, IPS or army officer is given from a fund of India; but in those times there was no such contingency fund or consolidated fund . The treasury existed, but it was used only for the expenses of the leader---no other person could take a salary from it .
\end{enumerate}

\begin{DataFact}
What then was the source of salary? Bribe, corruption and loot . It was not a fixed income .
The official had full authority to take a bribe because he was an official; he could coerce the people to give bribes precisely because he was an official .
Remember, this is a model---a model having certain features---not reality .
\end{DataFact}

\paragraph{Ideal Type as Mental Construct}
Remember this very important thing: this is an ideal type . An ideal type is a mental construct . Mental construct means that in reality you would not find the actual type . These are only highly accentuated or exaggerated features of patrimonialism .
What Marx presented to us in the name of capitalism was not capitalism; it was only exaggerated features of capitalism that made us believe we were studying capitalism .

\paragraph{Patrimonial Bureaucracy: Examples}
The Mughal Empire of India was a patrimonial bureaucracy . Remember this word .
\begin{itemize}
    \item Bureaucracy is not exclusive to modern society . Even in traditional societies we had bureaucrats---the treasurers you have read about, the \textit{dhamma mahamatras} who spread the culture and religion in Ashoka's time, the treasurers of the Gupta empire .
    \item The Mauryan empire is largely believed to be an empire having a great degree of bureaucracy: proper administration for the navy, the treasury, the administration of horses and elephants, with committees; and the \textit{Arthashastra} clearly outlined the working of that bureaucratic machinery---it was all about the economic and political administration of the Mauryas .
    \item Even ancient Egypt, the Chinese empire and the Babylonian regime are patrimonial bureaucracies .
\end{itemize}

\begin{ComparisonBox}
\textbf{Empire versus nation:} during the Mughal regime there was no conception of India, no conception of nation as such---it was an empire, not a nation .
In an empire you are indebted to the leader of the empire . In a nation it is not the rule of the leader, it is the rule of law; the leader is not over and above the law .
In the Mughal empire the leaders were over and above all the laws; the administration and military were personal instruments, like slaves to those leaders .
\end{ComparisonBox}

\paragraph{Sultanism}
The extreme form of patrimonialism is sultanism . We do not have to go into these details, but you should have these keywords---because when you read editorials talking about the falling state of democracy, or that democracy in India is not the same that we had, you should have a sense of what they are actually talking about: what they think the ideal type of democracy is, against which they have measured this lower degree and concluded that it is a failed or flawed democracy .

\begin{GovtPolicy}
The scholar Christophe Jaffrelot calls the current Prime Minister's rule that of sultanism---that it is a sultan ruling India, not a democratic leader . This is how the controversy develops and why people start calling him fascist .
Why sultan? Because when media, Supreme Court, High Court, CBI, CVC, NITI Aayog and the press---when every instrument of administration appears to have become the mouthpiece of the leader, it gives the impression that it is a sultan ruling us and not a democratic leader . Sultanism means that the administration has become the personal instrument of the master .
\end{GovtPolicy}

\begin{QuickRevision}
\begin{itemize}
    \item Four ideal types of traditional authority: gerontocracy, primary patriarchalism, patrimonialism (prebendalism), sultanism .
    \item Gerontocracy = rule of the elderly (Sparta, village panchayat, Plato's bar of 50); the legitimising factor is age .
    \item Primary patriarchalism = rule of the father within the household; patrimonialism = the same extended to the state .
    \item Four features of patrimonialism: loyalty to leader not nation; appointment by contacts not merit; administration as personal instrument of the master; no separation of public and private ends, with salary coming from bribe, corruption and loot .
    \item Ideal type is a mental construct of accentuated features; Mughal, ancient Egyptian, Chinese and Babylonian regimes are cited as patrimonial bureaucracies, and Jaffrelot's ``sultanism'' is the extreme case applied to contemporary India .
\end{itemize}
\end{QuickRevision}

\subsection{Charismatic Authority}
Charisma is something that has wings . If you see charisma in a leader, there are very high chances that you may see the same charisma in his son or daughter---the charisma has flown to the other leader .
For example, Rajiv Gandhi and Indira Gandhi---you see a similar charisma in them . You may not see the same in Rahul Gandhi today, but that does not make him someone who is not from the Nehru family; he still is from the Nehru family .

Charismatic authority simply is authority that gains legitimacy on the basis of the heroic, extraordinary qualities that the leader possesses .

\begin{KeyConcept}
Charismatic authority comes in the times of crisis . Whenever society is going through a period of crisis, you will see charismatic authority developing .
You start seeing charisma in that person only because you know the context---it is the crisis in society that has made you look out for someone who could bring back the same power or the same pride .
Such crises are not frequent; they happen once in hundreds of years . That is why such leaders are born once in hundreds of years, and why in hundreds of years you remember only two or three leaders .
\end{KeyConcept}

\paragraph{Illustrations from the Lecture}
\begin{itemize}
    \item \textbf{Hitler}---this is what Max Weber believed when Germany lost World War II: he was looking for someone who could bring the same pride to Germany once again . And afterwards Hitler came in 1933, and people considered him a charismatic personality . \textit{[as stated in lecture]}
    \item \textbf{Gandhi}---it was the crisis situation India was going through . We were oppressed under the British empire and we wanted freedom, liberation .
    \item \textbf{Lord Buddha}---the lower castes were oppressed and were looking for someone who could free them from the oppressive caste structure . It was mainly the Vaishyas, who belong to the third category in the status group, who wanted to improve their status, because their status had been kept very low by the Brahmins . Lord Buddha came in and told them there is no such thing as caste or Vaishya; he challenged the Brahmins .
\end{itemize}

\paragraph{Features of Charismatic Authority}
\begin{enumerate}
    \item Charismatic authority is a historic and transitory phenomenon . When society is in a transition phase, the leader comes to help society move through that transition into a smoother, better world: Gandhi comes to help India become free; Lord Buddha comes to help the lower castes become free of oppressive caste structures; Hitler comes to liberate Germany from whatever he presents as the barrier to a great Germany .
    \item They develop during the times of crisis .
    \item They legitimize their authority on the basis of ancient traditions . Hitler would legitimize himself by taking recourse to the great Aryan rule---the white supremacy, Aryanhood, and the claim that this Aryanhood is in threat . Gandhi would talk about ancient Hindu texts to show that Hindus have always lived in an egalitarian system---he would talk about the Rig Veda and the Bhagavad Gita, to ensure that we live in a peaceful, harmonious fashion .
    \item The charismatic leader may be a hero or he may even be a tyrant .
\end{enumerate}

\paragraph{Routinization of Charisma}
Now where is the problem? The person came during the times of crisis and his rule was transitory . What happens when he dies? 

\begin{KeyConcept}
Write this under the heading \textbf{Routinization of Charisma} .
On the death of the charismatic leader, the charisma gets routinized---that is, transferred---either to the next generation, or it takes the form of a traditional authority or a legal-rational authority .
\end{KeyConcept}

\begin{itemize}
    \item \textbf{Example 1---Prophet Muhammad:} When Prophet Muhammad died, some people believed his son possessed the charisma; some believed the half-brother, Abu Bakr, possessed it . Some believed one man was the legitimate heir to Prophet Muhammad, others said no . The Shia-Sunni divide came from there . \textit{[as stated in lecture]}
    Shia and Sunni became sects, with their own boards and institutions---Shia board, Sunni board---and different administrative organisations catering to the interests of Shias or Sunnis . Not just that: we have nation states which are Shia and Sunni today . Saudi Arabia as a Sunni state; Lebanon, Jordan, Iraq, Iran and Syria as Shia states, with Syria being Alawite Shiite . \textit{[as stated in lecture]}
    The charisma also assumed the form of a text---the Holy Quran---which is traditional authority: ancient tradition and text . The boards and organisations are legal-rational authorities .
    \item \textbf{Example 2---Lord Buddha:} Lord Buddha died, and Buddhism got developed after that as a religion . When Buddha was alive there was no religion called Buddhism; only after his death did Buddhism come in . It also came as text---the three Pitakas in Buddhism .
    \item \textbf{Example 3---Guru Nanak:} On the death of Guru Nanak, Sikhism developed as a religion . The same charisma that Guru Nanak possessed also got passed on to the other nine gurus . After all the gurus, the last guru was the Guru Granth Sahib---traditional authority, ancient text and tradition . This text possessed the charisma of all the gurus .
    We also have Gurdwaras, Gurdwara Prabandhak Committees, the Akali Party---political parties, that is legal-rational authority---Punjab as a state, and the Khalistan movement . All of this is charisma continuing to travel .
\end{itemize}

\begin{QuickRevision}
\begin{itemize}
    \item Charismatic authority = legitimacy resting on heroic, extraordinary qualities; charisma ``has wings'' and can pass to descendants (Indira and Rajiv Gandhi) .
    \item It is a historic and transitory phenomenon, arises in times of crisis, legitimises itself through ancient traditions, and its bearer may be a hero or a tyrant .
    \item Crises come once in hundreds of years---hence only two or three such leaders are remembered (Gandhi, Buddha, Hitler as the lecture's cases) .
    \item Routinization of charisma: on the leader's death charisma transfers to the next generation or converts into traditional authority (Quran, Guru Granth Sahib, three Pitakas) or legal-rational authority (boards, committees, parties, states) .
\end{itemize}
\end{QuickRevision}

\subsection{Legal-Rational Authority}
Legal-rational authority is authority that is legitimized on the basis of rules and regulations . But the question is: what are rules and regulations? We had rules and regulations even during the Mughal empire .

\begin{CriticalPoint}
People may say: but we did not have a constitution . So what---that is just one form . It is not that unless you have a constitution you do not have rules .
In Britain there is no written constitution; they are running by an unwritten one . So you cannot say they do not have formal rules and regulations .
\end{CriticalPoint}

Therefore: legal-rational authority is legitimized on the basis of formal codes, rules and regulations . Weber did not define what a formal code is because all of these are ideal types only .
Remember one most important thing: according to Max Weber, legal-rational authority represents the modern bureaucracy .

\paragraph{What Modern Bureaucracies Have}
\begin{itemize}
    \item A formal code of conduct and written rules of appointment .
    \item Appointment not on the basis of contacts but on the basis of merit .
    \item As an administrator you are not loyal to the minister---you are loyal to the Constitution . You are not even loyal to the President; even if the President tells you something, you may say no, because your loyalty is to the Constitution . You are not a loyalist to anyone .
    \item Modern bureaucracy is highly rational, not emotional; fully bound by laws and legal regulations . It performs \textit{zweckrational} social action and is driven by formal rationality .
\end{itemize}

Did we have this form of bureaucracy in the Mughal empire or the Mauryan empire? No---and why not will become clear when bureaucracy is taught as the next topic .

\begin{KeyConcept}
In conclusion: Max Weber believed that legal-rational authority is the marker of modernity .
All those traditional societies which were run by traditional bureaucracies could not develop the modern form of economy, the modern form of judiciary, or the modern form of capitalism .
\end{KeyConcept}

\begin{CriticalPoint}
Remember: capitalism existed then and exists now; bureaucracy existed then and exists now . The forms of capitalism and bureaucracy have changed .
Rules may be written or unwritten; officials may not be appointed by competitive examination or on merit---they may be appointed on the basis of the spoils system---but still they are bureaucrats, they are officials .
There may not be a clear-cut separation between private and public ends, but bureaucracy is still there .
\end{CriticalPoint}

\begin{QuickRevision}
\begin{itemize}
    \item Legal-rational authority is legitimised by formal codes, rules and regulations---a constitution is only one form of these (Britain runs on an unwritten one) .
    \item It represents modern bureaucracy: formal code of conduct, written rules of appointment, merit-based recruitment, loyalty only to the Constitution, rule-bound and unemotional conduct .
    \item It is driven by formal rationality and performs \textit{zweckrational} social action .
    \item Conclusion dictated in class: legal-rational authority is the marker of modernity; traditional societies run by traditional bureaucracies could not develop modern economy, judiciary or capitalism .
\end{itemize}
\end{QuickRevision}

\subsection{Bureaucracy: Definition and Scope}
Here comes the difference between political science, public administration and sociology . The way you are going to read bureaucracy is very different from the way a student of public administration or political science reads it . Max Weber is ours---he was a sociologist . He was never a political scientist and never a public administration academic .

\begin{KeyConcept}
Bureaucracy is an organizational model, rationally designed to achieve the task in an efficient manner . Write it exactly like this .
Basically it is a machine---a machine used to bring efficiency in the organization . This is what Max Weber calls machine-like efficiency .
Bureaucracy is a model . It is not reality .
\end{KeyConcept}

\paragraph{Bureaucracy Is Not Limited to Government}
Bureaucracy is not limited to the government . Since it is an organizational model, you see bureaucracy everywhere :
\begin{itemize}
    \item In schools, colleges and universities 
    \item In hospitals 
    \item In political parties, in the theatre, in hotels and clubs 
    \item In shopping malls 
    \item In political administration 
    \item And even in religious organisations 
\end{itemize}

\begin{CriticalPoint}
When public administration teaches Max Weber, it limits him only to political administration---but Max Weber himself never limited bureaucracy to political administration .
\end{CriticalPoint}

\paragraph{Why an Ideal Type Is Needed}
Since we know we have bureaucracy in every sphere of society, the question is how to study such a vast reality . A single man cannot study such a vast reality; therefore a model has to be constructed, and that model is the ideal type .

\begin{KeyConcept}
Next line: since bureaucracy is a vast reality, a model---that is, an ideal type---is constructed to highlight the most accentuated features or typical characteristics of bureaucracy .
Accentuated features means typical characteristics . What we are going to study in the name of bureaucracy is a model having certain features; it is not the absolute reality .
\end{KeyConcept}

\begin{QuickRevision}
\begin{itemize}
    \item Bureaucracy is an organizational model rationally designed to achieve tasks efficiently---a machine delivering machine-like efficiency .
    \item It is a model, not reality, and it is not confined to government: schools, universities, hospitals, parties, hotels, malls, theatres and religious organisations all display it .
    \item Weber was a sociologist, so the sociological reading of bureaucracy is wider than the public-administration reading .
    \item Because the reality is vast, an ideal type is constructed to highlight the most accentuated or typical characteristics .
\end{itemize}
\end{QuickRevision}

\subsection{Features of the Ideal Type of Bureaucracy}
\begin{enumerate}
    \item \textbf{Division of Labour:} When you become an IAS or IPS officer you will be allocated a district . In that district you will have many other officers as well---some governing law, some revenue collection, some postal services, some trade or military defence . That means bureaucracy is a model running on division of labour, and division of labour ensures efficiency .
    Write down: division of labour is implemented to ensure efficient performance of the tasks .
    The same is true in any organisation---a technical team, a content team, a video team, a graphics team . Division of labour is there, and that is what bureaucracy is all about .

    \item \textbf{Merit-Based Recruitment:} When you become an IAS, that is based on merit, not on contacts . If we have a technical team, that is purely on merit, because we want good performance . We cannot afford people who come solely by contacts; we need efficient workers . Merit is something that has to be judged . How do we judge someone's merit? We look for experience . A person may even be asked to write a test on a specific thing .

    \item \textbf{Written Rules:} Rules are codified . Nothing will be done orally . For example, if an employee takes a taxi home and comes the next day for reimbursement, he would be asked to produce and submit everything in return . Nothing is done orally . Even in political administration specifically, you will do everything in written form---there will be letters signed by district collectors that are passed down . Even in universities everything is based on written documents; everything is written and codified .

    \item \textbf{Strict Separation of Private and Official Income:} Whatever things are available at your disposal as an administrator, you cannot use those things for your personal gain . When you become an IAS officer you will have a computer on your table---you cannot take that computer to your home . It stays there . You will have a pen, stationery, a printer, a very nice table---these are official . You may even write this as strict separation of private and official ends rather than income .

    \item \textbf{Officials Cannot Improvise:} You cannot change the rules . Rules are already laid out; you just have to act on them . Suppose as an officer you create a special exception for someone very close to you---a close relative comes to you for help and you change the rules specially to help that person out . That cannot be done . You cannot afford to be emotional or creative, or create a special grant or exception for someone .

    \item \textbf{Impersonal Relations:} Impersonality is also one feature of bureaucracy, and it is very similar to the rule that officials cannot improvise . When you talk to your client, that has to be a purely impersonal relation . This is what modern society is all about: there is no place for emotion . It is not \textit{Gemeinschaft}, it is \textit{Gesellschaft} . It is not simple, mechanical solidarity having similar consciousness; it is specialised division of labour and organic solidarity . A \textit{blasé} attitude---not unconcerned exactly, but basically not emotional . Whether it is your mother, your father, your son, your daughter, your friend, or someone you do not know---all are equal to you . That is impersonal relation .

    \item \textbf{Hierarchy---Chain of Command:} Whether you are a teacher in a school, a professor in a university, a political administrator under government machinery, or a staff member at a shopping mall, all of us have to obey the orders given by our superiors . A teacher obeys the principal; a professor obeys the vice chancellor; an assistant collector obeys the collector; the collector obeys the chief secretary . There is a long chain of command that you cannot escape .
    If you do not obey, that has to be justified only on the basis of rules . You cannot simply say ``I will not obey you''; you must cite a justified rule---``I cannot do this, officer, because this is beyond the ambit within which I can work'' . So if an MLA comes to you and asks for a favour, you are free not to obey him . Even with higher officials, you are free not to obey them if that refusal is based on rules . But if the order is within the rules, you will have to obey---you have no choice .
\end{enumerate}

\begin{KeyConcept}
Write down one more line . Max Weber says that the endeavour---the attempt of all modern world societies is to approach or get closer to the ideal type of bureaucracy .
It is the modern bureaucracy, not bureaucracy as such . We had bureaucracy even in the Mughal empire, but that was not merit-based, not purely impersonal, and had no strict separation between private and public ends .
\end{KeyConcept}

\begin{QuickRevision}
\begin{itemize}
    \item Seven features of the ideal type: division of labour, merit-based recruitment, written (codified) rules, strict separation of private and official ends, no improvisation by officials, impersonal relations, and hierarchy/chain of command .
    \item Division of labour is implemented to ensure efficient performance of tasks; merit is judged through experience or testing .
    \item Impersonality places bureaucracy in \textit{Gesellschaft} and organic solidarity, not \textit{Gemeinschaft}---mother, friend and stranger are treated alike .
    \item Disobedience is legitimate only when justified by a rule---which is why an MLA's or even a senior's illegitimate order can be refused .
    \item Weber: the endeavour of all modern societies is to approach the ideal type of modern bureaucracy; nobody has reached it .
\end{itemize}
\end{QuickRevision}

\subsection{Bureaucracy in Reality: The Distance from the Model}
Since we know that what we have read is not bureaucracy but a model of bureaucracy, the reality is different---and this model can be used to understand the differences between reality and the model .

\vspace{6pt}
\noindent
\begin{tabularx}{\textwidth}{>{\hsize=0.45\hsize}Y >{\hsize=0.55\hsize}Y}
\toprule
\rowcolor{AccentDark}
\thead{Feature of the Model} & \thead{What the Lecture Points Out in Reality} \\
\midrule
Strict separation of private and public ends & In every bureaucracy of the world there is corruption; public power is used to serve private ends . Officers take bribes because permission comes only from them; earnings happen under the table . \\
\midrule
Merit-based recruitment and hierarchy & The chain of command is not sacred; it has been violated many times . People become senior or supersede others on the basis of contacts or personal affinal relations developed with the senior-most person in the chain . \\
\midrule
Impersonality and equal treatment & In IAS, IPS and many bureaucracies you would not see lower-caste people promoted to the upper rungs; largely upper-caste men, or at any rate men, are promoted to the upper level . \\
\midrule
Universal access to the top & Lower castes and women gain entry but do not reach the top level . If one or two do, it becomes an exception---and an exception is not a general rule . \\
\bottomrule
\end{tabularx}
\vspace{6pt}

\begin{ExampleBox}
You will see corruption even in the Supreme Court and in universities . Only recently, the head of a department of political science was chosen when someone else---the senior-most according to the seniority principle---should have been there .
Mughal bureaucracy was still bureaucracy, although not on this model: it was patrimonial bureaucracy . Things at their disposal---horses, elephants, cavalry---could be used for personal gain . Today, if you get a horse as an IAS officer from the Government of India, you cannot go for a ride on it .
\end{ExampleBox}

\begin{QuickRevision}
\begin{itemize}
    \item The ideal type is a measuring rod: the gap between model and reality is itself the finding .
    \item Corruption everywhere shows the absence of strict separation of private and public ends .
    \item Hierarchy exists but does not run in a streamlined fashion---contacts and personal relations break the seniority principle .
    \item Caste and gender ceilings mean the promise of impersonal, merit-based promotion is not realised; the exceptions do not make a rule .
\end{itemize}
\end{QuickRevision}

\subsection{The Iron Cage and the Irrationality of Rationality}

\begin{KeyConcept}
Write down: Max Weber believed that the bureaucratization of society is inevitable . The more we get developed or advanced, the more we get bureaucratized .
And therefore it is an iron cage which cannot be escaped by humanity---this is an iron cage of rationality .
Underline, highlight, italicise: it is developed as a result of a large number of people performing instrumental social action and exhibiting formal rationality .
\end{KeyConcept}

Therefore the modern world is a highly disenchanted world---a world having no space for values, customs, traditions, emotions and affection, but a space reserved only and only for efficiency, stability, calculability, rationality and predictability . That is formal rationality .

\begin{CriticalPoint}
And therefore it is not rationality that is going to drive the world; rather it is the irrationality of rationality that is going to occupy the world .
\end{CriticalPoint}

\paragraph{The Trajectory of Rationalization}
This is what the trajectory of societies has been historically . From the dark ages till the modern ages, it is nothing but rationalization that has been going on .
\begin{itemize}
    \item First we questioned religion . Authority was no longer based on religion; now it was based on constitutionalism . It was no longer tradition, customs, religion---rather it was a man coming to power purely on the basis of the Constitution .
    \item We questioned all the traditions . Patriarchy was questioned---Mary Wollstonecraft questioned patriarchy in the French Revolution . So the world would no longer be run only by men; women can also rule .
    \item The world would no longer be run by a religion, a gender or a caste . In India we have been ruled by caste for so long---by Brahmins . Today anybody can become the Prime Minister, Chief Minister or a bureaucrat irrespective of his caste .
    \item That means traditions, customs, values and emotions have been sidelined . What is rewarded is merit and efficiency---rationality, calibre, the ability to perform the task in the least time allotted . It does not matter whether you are Brahmin or Dalit, man or woman .
    \item Rationality, efficiency, calculability and predictability are the principles on which the world now runs . Rationalization has been going on, and this implies disenchantment---you are no longer enchanted with patriarchal rule, no longer enchanted with Brahmins who were believed to possess magical spirits .
\end{itemize}

\paragraph{What the Rational Individual Has Become}
This is rationality: a fully planned day, routine, scheduled, completely disconnected from emotions---not talking to parents, not talking to friends, not going out with friends, no fun time, no entertainment . All the values gone; completely rational, goal-oriented, instrumental man or woman is what we are turning into . And this is what Max Weber cautioned us against .

\paragraph{Examples of the Irrationality of Rationality}
\begin{ExampleBox}
\begin{itemize}
    \item \textbf{MCQ-based selection:} today people become professors, toppers or bureaucrats on the basis of MCQs, because efficiency is what we strive for---maximum result in minimum time . But are you creating creative individuals, or rote learners? People become rote learners, because there is no written examination; you cannot check them on analysis, only on which option is right, which can easily be rote learned . Go through previous year questions, rote learn in two or three days, tick the same thing, and once you are out you forget everything . We are creating rote learners, not learned individuals .
    \item \textbf{TRP in newsrooms:} TRP is rewarded, and what are we creating in its name? Fake news and yellow journalism . When TRP is rewarded, fake news spreads, because fake news gets maximum viewership .
    \item \textbf{The bureaucrat who cannot serve:} you become a bureaucrat to serve the public, but you cannot serve the public because you are bound by the Constitution and by rules . The rules have come in between you and the public . You cannot improvise, cannot go beyond the rules, cannot help anyone within a minute . The person before you may be dying, but the rules ask you to first get a report written . The rules are more important .
    \item \textbf{McDonald's:} you eat there because it is cheaper and saves time . The consequence is obesity, non-nutritious food, environmental degradation and generation of plastic; the world gets more and more toxic .
    \item \textbf{Cars and transport:} produced in the name of efficiency, leading to environmental degradation---and every car is produced by manual labour in the mines, which is the most dangerous work for anyone's health, where health standards are not followed . You only see the finished product, not what is happening backstage . This is commodity fetishism---you do not realise the labour that has gone into it; you only fetishize commodities .
    \item \textbf{GDP and the environment:} we calculate a nation's development in terms of GDP, but this GDP comes at the cost of gross national happiness and of the environment---which is why another programme, the SDGs, became necessary . Potable water is decreasing, groundwater has receded, global warming has melted ice sheets, pollution is carried by jet streams into the darker regions, and even Antarctica is going to be over very soon .
\end{itemize}
\end{ExampleBox}

\begin{PYQLink}
Alvin Toffler, author of \textit{Future Shock}, is described in the lecture as a very great sociologist whose predictions of the 1960s are actually happening today .
Sociology is a very young discipline---barely a hundred years old---but the claim made in class is that after two hundred years sociology will replace every social science, because sociology is not philosophy .
\end{PYQLink}

\begin{QuickRevision}
\begin{itemize}
    \item Bureaucratization is inevitable: the more advanced the society, the more bureaucratized---an inescapable iron cage of rationality .
    \item The cage is built by large numbers of people performing instrumental social action and exhibiting formal rationality .
    \item The modern world is therefore disenchanted: no space for values, customs, tradition, emotion or affection---only efficiency, stability, calculability, rationality and predictability .
    \item The result is not rationality but the irrationality of rationality: MCQs producing rote learners, TRP producing fake news, rule-bound officers unable to help a dying person, McDonald's producing obesity and plastic, GDP produced at the cost of happiness and environment .
    \item Commodity fetishism names the failure to see the labour behind the finished product .
\end{itemize}
\end{QuickRevision}

\subsection{Weber's Ambivalence and the Bridge to Protestant Ethics}

\begin{CriticalPoint}
Max Weber---and specifically Max Weber---did not give any opinion on bureaucracy . He was ambivalent; ambivalent means he was confused .
He did not really know what to call this bureaucracy: is it efficient and should we be happy about it, or is it leading to irrationality and we should not be happy about it? What should be done? Can any individual do anything about it? 
You would say you can change the system . But can you? You become systematized when you enter the system .
\end{CriticalPoint}

That is why very few IAS, IPS or other officers become famous in their fields only because they have gone beyond the rules . They did not limit themselves to the bureaucratic regime that asked them to obey the seniority chain of command; they actually went beyond it, and that is why they became the highlight of the newsrooms . When you do not obey, you become the highlight---and that is an exception .

\paragraph{Where Does All of This Come From?}
Max Weber is telling us through his book: whatever changes you have read so far in history---formal rationality, \textit{zweckrational} social action, legal-rational authority---where are these things coming from?  What is the source of a man, an emotional man, thinking rationally?  What is the source of a man not being concerned about any person but only about his goal, his plan, his schedule? 

\begin{KeyConcept}
The source of all these things is Protestant theology, which we also call the Reformation .
There were two historic events in world history: one was the Renaissance, and the second was the Reformation---the Protestant Reformation . Everything got changed with this only .
And remember, this book is one of the biggest critiques of Marx .
\end{KeyConcept}

The next question is how Marx, Weber and Protestantism---four seemingly different realities---are connected to each other, and how they connect to the way we act rationally right now .

\begin{QuickRevision}
\begin{itemize}
    \item Weber gave no verdict on bureaucracy---he was ambivalent about whether it is efficiency to be welcomed or irrationality to be feared .
    \item Individual escape is doubtful: you become systematized on entering the system, and the officers who become famous are the exceptions who went beyond the rules .
    \item The source of formal rationality, \textit{zweckrational} action and legal-rational authority is Protestant theology, i.e.\ the Reformation .
    \item The Renaissance and the Reformation are the two historic events; the \textit{Protestant Ethic} is one of the biggest critiques of Marx .
\end{itemize}
\end{QuickRevision}

\subsection{Doubts Answered in Class}
\paragraph{Is the irrationality of rationality simply the loss of the human element because of rules?}
Yes . The irrationality of rationality implies losing the human element . Disenchantment is nothing but the irrationality of rationality; it leads to it .

\paragraph{Did Weber, in the end, become a positivist?}
He did not become a positivist . Remember this . He never gave any law . Why can you not give a law as Weber? Because Marx and Durkheim are interested in looking at the world as an objective reality---that there are laws in this world which are running the world .

\paragraph{How Marx Was Presented in This Contrast}
Marx explains through the laws of dialectics . He learned from Hegel that the world moves, but for Marx it is not the idea, it is the material condition . In assuming laws of dialectics he has assumed the world to be objective and monocausal---a cause-effect relation .
Contradictions developed in ancient society, then in feudalism among lords and masters, and the bourgeoisie developed within it . As technology developed, classes developed . Marx sees a single cause-effect relation and no other reason behind the formation of class, the formation of the mode of production, or the evolution of society .

For Marx everything is governed by the economy; the economy is the sole engine . Religion, law and the state are all helping the economy and justifying the condition of the poor and oppressed---telling them they will get heaven in the next life .
Today, according to this reading, the state favours corporations and lords rather than human rights, because those people have power . Society is the organ, the economy is the engine .

\paragraph{How Weber Reads Action Differently}
Weber asks whether an action is merely instrumental and economic, or affectual and emotional; each action must be understood as a pattern that was historically evolving .
There was a time when people did not talk in terms of money . Whatever we produced, we divided with one another; we sat under the sky, singing and dancing . Our action was not instrumental---it was affectual social action and traditional social action . Even festivals were not ``festivals'' as such---they were the mode of living .

\begin{QuickRevision}
\begin{itemize}
    \item Irrationality of rationality = loss of the human element under rules; disenchantment and it are the same movement .
    \item Weber never became a positivist and never gave a law, because for him there is no single objective law of the social world .
    \item Marx by contrast is monocausal and dialectical: material conditions, contradictions, technology and class, with religion, law and state serving the economy .
    \item Weber insists on reading each action historically and notes that pre-modern action was affectual and traditional, not instrumental .
\end{itemize}
\end{QuickRevision}

\sectiondivider

\section{The Protestant Ethic and the Spirit of Capitalism}

\subsection{Irrationality of Rationality: The Student's Own Case}
What is the main concern that all of us have today? We are being driven by rational forces---forces that are not in our hands . Take your own example .
\begin{itemize}
    \item Why are you studying sociology? You are not studying sociology for sociology . That is not affectual social action; it is instrumental social action . So this is rationality .
    \item But this rationality is irrational, because you are studying sociology for UPSC . In that sense you do not know sociology .
    \item You study to make keywords and to make notes; you study to answer the questions . All through the lecture you are concerned with the keywords you are going to write in that answer .
    \item The moment the lecture is over the question comes: ``Sir, how to answer this question? How to get 6 marks out of 10? How will I get 12 out of 20, how will I get 300?'' Even after a lecture on how rationalization has trapped us in the iron cage of rationality, the immediate question is about scoring .
\end{itemize}

\begin{CriticalPoint}
That means you are simply manifesting instrumental social action . Is this in your hand? Can you change it? This is not in your hand .
You have already become part of the system . You cannot change it .
This is the irrationality of rationality: you lose the charm of everything you listen to or see, because you become highly instrumental .
\end{CriticalPoint}

\begin{QuickRevision}
\begin{itemize}
    \item Studying sociology for UPSC rather than for sociology is instrumental (\textit{zweckrational}) social action .
    \item The action is rational in means but irrational in outcome---the charm and substance of what is learned is lost .
    \item The individual cannot escape it because he has already become part of the system .
\end{itemize}
\end{QuickRevision}

\subsection{The One-Crore Thought Experiment: Formal versus Practical Rationality}
Before understanding Max Weber's \textit{Protestant Ethic}, a few things must be very clear . All the things you have read so far are rivers; this is the ocean into which all the rivers are going to merge . Now you will understand social action better, rationality better, and why we read authority at all . Everything is in this book .

Suppose I give you one crore rupees and ask you to do just one thing with it, and to give me an honest reply---brutally honest, because this is not a personality test of UPSC where you have to fake .

\begin{CriticalPoint}
An aside made in class: UPSC does conduct a personality test, but the very act of testing personality would not let the real personality come out . This is the inherent contradiction within the system .
\end{CriticalPoint}

Basically all of you said you would invest . Which kind of rationality is being exhibited there? 

\vspace{6pt}
\noindent
\begin{tabularx}{\textwidth}{>{\hsize=0.3\hsize}Y >{\hsize=0.25\hsize}Y >{\hsize=0.45\hsize}Y}
\toprule
\rowcolor{AccentDark}
\thead{Choice with One Crore} & \thead{Kind of Rationality} & \thead{Explanation Given in Class} \\
\midrule
Investing the money & Formal rationality & You weigh everything, calculate everything, and try to maximise utility out of what is in your hand . \\
\midrule
Buying a Mercedes C-Class---fifty lakh in the car and the rest in its petrol & Practical rationality & Like picking up a two-thousand rupee note and spending it . No calculation of maximum utility, simply spending it . \\
\bottomrule
\end{tabularx}
\vspace{6pt}

Your investment in property is a manifestation of instrumental social action, because you are investing for the family, for the future---you have all those people in your head, and one day you have to look after them as well . You would invest in property or in the stock market, thinking that the more you get, the more beneficial it would be for your future or your family . That is perfectly \textit{zweckrational} social action .

\begin{KeyConcept}
But what is the origin of these things? If three hundred years back I had given you one crore rupees, would you have invested it in the stock market or in property? No .
This is the main theme of the \textit{Protestant Ethic} . In one line: the fact that you would invest one crore rupees today means, Max Weber says, that you should be thankful to the Protestants . This attitude of investing money comes only from Protestantism .
\end{KeyConcept}

\begin{QuickRevision}
\begin{itemize}
    \item Formal rationality = calculating and maximising utility (investing the crore); practical rationality = spending it as it comes (the Mercedes) .
    \item Investment for family and future is \textit{zweckrational} (instrumental) social action .
    \item The disposition to invest is historically recent---Weber traces it to Protestantism .
    \item The \textit{Protestant Ethic} is the ``ocean'': social action, rationality, ideal type and authority all flow into it .
\end{itemize}
\end{QuickRevision}

\subsection{A Detour: How to Read History for the Examination}
To those who do not have even a little idea about the word Protestant: if you know Buddhism, Buddhism today is a religion, but it never developed as a religion---it actually was a protest . Buddhism developed as a protest against Hinduism . Students who have not read ancient history should, if nothing else, definitely read Buddhism .

\begin{PYQLink}
When you read ancient history, do not perform affectual social action; do not get emotional about ancient history . You have to be highly rational .
You cannot afford to read the whole of ancient history from R.S.~Sharma . The right method is: go through previous year questions, check the weightage of questions area-wise, and in a meticulous, planned fashion read the points from the book that gives maximum result in minimum time .
Throughout this you are performing \textit{zweckrational} selection only .
\end{PYQLink}

\paragraph{Ancient History---What to Read}
\begin{enumerate}
    \item Harappa---Harappan civilisation, that is the Indus Valley Civilisation; both are one and the same .
    \item Buddhism 
    \item Jainism 
    \item The Mauryan Empire 
    \item The Guptas 
    \item Scientific achievements during ancient India---trigonometry and the rest; this is the last chapter of R.S.~Sharma .
    \item Philosophies of ancient India---Nyaya, Mimamsa, Vaisheshika; the Darshanas, the six Shad Darshanas . This is the second-last chapter .
\end{enumerate}
You do not have to read the Greeks, Shakas, Parthians and Kushans; you may even skip the Satavahana dynasty and such topics . If you do only this much, it will take you only three days---maximum .
Whoever is preparing should start ancient history right now, because ancient history is generally considered boring and whatever you think is boring should be read first, since otherwise you will not read it later at all .

\paragraph{Medieval and Modern History}
\begin{itemize}
    \item \textbf{Medieval:} no Sultanate . Only the Mughals, plus Chola and Pallava (according to some books these are part of ancient history, according to others early medieval), and Bhakti and Sufi . Only these four for medieval .
    \item \textbf{Modern:} you do not have to read early modern history---the period from the 1750s to 1857 before the revolt . Just read the major events in modern history along with the major acts . That is it .
\end{itemize}
Always keep previous year questions beside you while reading these things . You cannot remember everything, and therefore you should not even try to become a scholar in modern history . Last year in modern history many questions came which are not present anywhere---not in NCERT, not in Spectrum, not in any general book read by competitive aspirants . What is in your hand you can do; what is not in your hand becomes part of the competition .

\paragraph{Art and Culture}
Art and culture: NCERT only . From it, Buddhist art is definitely important, along with paintings, Indo-Islamic architecture, folk dances and classical dances .

\begin{QuickRevision}
\begin{itemize}
    \item Read history instrumentally: previous year questions first, weightage next, then the minimum book that gives maximum result .
    \item Ancient: Harappa/IVC, Buddhism, Jainism, Mauryas, Guptas, scientific achievements, philosophies (Shad Darshanas)---skip Greeks, Shakas, Parthians, Kushans, Satavahanas .
    \item Medieval: Mughals, Cholas, Pallavas, Bhakti and Sufi only . Modern: major events plus major acts, skipping the pre-1857 early modern period .
    \item Art and culture: NCERT only---Buddhist art, paintings, Indo-Islamic architecture, folk and classical dances .
\end{itemize}
\end{QuickRevision}

\subsection{Buddhism as Protest: The Analogy for Protestantism}
Lord Buddha challenged Hinduism . In Hinduism we had a varna structure: Brahmins at the top, then Kshatriya, then Vaishya, and Shudra at the bottom . This was graded inequality, and Lord Buddha challenged this inequality .
\begin{itemize}
    \item Lord Buddha himself was a Kshatriya; he belonged to a family of princes . Mahavir Jain was also a Kshatriya . Both of them were Hindus .
    \item Buddha challenged Hinduism because of corruption, because of inequality, because of the monopoly of Brahmins, and for many other reasons .
    \item From there on people started reposing their faith in Buddha . They saw him as someone very charismatic, someone who could challenge the Brahmins, someone who could give a new way of life and tell them how to lead their lives, what should be done and what should not be done, how to spend their precious life on earth .
    \item A new way of life was told to the masses by Buddha, and from here on Buddha started gaining that halo behind his head . People actually made Buddha what Buddha was .
\end{itemize}

\paragraph{Reformation in Hinduism}
\begin{ExampleBox}
Reform in Hinduism began as a response to the challenge posed by Buddha . Hinduism had to reform itself because Brahmins realised that most of their followers were becoming Buddhists, and they wanted to retain them .
One of the reforms was that the use of animals in agricultural activities would stop, and the eating of beef would stop .
Within Hinduism, to eat beef was earlier a natural thing, considered at par with the values of Hinduism . Only after the challenge posed by Buddha did beef eating, animal eating, the sacrifice of animals and the use of animals in agriculture get suspended .
\end{ExampleBox}

\paragraph{Reformation within Buddhism}
Within Buddhism too there were disputes after the death of Buddha . People began arguing about which beliefs should continue to exist and which should not .
Because of two opposing groups, Buddhism split into multiple factions---initially into two, and then a third developed: Mahayana and Hinayana . This was nothing but reformation in Buddhism .
Lord Buddha told the people that the Brahminical monks tell you not to do these things but they themselves are indulging in them: they ask you to maintain celibacy but do not; they ask you not to eat non-vegetarian food but they eat it; they ask you not to amass wealth but they have amassed enormous wealth---all the corrupt things through the back doors . People got enlightened, which is why that period is called the Indian enlightenment period---\textit{nirvana} .

\begin{KeyConcept}
A similar fate befell Christianity in the early sixteenth century . It was the rule of the Catholic Church, just as we had the rule of Hinduism until Buddha came .
We had the rule of Catholicism until the ``Buddha of Germany'' came---that is, Martin Luther . Martin Luther acted just like Buddha .
Just as Buddhism is an offshoot of Hinduism, Protestantism is an offshoot of Catholicism . To understand Protestantism you have to understand Catholicism .
\end{KeyConcept}

\begin{DataFact}
Some facts given in class about Protestants:
\begin{itemize}
    \item The founder-owners of McDonald's---Richard McDonald and Ray Kroc---were Protestants . McDonald's was the first fast food chain in the world; the fast food culture started with it, and KFC, Domino's, Pizza Hut, Burger King and the rest came by copying that concept . This is why George Ritzer studied McDonald's and why McDonaldization became such an important topic .
    \item Warren Buffett---although he says he is agnostic and does not believe in God---is a Presbyterian, and Presbyterianism is a sect within Protestantism, so he is largely a Protestant .
    \item Bill Gates was born in a Protestant church; Benjamin Franklin was a Protestant .
    \item There is something special about Protestants that no other people possess .
\end{itemize}
\end{DataFact}

\begin{QuickRevision}
\begin{itemize}
    \item Buddhism began as a protest against Hindu graded inequality (Brahmin-Kshatriya-Vaishya-Shudra); Buddha and Mahavir were both Kshatriyas and both Hindus .
    \item Buddha's challenge forced reformation within Hinduism (end of beef eating and animal use in agriculture) and Buddhism itself split into Mahayana and Hinayana .
    \item The structural parallel: Catholicism: Protestantism :: Hinduism: Buddhism; Martin Luther is the ``Buddha of Germany'' .
    \item Modern illustrations of Protestant distinctiveness: McDonald's founders, Warren Buffett (Presbyterian), Bill Gates, Benjamin Franklin .
\end{itemize}
\end{QuickRevision}

\subsection{The Doctrines of Catholicism}
In the early sixteenth century, in Germany and in Europe, the Catholic Church was ruling, and the Catholic Church had these doctrines :
\begin{enumerate}
    \item You can attain salvation by purchasing indulgences and by good work .
    \item You have to pass through holy rituals like confirmation and baptism .
    \item You must work until your immediate needs are met; once your immediate needs are fulfilled you do not have to work any more .
    \item Profit is sin .
\end{enumerate}

\paragraph{What Was an Indulgence?}
\begin{ExampleBox}
At that point of time the Church used to print a letter of indulgence . Suppose you have committed a sin and you want to atone---to rid yourself of the sin and secure your place in heaven, since your entry into heaven is now debarred .
The Church used to say: whoever has committed a sin, do not worry, we will give you a letter signed by God---but you will have to give us some money for it, and then all your sins will be washed away .
People used to come and stand in very long queues waiting for the priest .
\end{ExampleBox}

Good work meant: if you follow the advice of the Church and live according to the Catholic doctrine, you are doing good work . If you are listening to the Church and following the command of the priest, you are doing good work; otherwise you are not .
Confirmation was a holy ritual of which everybody had to be a participant member . A priest with a long cap on his head would hold a stick, put his hand on your forehead, you would bow, and you would go away with the letter of indulgence saying that your seat in heaven is secured .
Baptism involves the holy water in which you have to immerse yourself .

On the third doctrine: if I have worked enough to meet my immediate needs, I should stop working . Because if I am working even after meeting my immediate needs, that means I am working out of greed . To escape greed, work only until your needs are fulfilled; whatever you do after that is sin .
On the fourth doctrine: if you are working to earn profit, you are going to commit a sin . And if by chance you have committed a sin, come back to us---we will give you a letter of indulgence .

\begin{QuickRevision}
\begin{itemize}
    \item Four Catholic doctrines: salvation through indulgences and good work; compulsory holy rituals (confirmation, baptism); work only until immediate needs are met; profit is sin .
    \item The letter of indulgence was a purchased, signed guarantee of heaven; good work meant obedience to the Church .
    \item Working beyond immediate needs was condemned as greed---which is precisely what blocked the accumulation of capital .
\end{itemize}
\end{QuickRevision}

\subsection{Martin Luther and the Protestant Reformation}
One fine day Martin Luther came, and he challenged the Catholic Church . The way he challenged was also very unique: he wrote 95 theses .

\begin{GovtPolicy}
The date given in class for Luther's act at Wittenberg, Germany was stated as October 31, 1570, while the schism that followed was dated to the event of 1517 . \textit{[as stated in lecture]}
Martin Luther wrote, in 95 sheets, what the problems with the Catholic doctrine were: how the Catholic people were looting the common masses, how the Catholic Church people were corrupt, and in what fashion they were amassing wealth and looting people in the name of God .
These 95 sheets were nailed on the door of the main church .
You do not have to remember these details for the answer; this is background you should know .
\end{GovtPolicy}

Luther refused to recant: he said he would not take back his words, because whatever he had done was out of pure conscience . The word \textit{conscience} was heard by those people for the first time---they knew only God, heaven and salvation . Man, in that world, was not a conscious being; man was simply sent on earth by God .

\paragraph{What Luther Proposed}
\begin{enumerate}
    \item You can attain salvation not by purchasing letters of indulgence, but only if you have faith in God . That means there is no need for letters of indulgence, baptism or confirmation .
    \item These priests are intermediaries between us and God---and why should we believe them? If God has made all of us, then all of us are the children of God . Why should these people possess some special quality, as if they are the guardians of God? How can they become the custodians? We have to have a direct relation with God .
    \item Go back to the teachings of the Bible, and do not make your own interpretations .
\end{enumerate}

\begin{ComparisonBox}
Something of this sort happened in the Bhakti movement in India . Bhakti and Sufi movements were all about our direct connection with God, not relying on priests or an intermediary class between us and God .
\end{ComparisonBox}

\paragraph{The Bible in German}
The basic reason people used to depend on the priest was that the Bible was written in Latin, and Latin was the language of the elites---just as Sanskrit was the language of elites here . Since it was an elite language, common people could not read it themselves, and this established a dependency relation between the people and the priest .

Martin Luther is the first person to translate the Bible into German . When he translated it, common people could read the German Bible---the New Testament---all by themselves .
In Germany the culture was that if you purchased a book, you had to read it out openly to the people around you---not carry it home in an envelope and read it peacefully as we do today .
So the moment some people purchased the German version of the Bible, they started reading it out, and people listened---and they got to know that nothing of the sort was written in the Bible in the name of which the Catholic priestly class had been looting them through letters of indulgence, baptism and confirmation . They realised they had been fooled, and they were highly thankful to Martin Luther for enlightening them .

\paragraph{From Protest to Protestantism}
People started following Martin Luther like anything; to them he was a charismatic personality, and they obeyed whatever he told them .
Luther's life was now at risk, because he had committed an act of heresy and was going to be excommunicated . But the Church knew that this man now had so many followers that it was very reluctant to kill him right away .

Suddenly a new wave erupted in Germany: for the first time a man was being followed by many people---a man who was not the interpreter of the text of God, an ordinary man .
The followers of Martin Luther are today called Evangelicals; they are also called Lutherans, and what formed was Lutheranism---just as Buddhism developed after Buddha .

This was a reform within Christianity . The Christian church had to reform itself, and the result of all this was the Reformation of Christianity .
The Reformation was the result of protest, and the people who protested and followed Martin Luther---because of whom Christianity had to reform itself---came to be called Protestants; hence Protestantism .

\begin{KeyConcept}
This event led to the breakdown of the church . At one point of time there was only one church, one empire, one set of followers; after this event the church was divided into two---the Catholic Church and the Protestant Church .
That is, the event led to a schism---a breakdown or split---just as there was a schism within Buddhism (Hinayana and Mahayana) and within Islam (Shia and Sunni) .
\end{KeyConcept}

\begin{QuickRevision}
\begin{itemize}
    \item Luther nailed 95 theses on the church door at Wittenberg, listing the corruption of the Catholic Church, and refused to recant, invoking conscience .
    \item His three proposals: salvation by faith alone; rejection of priests as intermediaries and direct relation with God; return to the teachings of the Bible .
    \item Translating the Bible into German broke the dependency created by Latin, and Germany's custom of public reading spread the discovery quickly .
    \item His followers became Lutherans/Evangelicals; the protest produced the Reformation, Protestantism, and the schism of the church into Catholic and Protestant .
\end{itemize}
\end{QuickRevision}

\subsection{The Doctrines of Protestantism}
People became Protestants but they still wanted to attain salvation . The belief in heaven is not gone; only the method of going to heaven has changed . It is no longer holy rituals, purchasing indulgences or good work . The people who had to form the new doctrines were, besides Martin Luther, above all John Calvin, who was highly inspired by the writings of Martin Luther, Plato and Aristotle, and who formed his own doctrine which later became part of the Protestant ethic .

\paragraph{1. Predestination}
\begin{KeyConcept}
John Calvin said something very landmark that changed the discourse of history: who is going to go to heaven has already been decided . So what you do on earth does not matter .
The \textbf{elect} are those already chosen to go to heaven; those who are doomed to hell are also already chosen, and they are the \textbf{damned} .
This is the doctrine of predestination .
\end{KeyConcept}

Consider the parallel drawn in class: if I told you that who is going to become an IAS has already been chosen---that out of this batch twenty are destined---but I do not tell you the names, what will happen? It will create anxiety .
That anxiety among the Protestants is called \textbf{salvation anxiety} .

\paragraph{2. Calling}
Calling means that when you do any work, your act of doing work is a call . When you are working you are actually obeying God, and it is in the work that your spiritualism is reflected .
Your act of doing work is a reflection of obedience to the Divine---what in Hinduism we call Dharma: to fulfil your duty and your obligations .
When you are working you are fulfilling your obligations and duty not just towards your family but also towards the Holy, because this is why you have been sent to Earth .
Very important line: work was to them a spiritual experience .

\paragraph{3. Asceticism---Ascetic Materialism}
Ascetics are those people who make money but do not spend it on themselves .
But being ascetic meant a lot more than this to Calvin . Calvin said whatever earning you have made out of your work should not be spent on consumption; it should be spent on something that is more productive and something that will help other people as well .

So this was not just asceticism---it was \textbf{ascetic materialism}: do not spend on yourself, but invest whatever you have made in something productive or something that is going to help other people as well .
Ascetic means someone who does not spend for consumption or to display his wealth---the opposite of conspicuous consumption . You may earn in billions but lead a very simple life .
Examples given in class: the Brahma Kumaris, who are billionaires yet lead a very simple life; and Ramdev, described as completely ascetic .

\paragraph{4. Denial of Worldly Pleasures}
Denial of worldly pleasures simply means that when you work, you must be completely focused on work and not on other pleasures . Renounce pleasure; focus on work .
Why focus on work? Because it is only through work that you are going to go to heaven---work is a spiritual experience .

\begin{KeyConcept}
Hidden in these four doctrines is what you do every day in your life---and why the British came to India, and how world history got changed . Everything is related to Martin Luther and John Calvin .
\end{KeyConcept}

\paragraph{Profit Reconsidered}
\vspace{6pt}
\noindent
\begin{tabularx}{\textwidth}{>{\hsize=0.3\hsize}Y >{\hsize=0.35\hsize}Y >{\hsize=0.35\hsize}Y}
\toprule
\rowcolor{AccentDark}
\thead{Question} & \thead{Catholicism} & \thead{Protestantism} \\
\midrule
Is profit sin? & Profit is sin . & Profit is not sin---you get profit only by working, and work is a spiritual experience . \\
\midrule
What then is money? & Wealth beyond immediate needs is greed . & Money is the grace of the holy divine for the work you have done . \\
\midrule
What is the sign of salvation? & Rituals, indulgences and good work . & Material success---maximum wealth generation is a reflection of maximum holy divine grace, and hence a sign of being among the elect . \\
\bottomrule
\end{tabularx}
\vspace{6pt}

The parallel offered: in Hinduism we have Lakshmi Mata, the goddess of wealth---so when we work we are working to have Lakshmi, not for our profit .

So the Protestant doctrine was all about focusing on work and generating profit, with material success as the sign that you are going to go to heaven .
Given salvation anxiety, who will go to heaven? Only those who generate maximum wealth, because maximum wealth generation reflects maximum divine grace . So if I want to know whether I am destined for heaven, I must work and work until I make that much money .

\begin{QuickRevision}
\begin{itemize}
    \item Four Protestant doctrines: predestination (elect and damned), calling (work as obedience to the Divine, a spiritual experience), ascetic materialism (invest productively rather than consume), and denial of worldly pleasures .
    \item Predestination produces salvation anxiety; John Calvin, inspired by Luther, Plato and Aristotle, is its author .
    \item Profit ceases to be sin and becomes the grace of the divine; material success becomes the visible sign of election .
    \item The believer therefore works unceasingly to accumulate---not out of greed but to read his own salvation .
\end{itemize}
\end{QuickRevision}

\subsection{Weber's Question and the Ideal Type of Capitalism}
Max Weber was dealing with one basic question: how did capitalism develop---and not just why, but why did capitalism develop in only two pockets of the world, England and the Netherlands? 

\begin{DataFact}
If you look into the statistics, in Europe most of the people who are prosperous and the countries which have high GDP are those countries which have Protestants .
Till date, the countries with high GDP are those which are Protestant, and the people who are billionaires are Protestants .
There must be something special about Protestants that no other people possess .
\end{DataFact}

To answer this, Weber first had to explain the meaning of capitalism---because capitalism too is just a model . He knew he was dealing with a model that is not absolute reality .
Did we not have capitalism before these countries? Weber says we had capitalism even before England and the Netherlands, even before the Protestant Reformation, even before the industrial revolution . But there is something special about this capitalism which no other form had .

So Weber formed an ideal type of capitalism . According to him, capitalism is characterised by nothing but amassing wealth: the more wealth you accumulate, the more capitalist you are . Wealth accumulation is the accentuated feature, the essence of capitalism .
But wealth has been accumulated by many people in many fashions through the pages of history---so they too must be called capitalists . There is no reason not to .

\begin{QuickRevision}
\begin{itemize}
    \item Weber's research question: why did capitalism develop, and why only in England and the Netherlands? 
    \item Supporting observation: high-GDP European countries and billionaires are disproportionately Protestant .
    \item Capitalism is a model; the ideal type of capitalism has wealth accumulation as its accentuated feature .
    \item Since wealth was accumulated in many historical forms, those forms too count as capitalism---hence the need for four ideal types .
\end{itemize}
\end{QuickRevision}

\subsection{The Four Ideal Types of Capitalism}
Since wealth accumulation took place in different forms through history, Max Weber formed four ideal types of capitalism :
\begin{enumerate}
    \item Booty capitalism 
    \item Traditional capitalism (merchant capitalism) 
    \item Pariah capitalism 
    \item Legal-rational capitalism 
\end{enumerate}

\begin{ExampleBox}
A note on the word \textit{pariah}: it is spelt p-a-r-i-a-h . The origin of the word is from a tribe in India---the Pariah tribe---from where Max Weber took the word .
\end{ExampleBox}

\paragraph{Booty Capitalism}
Booty capitalism is the one where wealth is accumulated by plunder---wealth accumulation by stealing .
Who were these people? All the invaders---those who invaded India or any country---and pirates, such as those of Somalia . All of these are booty capitalists .
The Aryans who came to invade India, and the Dasyus who made the indigenous people slaves and looted them, were also booty capitalists . They looted a tribe, hoarded that tribe's wealth, got satisfied, and after some days invaded another tribe . This is how the expansion of empires was happening .

\paragraph{Traditional (Merchant) Capitalism}
Traditional capitalism happened among those who used to collect something very expensive or luxurious---for example the silk route . This is what we also call merchant capitalism .
Today's capitalism is all about mass production: you mass produce, things become cheaper, demand rises, and this demand impels another level of supply---mass production and mass consumption go hand in hand .
But in merchant capitalism there was no mass production . There was production for a class by a class: merchants catered to kings and nobility, and that was the silk route .
The nobilities produced things which were very expensive and could not be afforded by the common masses . That was also capitalism: wealth was accumulated by a select few who could afford it, or generated by those who could produce those expensive things . It was capitalism limited to the nobility or the elites .

\paragraph{Pariah Capitalism}
Pariah capitalism was the result of money lending . You lend your money to someone, and in that case too you accumulate more wealth .
This is what banks do today---a bank such as ICICI earns by lending . There is a hidden layer here, because reality is not of a pure type; we have layers .
Think about the Jews . Jews were hated by Europeans because money lending was the means through which they accumulated wealth . If you see many Jewish billionaires and millionaires today, the real source of their wealth was lending activities . Lending money to generate wealth was considered a sin in Catholic doctrine, and that is why Jews were hated initially---but this was also a form of making money .

\paragraph{Legal-Rational Capitalism}
The wealth accumulation we had after the industrial revolution was a revolutionary form, very different from any other form in history .
It was not a class producing for a class; not wealth generated by a few by lending money; not wealth generated by looting people .
It was mass generation of wealth, contributing to the rise of the wealth of nations; mass production catering to mass consumption; economies of scale . This is legal-rational capitalism---a very revolutionary form not present earlier in history .

\begin{KeyConcept}
Write down one line: according to Weber, the only feature that distinguishes the other forms of capitalism from present-day legal-rational capitalism is the \textbf{unrelenting pursuit of profit} .
One crore becomes two crore, two crore becomes three crore; produce cars, advertise, reinvest---the unrelenting pursuit of wealth . This was not there in the earlier forms .
\end{KeyConcept}

\vspace{6pt}
\noindent
\begin{tabularx}{\textwidth}{>{\hsize=0.3\hsize}Y >{\hsize=0.35\hsize}Y >{\hsize=0.35\hsize}Y}
\toprule
\rowcolor{AccentDark}
\thead{Form of Capitalism} & \thead{Source of Wealth} & \thead{Illustrations from the Lecture} \\
\midrule
Booty capitalism & Plunder and loot & All invaders; pirates of Somalia; Aryans and the enslavement of indigenous people by the Dasyus . \\
\midrule
Traditional / merchant capitalism & Luxury production for a class by a class & The silk route; merchants catering to kings and the nobility; no mass production . \\
\midrule
Pariah capitalism & Money lending at interest & Jews in Europe; modern banks such as ICICI; lending was a sin in Catholic doctrine . \\
\midrule
Legal-rational capitalism & Mass production for mass consumption, economies of scale, reinvestment & Post-industrial-revolution capitalism; defined by the unrelenting pursuit of profit . \\
\bottomrule
\end{tabularx}
\vspace{6pt}

\begin{QuickRevision}
\begin{itemize}
    \item Four ideal types: booty (plunder), traditional/merchant (luxury goods, silk route), pariah (money lending), legal-rational (mass production and reinvestment) .
    \item ``Pariah'' is borrowed from the name of an Indian tribe .
    \item The single distinguishing feature of modern legal-rational capitalism is the unrelenting pursuit of profit through reinvestment .
    \item All four are ideal types---models, not the reality itself .
\end{itemize}
\end{QuickRevision}

\subsection{Protestants and the Industrial Revolution: Weber's Sources}

\begin{DataFact}
In the industrial revolution, 70 to 80 per cent of the employees, staff, factory owners, managers and workers were Protestant .
The countries where the industrial revolution happened---largely England and the Netherlands---were at that point of time Protestant countries .
\end{DataFact}

\begin{KeyConcept}
Write down: without understanding the social action of Protestants, one may not understand the rise of rational capitalism .
\end{KeyConcept}

The next question in Weber's mind was: how can I understand Protestants---their social actions and the meanings they attach to work? Everything narrated so far as history was also read by Max Weber .

\paragraph{Benjamin Franklin}
Other than Luther and Calvin, Weber had also read one more man in America: Benjamin Franklin . Max Weber paid only one visit to America, in 1904---the same year in which he published his book, the \textit{Protestant Ethic} . \textit{[as stated in lecture]}

Franklin had written many books on how to become wealthy---as today one reads books like \textit{Rich Dad Poor Dad} . He had published three books: \textit{Advice to a Young Tradesman}, \textit{The Way to Wealth}, and his \textit{Autobiography} . Max Weber read these three books .
So Weber read John Calvin, read Martin Luther, read these three books, and also read the Industrial Revolution---and found a connection in all of them .

\begin{KeyConcept}
Next line: Weber read books written by Benjamin Franklin to understand the meanings and motives attached by Protestants to their work .
In Franklin's books Weber read the phrases we take for granted today: \textit{money begets money}, \textit{time is money}, \textit{work is worship}---which later came to be regarded as world-revolutionary thoughts to be published in every corner of the school .
\end{KeyConcept}

\paragraph{Weber's Mother}
Max Weber's mother was also a Protestant, and he dedicated the book to his mother .
His mother was a social philanthropist . She used to earn a lot of money but never used to invest it on herself---she led an ascetic lifestyle, exactly on the lines of Protestantism . She used to invest in some things, and she also used to donate .
The comparison drawn: just as today we have the Bill and Melinda Gates Foundation---donating is the sign of the billionaire .

\begin{QuickRevision}
\begin{itemize}
    \item 70--80\% of those involved in the industrial revolution were Protestants; England and the Netherlands were Protestant countries .
    \item Dictated line: without understanding the social action of Protestants one cannot understand the rise of rational capitalism .
    \item Weber's sources: Luther, Calvin, the Industrial Revolution, and Benjamin Franklin's three books (\textit{Advice to a Young Tradesman}, \textit{The Way to Wealth}, the \textit{Autobiography}); one visit to America in 1904, the year of publication .
    \item Franklin's phrases---money begets money, time is money, work is worship---supplied the meanings Protestants attached to work; Weber's own mother, an ascetic Protestant philanthropist, is the book's dedicatee .
\end{itemize}
\end{QuickRevision}

\subsection{The Ideal Type of Protestant Ethics}
Finally Weber developed an ideal type of Protestant ethics . He knew that what he was going to understand about Protestants would be his understanding only---not what Protestantism is in itself---but it was his attempt to understand the meanings and motives Protestants attach to work .
\begin{enumerate}
    \item \textbf{Frugality}---be frugal, save your money, because money begets money . This is what he had learned from Benjamin Franklin .
    \item \textbf{Industriousness}---be industrious, do hard work, and never complain in the work; because work is a spiritual experience, and if you complain about the work you are complaining against the wish of God . You have to work unrelentingly .
    \item \textbf{Work is worship} 
    \item \textbf{Time is money} 
    \item \textbf{Calling}---it is in the work only that you gain spiritual experience .
    \item \textbf{Predestination} 
    \item \textbf{Ascetic materialism} 
    \item \textbf{This-worldly orientation}---orientation to this world, not to the other world . Focus on your work; do not think about going to heaven . To go to heaven is a result of doing work, so focus on work first . You cannot go to heaven merely by standing on the road with faith in God and taking His name .
\end{enumerate}

\begin{KeyConcept}
Weber clearly says: I am not saying this is what Protestantism is . I am saying this is a model of Protestantism .
\end{KeyConcept}

\begin{QuickRevision}
\begin{itemize}
    \item Eight elements of the ideal type of Protestant ethics: frugality, industriousness, work is worship, time is money, calling, predestination, ascetic materialism, this-worldly orientation .
    \item Frugality and time-discipline come from Franklin; calling and predestination from Calvin .
    \item Complaining about work is complaining against the wish of God---hence uncomplaining, unrelenting labour .
    \item Weber presents this explicitly as a model of Protestantism, not Protestantism itself .
\end{itemize}
\end{QuickRevision}

\subsection{Elective Affinity and Adequate Causality}
Weber now had two models floating before him---the ideal type of Protestant ethics on one hand and the ideal type of rational capitalism on the other . Looking at them together, he saw all the elements in the Protestant ethic which ultimately and unintentionally led to the rise of rational capitalism .

It was only in Protestantism that the spirit of capitalism was born . The spirit lay in the work ethic of the Protestants, which is also called the \textbf{Puritan work ethic} .
What does the capitalist want? Capitalists always reward productivity; they reward hard work, discipline and punctuality; they reward frugality and justness . All these elements are the spirit of capitalism, and they were present in the Protestant .

\begin{KeyConcept}
Those people, unconsciously, through their social action of working harder and harder to secure their seat in heaven, ended up creating the industrial revolution and giving jobs and employment to so many people .
They were working only to secure a place in heaven or to get the grace of God . They never knew what capitalism or the industrial revolution would become .
That is why Max Weber said that rational capitalism was an accident in history . It was never planned . They were never working to set up factories .
They never wanted to exploit anyone . It was never their intention to exploit people the way Marx looks at it, never to loot people, never to appropriate wealth or surplus value, never to make people work longer hours .
Therefore: capitalism was the unintended consequence of the Protestant Reformation . Martin Luther never knew this would happen; John Calvin never knew .
\end{KeyConcept}

\paragraph{Is This a Monocausal Claim?}
Is Weber saying Protestantism is the only cause behind the rise of capitalism? Is he drawing a monocausal relationship between the two ideal types? No . He never said anything like this .
Because reality is very vast, what we can do is develop models; and when you draw a relationship between models you cannot be determinist about the causal relationship .
He says capitalism was an accident in history, an unintended consequence of the acts of people who were working only to secure their place in heaven . So he cannot make a monocausal explanation . But one thing he can definitely see is that there is a relation between them .

\begin{KeyConcept}
\textbf{Adequate causality:} if $x$ happens, there is a probability that $y$ has happened because of $x$ . This is not the determinism of natural scientists (``if $x$ then $y$'') . Adequate causality also implies multi-causality---one factor may be adequate among others .

\textbf{Elective affinity:} Weber sees Protestant ethics and rational capitalism as partners which cannot be separated . The affinity we have between hydrogen and oxygen, which is why we have water---a similar kind of affinity exists between Protestant ethics and rational capitalism .
It is a subtle relationship: it is not that one leads to the other, both are interdependent and conjoined, like conjoint twins---two sides of the same coin, a hand-in-hand relation in which you cannot see a causal relation .
\end{KeyConcept}

\paragraph{The Other Factors}
Weber says asceticism is the necessary condition---the other conditions in the ideal types are only sufficient . You cannot be a capitalist if you are not ascetic; you cannot be a capitalist if you are spending on yourself, because you have to reinvest your money . Besides asceticism, other factors responsible were :
\begin{enumerate}
    \item Asceticism---the necessary condition 
    \item Formal laws---without laws there cannot be regulation in industry 
    \item Appropriation of the physical means of production---which Marx believed was the only condition 
    \item Market and trade 
\end{enumerate}

\begin{CriticalPoint}
But trade, market and laws have existed in all parts of the world . Why then only in England and the Netherlands? Because there is something special about the meanings, which are hidden and have to be interpreted by a sociologist . Max Weber is an interpretative sociologist .
\end{CriticalPoint}

\paragraph{Turning Marx Upside Down}
\vspace{6pt}
\noindent
\begin{tabularx}{\textwidth}{>{\hsize=0.3\hsize}Y >{\hsize=0.35\hsize}Y >{\hsize=0.35\hsize}Y}
\toprule
\rowcolor{AccentDark}
\thead{Dimension} & \thead{Karl Marx} & \thead{Max Weber} \\
\midrule
Base of society & The economy is the base . & Religion is the base . \\
\midrule
Position of religion & Religion sits in the superstructure and justifies the economy . & Religion led to the development of the economy . \\
\midrule
Form of causation & Monocausal and determinist . & Adequate causality, multi-causality, elective affinity---a dotted line, not a solid arrow . \\
\bottomrule
\end{tabularx}
\vspace{6pt}

Basically what Weber said is that the economy comes out of religion . This was a critique of Marx: Marx believed religion justified the economy; Weber turned Marx upside down and believed it is religion which is the base that led to the development of the economy .

\begin{QuickRevision}
\begin{itemize}
    \item The spirit of capitalism---the Puritan work ethic---was born in Protestantism, and capitalists reward exactly its virtues: productivity, hard work, punctuality, frugality .
    \item Rational capitalism was an accident of history: the unintended consequence of the Protestant Reformation, with no intention to exploit .
    \item Weber is not monocausal; he claims adequate causality (probabilistic, multi-causal) and elective affinity (conjoint twins, hydrogen and oxygen), not determination .
    \item Asceticism is the necessary condition; formal laws, appropriation of the physical means of production, and market/trade are the other, merely sufficient, conditions .
    \item Weber turns Marx upside down: for Marx economy is base and religion superstructure; for Weber religion is the base .
\end{itemize}
\end{QuickRevision}

\subsection{Why Not India? The Ideal Type of Hinduism}
The next question is why the model of legal-rational capitalism never developed in any other pocket---why it never developed in India . To answer this, the social action of Hindus has to be understood, and the social action of Muslims has to be understood in order to understand why Islam never led to the rise of capitalism . What is so special about Protestantism only, that the spirit of capitalism was absent in every other religion? 

\begin{CriticalPoint}
We study Hinduism rather than Islam or Judaism . Why? India is not only about Hindus---we have Muslims, Jews, Buddhists and Jains .
The reason given: Indian Society/Indian Sociology, Paper 2, is also about Hindus; we have to study the caste system, which is not a system of Islam or Buddhism---these have actually challenged the caste system, yet we unfortunately study caste .
That is why scholars such as Dipankar Gupta and Yogendra Singh call Indian Sociology Hindu Sociology, not just Indian Sociology .
\end{CriticalPoint}

\begin{KeyConcept}
To Max Weber, Hinduism is Indian Catholicism---the same Catholicism against which Protestants developed .
That is why Buddha was taught: Buddha developed against Hinduism just as Luther developed against Catholicism . There is a relation between the Brahminism that we had and the Catholicism that they had; both share almost similar doctrines .
So Weber developed an ideal type of Hinduism as well, saying: I am just an individual, I cannot know Hinduism in totality, but I would highlight certain accentuated features . This is his method: ideal type, adequate causality, elective affinity, verstehen .
\end{KeyConcept}

\paragraph{Elements of the Ideal Type of Hinduism}
\begin{enumerate}
    \item \textbf{Karma-dharma theory}---you are born a Brahmin as the fruit of past karmas; you have become a Shudra because of past deeds and sins committed in a previous life . Karma means your deeds; dharma means your duty .
    \item \textbf{Transmigration of soul}---the soul is a transitory phenomenon . The body is made up of essential elements---earth, fire, air---but it is the soul that decides your fate, and the soul has to be kept pure . Nothing is kept in the body or in worldly pleasures .
    \item \textbf{Belief in rebirth}---since there is belief in the soul, there is belief in rebirth .
    \item \textbf{No aspiration}---Hindus do not develop aspirations, because we are destined to perform this duty according to varna and dharma . Hence there is no concept of social mobility .
\end{enumerate}

\paragraph{The Dharma of Each Varna}
\vspace{6pt}
\noindent
\begin{tabularx}{\textwidth}{>{\hsize=0.3\hsize}Y >{\hsize=0.7\hsize}Y}
\toprule
\rowcolor{AccentDark}
\thead{Varna} & \thead{Dharma as Stated in the Lecture} \\
\midrule
Brahmin & To assist the kings, to give them advice, to educate the masses, to preach religion . \\
\midrule
Kshatriya & To safeguard the interest of the people and to secure them from war of any kind . \\
\midrule
Vaishya & To indulge in merchant activities and to generate wealth . \\
\midrule
Shudra & To serve the upper three varnas . \\
\bottomrule
\end{tabularx}
\vspace{6pt}

A Shudra cannot do what a Vaishya is doing; a Brahmin cannot do what a Kshatriya is doing; a Kshatriya cannot do what a Brahmin is doing . The division of labour has been done through the varna system in India .
Ambedkar calls it not just division of labour but division of labourers, because the labourers themselves have been divided according to their varnas . Your specialisation comes out of your varna, not out of your merit or talent .

\paragraph{Consequences}
The child of a Dalit is born a Dalit and performs what is considered the occupation of a Dalit---it reproduces . The child of a Brahmin gets education at universities; the child of a Kshatriya performs a similar kind of occupation to his father . Aspirations are not high; they may be a bit different but stay in the same range .
If you do not have aspiration and ambition, you are not going to work beyond your varna . And you also have belief in the afterlife---so what would you do with money? 

\begin{ExampleBox}
One ancient Indian philosopher, Charvak, propagated a materialist philosophy: earn money, even if you have to steal, and live your life with full luxury---whether it comes by theft, by loot or by plunder . That is Charvak's materialism .
\end{ExampleBox}

\begin{CriticalPoint}
With such beliefs, can we ever develop the spirit of capitalism? We wait for Sunday; Monday feels like a curse---there are even songs about it .
So we do not have labour commitment---we are not committed to labour . But Protestants are not like this: the more they work, the more they feel they are getting closer to their destination, that is heaven .
They enjoy and feel pleasure in the work; we curse work and are concerned only about the paycheck . We are more oriented to the party, not to the work . This is why Weber believed Hinduism is Indian Catholicism .
\end{CriticalPoint}

\begin{QuickRevision}
\begin{itemize}
    \item Weber studies Hinduism because Indian sociology itself is largely built on caste (Dipankar Gupta and Yogendra Singh call it Hindu Sociology) .
    \item Hinduism = Indian Catholicism; Weber constructs an ideal type of it too .
    \item Four elements: karma-dharma theory, transmigration of the soul, belief in rebirth, absence of aspiration and hence of social mobility .
    \item Varna dharma fixes occupation---Ambedkar calls it division of labourers, not division of labour; specialisation follows varna, not merit .
    \item Result: no labour commitment; work is cursed rather than experienced as spiritual, so the spirit of capitalism could not arise .
\end{itemize}
\end{QuickRevision}

\subsection{Answer-Writing Guidance and Doubts from This Class}

\begin{PYQLink}
What you have to do is understand the ideal type of Protestantism---known as the Protestant work ethic---and the ideal type of rational capitalism, and then draw a relation between these two .
In the answer you simply explain these ideal types and draw the relation . You do not have to explain all the history---Martin Luther, John Calvin, Benjamin Franklin . That history is only for your understanding, and it will help you in other units .
The history-versus-sociology distinction here is itself usable: the narrative of Luther and the Church is history; the interpretation of meanings is sociology .
\end{PYQLink}

\paragraph{The Printing Press: A Point Added at the End}
\begin{GovtPolicy}
Before Martin Luther there were many people who had challenged the Catholic Church . Why then is Martin Luther singled out? Because Martin Luther took advantage of the printing press .
The printing press was a very recent technology . By translating the Bible into German, Luther had two lakh copies published---and to publish two lakh copies in the 1530s was a big thing . They were colourful copies, not black and white, and people purchased them .
The printing press as a technology actually led to awareness about the Bible; without it people would not have known what the Catholic Church does and what actually lies in the Bible .
Therefore technology is an agent of social change, and religion in itself is an agent of social change---it was only because of Protestantism that capitalism took shape, so Protestantism had an adequate causal relation with the rise of the industrial revolution .
\end{GovtPolicy}

\paragraph{Were Protestants Only in England and the Netherlands?}
Largely, yes . Just as Hinduism is largely in India---we say largely . In the social world there is nothing like ``only'' . Some Protestants exist elsewhere too; we have to look at what is largely the case .

\paragraph{Is Weber a Pseudoscientist, Since All His Theories Are Irrefutable?}
You cannot call Weber a pseudoscientist, because Weber never said ``this is why capitalism developed'' . He never established a determinist relationship; he is not a monocausal scholar . And if he is not a determinist, how can you call him a scientist or a pseudoscientist in that sense? 
But you can definitely still call him a scientist: his research was very rigorous, aimed at understanding social action in an objective fashion, forming ideal types of reality and studying reality through models---which is what scientists do---while maintaining objectivity and not letting his values come in .
He clearly says sociology is a science of social action in so far as it is concerned with its causes and consequences .

You may call Marx a pseudoscientist, because he was a determinist---although when people criticised Marx as an economic determinist, Engels came forward and defended the thesis, explaining why he is not .

\paragraph{According to Weber, Was Capitalism Not Exploitative?}
The question of exploitation does not arise at all . Weber is looking at capitalism in a very different fashion .
When you say capitalism is exploitative, you are looking at capitalism through the Marxian perspective---not in an objective fashion . Remove those goggles, put on Weber's goggles, and you see a different reality . Put on Durkheim's glasses and you see a very different idea of capitalism; put on Adam Smith's glasses and capitalism appears different again .
There is no absolute reality in the social world---these are all perspectives .

\begin{QuickRevision}
\begin{itemize}
    \item For the answer: explain the ideal type of Protestant work ethic and the ideal type of rational capitalism, then draw the relation; leave the history out .
    \item Luther is singled out because he used the printing press---two lakh colourful German Bibles in the 1530s; technology and religion are both agents of social change .
    \item Protestantism was ``largely'', not ``only'', in England and the Netherlands---nothing in the social world is ``only'' .
    \item Weber is a scientist but not a determinist, so the charge of pseudoscience does not apply; Marx, being determinist, is more exposed to it, and Engels defended him .
    \item Exploitation is a Marxian category; for Weber capitalism is wealth accumulation, and there is no absolute reality in the social world---only perspectives .
\end{itemize}
\end{QuickRevision}

\sectiondivider

\section{Critique of Max Weber and the Comparison of Marx, Weber and Durkheim}

\subsection{Placing the Protestant Ethic at the Centre of Weber's Work}
In Weber we have read social action, ideal type, authority, bureaucracy, and finally the \textit{Protestant Ethic and the Spirit of Capitalism} . In the \textit{Protestant Ethic} we can see all of them---bureaucracy, authority, ideal type, verstehen and social action . These were the rivers which were destined finally to meet the ocean .

\begin{KeyConcept}
Write it in the past tense: work was a religious or spiritual experience to Protestants .
If work is a spiritual experience to you as a Protestant, and you are working continuously, relentlessly, incessantly in a non-stop fashion because you think you are engaged in a spiritual experience, you are bound to create a system called capitalism .
When you are working you are not working to earn money---it is like going to a temple . This kind of belief is very unique and not found just anywhere in the world .
They were working out of passion, not to earn money, but to secure their seat in heaven .
\end{KeyConcept}

\begin{QuickRevision}
\begin{itemize}
    \item The \textit{Protestant Ethic} is the point at which social action, ideal type, authority and bureaucracy all converge .
    \item Work was a religious/spiritual experience to Protestants---worship, not wage-earning .
    \item Continuous work driven by that belief is what produces the system called capitalism .
\end{itemize}
\end{QuickRevision}

\subsection{Catholicism, Protestantism and Capitalism: The Three-Box Comparison}
Draw three boxes: Catholicism, Protestantism and Capitalism . What is written in the first box will not match the third; what is written in the second will match the third fully . That is where you see the congruence, correspondence and affinity between the two models .

\vspace{6pt}
\noindent
\begin{tabularx}{\textwidth}{>{\hsize=0.28\hsize}Y >{\hsize=0.36\hsize}Y >{\hsize=0.36\hsize}Y}
\toprule
\rowcolor{AccentDark}
\thead{Question} & \thead{Catholicism} & \thead{Protestantism} \\
\midrule
How is salvation attained? & By rituals and by purchasing indulgences . & By predestination---the elect have already been chosen; no indulgence, no baptism, no ritual is needed . Complete rejection of the Catholic path . \\
\midrule
How long should one work? & Work until your immediate needs are met . & Calling---keep working . Can you stop worshipping once you feel satisfied? Worship is something you engage in forever and never feel satisfied with . Work is a spiritual experience . \\
\midrule
What is sin? & Profit is sin . & Idleness is sin . Sitting idle is the sin; earning is not . \\
\midrule
What should be done with wealth? & Distribute your wealth to the poor . & Ascetic materialism---asceticism yes, but manifested by investing money in something productive . Distributing a crore among the poor is not productive; they will come back tomorrow when it is finished . \\
\bottomrule
\end{tabularx}
\vspace{6pt}

\paragraph{Why the Protestant Was the Productive Worker}
If you are every time engaged in worship, always thinking that a sin has been committed, will you ever concentrate on work? Are you a productive worker to the capitalist? No .

The Protestants believed they had already been chosen, but did not know who among them were the chosen ones---hence salvation anxiety .
They therefore felt that the more wealth they generated, the more their seat was secured . In wealth they saw the grace of God and the opportunity to go to heaven . So wealth was not generated out of greed; it was generated out of their wish to secure their seat in heaven .
Since the wealth comes from work, they were productive to the capitalist class which needs hard-working workers .

Calling and predestination were values that gave the capitalist class productive workers . Because idleness is sin, if the boss said work one more hour, they would happily do it and would not complain . In the work they saw God .
So: maximum hours given to work; and whatever was earned was largely invested in productive things---that is reinvestment and productivity, and that is precisely the spirit of capitalism .

\begin{KeyConcept}
The spirit of capitalism = hard work, efficient people, maximum hours of work, not sitting idle, reinvesting, being productive .
Weber clearly saw that so many things match: it seems as if each was made for the other .
And Hinduism has no such relation with capitalism at all---Hinduism is nothing but Indian Catholicism, so capitalism will never come from it, according to Weber . We do not have the labour commitment these people have; it is their belief that makes them different from the rest of the world---people like Bill Gates, Warren Buffett, McDonald's .
\end{KeyConcept}

\begin{QuickRevision}
\begin{itemize}
    \item The Catholic box contradicts the capitalist box; the Protestant box matches it point for point .
    \item Salvation by predestination replaces rituals; calling replaces working-till-needs-are-met; idleness replaces profit as sin; ascetic materialism replaces charity to the poor .
    \item Salvation anxiety converts the search for a sign of election into unlimited productive labour and reinvestment .
    \item The spirit of capitalism: hard work, efficiency, maximum hours, no idleness, reinvestment, productivity .
\end{itemize}
\end{QuickRevision}

\subsection{Why Not the Jews, Muslims and Sikhs?}
A student asked: Jews were also famous as hard-working people, so why did capitalism not rise there? 
\begin{itemize}
    \item \textbf{Jews}---the problem is that they invested heavily in lending . Most of the wealth they accumulated came out of money-lending practices . When you lend, your practice is to lend money, wait for it to garner interest and then collect it back---so you are sitting idle most of the time, and you are not able to think beyond that one given task of generating wealth . That is what pariah capitalism was, and Weber gives the example of the Jews for it .
    \item \textbf{Islam}---in Islam you will see many things which are haram, forbidden, and those forbidden things do not let the same spirit of capitalism develop, according to Weber . Islam, he says, practises hedonism; luxury is cursed, and the participation of women in the work sphere is cursed . This is not offered as the reality---this is Max Weber's ideal type of Islam .
    \item \textbf{Sikhism}---if you look into the practices of Sikhs, most of the wealth they generate is donated in langar or in charity . During the lockdown it was Sikh people who helped homeless people . But when you donate most of your wealth in charity, you are not an ascetic materialist; this is not what a Protestant would do .
\end{itemize}
A Protestant would instead invest the same amount in corona bonds---Covid bonds, as happened in Europe---rather than sending money to a relief fund . We Indians will do the latter because we are emotional; Protestants will never do these kinds of things .

\begin{QuickRevision}
\begin{itemize}
    \item Jews: money lending = pariah capitalism; lending involves sitting idle and cannot generate the work-ethic spirit .
    \item Islam: forbidden categories and what Weber's ideal type calls hedonism---cursing luxury and women's work participation block the spirit .
    \item Sikhism: wealth given in langar and charity is not ascetic materialism, because charity is not productive investment .
    \item The contrast is bonds versus relief funds---productive reinvestment versus emotional donation .
\end{itemize}
\end{QuickRevision}

\subsection{From Calling to Compulsion: Secularization}
If work to a Protestant is a spiritual experience, what kind of social action is being performed? You will see traditional, affectual and value-rational social action---but you will not see \textit{zweckrational} action .
You are not doing work because you have to make so much wealth and have investment plans, a bungalow, a car, a foreign trip, or ambitions to become CEO . These calculations do not go into it .
It is purely a spiritual experience, and when you are manifesting spiritualism there is no calculation of means and ends . You are simply spiritual without being utilitarian . Where work is a spiritual experience, there is no space for utilitarianism .

\begin{KeyConcept}
Puritans means the members of the Protestants---Puritans and Protestants are synonyms .
Weber says, write it down: \textbf{``Puritans used to work in calling. But today they are forced to do so.''} Everything is hidden in this line .
\end{KeyConcept}

\paragraph{What the Line Means}
The Protestants were simply expressing their spiritual feeling . They saw work as a spiritual experience and kept on working because to them work was nothing but worship .
Ultimately the values required for the generation of rational capitalism were visible in the Protestants, and that is why rational capitalism developed in England and the Netherlands---because of their productive expenditure and investment practices, their hard work, and their treatment of idleness as sin .

But what happened after that? The money that was a sign for you that you are going to go to heaven had now become an end in itself . The same money which was a token, a grace of God, had become an end .
Now you were working not to manifest spirituality, not working in calling, but to generate wealth---even if that generation of wealth comes at the cost of the exploitation of others .
Gradually man became greedy, because he saw that people were giving him so much respect and status . So he wanted more and more wealth, and wealth generation became completely a manifestation of \textit{zweckrational} social action . From there on there were no more values, no tradition, no affectual social action with God---everything was with the money .

\begin{ComparisonBox}
Here Max Weber is shaking hands with Karl Marx . If you remember, Marx said that initially we had $C_1\text{--}C_2$, then $C_1\text{--}M\text{--}C_2$, and finally in capitalism we have $M_1\text{--}C\text{--}M_2$---money becoming the end .
The two are very close to each other when Weber makes this statement .
\end{ComparisonBox}

\begin{KeyConcept}
Write down one line: Max Weber believed the money which was a sign of holy grace had now become an end in itself .
And therefore it was no longer value-rational social action; rather it was completely \textit{zweckrational} social action .
It was no longer substantive rationality; rather it was completely formal rationality .
\end{KeyConcept}

\begin{ExampleBox}
In simple language: you are a person who worships God every day before opening your shop . On the opening day you conducted the rituals, fed the priest, did everything . Every day you worshipped God, and you treated customers as manifestations of God .
But ultimately, when you started generating wealth, it stopped mattering if a customer went away unhappy . Everything is forgotten; the idol gathers dust and it makes no difference . Money had now become an end in itself .
\end{ExampleBox}

\begin{KeyConcept}
This is what Max Weber called the advent of \textbf{secularization} .
Secularization is a process---hence the ``-ization''---characterised by a decline in religious beliefs or values .
If money-making has become an end in itself, then religious beliefs are declining . That is why Weber says Puritans used to work in calling, and today they are forced to do so: the concept of calling is finished .
\end{KeyConcept}

\begin{QuickRevision}
\begin{itemize}
    \item Work as spiritual experience is value-rational, affectual and traditional---never \textit{zweckrational}, because there is no means-ends calculation .
    \item Puritans = Protestants; Weber's line: ``Puritans used to work in calling; today they are forced to do so'' .
    \item Money that was a sign of holy grace became an end in itself; value-rational action became \textit{zweckrational}, substantive rationality became formal rationality .
    \item Weber here converges with Marx's $C\text{--}M\text{--}C$ to $M\text{--}C\text{--}M$ formula .
    \item The outcome is secularization: a process characterised by the decline of religious beliefs and values .
\end{itemize}
\end{QuickRevision}

\subsection{Micro to Macro: Weber as a Substantive Sociologist}
Is Weber studying social action or structures? Is he a micro scholar, or is he also talking about macro structures? 

We started our journey with the social action of individuals---a small man and the attempt to understand his meanings through verstehen .
Gradually, Weber says that a large number of social actions of people resulted in the formation of social structures . If millions are working with almost the same belief, it leads to the rise of capitalism---and these millions were in England and Holland, that is the Netherlands .
So individuals worked to create capitalism, and now that same capitalism is forcing those individuals to generate wealth . It was their action which led to the rise of capitalism; now capitalism is forcing them . Capitalism as a structure is no longer in their hands .

\begin{CriticalPoint}
When you read only that second statement in isolation, Weber sounds just like a structuralist---like Durkheim on social facts: external to and coercive of actors .
Capitalism becomes external to and coercive of actors; these are social facts which have now assumed a reality and are constraining individuals . You would not be able to think beyond generating wealth; wealth-making has become the end and spirituality has become a token symbol .
The same applies to bureaucracy: when you become a bureaucrat you will be bureaucratized . Rules and regulations will constrain you; you will not be able to help an individual beyond the given rules, and you will feel the constraining power of these structures .
\end{CriticalPoint}

\begin{KeyConcept}
Weber used to tell us that the subject matter of sociology is the study of social action . But then it became clear that it is the study of social action, its causes and its consequences .
What are the consequences of social action? Large-scale social structures . Capitalism is the result of a large number of people manifesting initially value-rational and finally \textit{zweckrational} social action; it has assumed a reality of its own type and now exerts control over individuals .
\end{KeyConcept}

\paragraph{Bureaucracy in the Same Line}
The social action of a large number of people led to the rise of bureaucracy . In bureaucracy---impersonality, hierarchy, division of labour, merit-based recruitment---everything becomes impersonal, rule-bound, legal, rational and calculable, all to ensure efficiency .
In the end it is efficiency that has become an end in itself . And this efficiency was actually coming from the Protestants, who were deemed the most efficient people because they worked hard without complaint and treated idleness as sin . They would do whatever was necessary to be seen as efficient, because they wanted to secure their seat in heaven---salvation anxiety made everyone want to be as efficient as possible .

\begin{KeyConcept}
Weber concluded that rational bureaucracy---the modern form, as against the patrimonial bureaucracy of ancient times---and rational capitalism are the result of the social action of Protestants .
The bureaucratic structure we adopted came from Europe; the Indian civil service came from Europe; and in Europe it came from the Protestants .

Therefore capitalism is just one part of the larger process of rationalization . It is not so important to study capitalism as it is to study the process of rationalization .
To study capitalism as a mode of production, and to regard the world as nothing but the evolution of modes of production, would be a very myopic study .
Capitalism is the rationalized form of the economy; bureaucracy is the rationalized form of law; modern education is the rationalized form of education . The larger process to be studied is rationalization .
\end{KeyConcept}

\begin{QuickRevision}
\begin{itemize}
    \item Weber begins with micro social action and moves to macro structures---this is why his sociology is called comprehensive .
    \item Large numbers of similar social actions create structures (capitalism, bureaucracy); the structures then become external to and coercive of the actors, sounding Durkheimian .
    \item Subject matter of sociology therefore = social action, its causes AND its consequences (large-scale structures) .
    \item Efficiency becomes an end in itself, and its source is the Protestant work ethic .
    \item Capitalism is only one part of the larger process of rationalization---which is the real object of study; bureaucracy is rationalized law, modern education rationalized education .
\end{itemize}
\end{QuickRevision}

\subsection{Critique of the Protestant Ethic Thesis}
What follows are not really critiques; they are largely misinterpretations of Weber, because Weber himself never claimed what they deny .

\paragraph{1. R. H. Tawney}
R. H. Tawney is an economic historian . He says he can never be certain that Protestantism led to the rise of capitalism . He believes the spirit of capitalism was absorbed by the Protestants---it is not that the work ethic of the Protestants gave rise to the spirit of capitalism; rather, the spirit of capitalism made them work efficiently .

Where then does the spirit of capitalism come from? Tawney says its origin is not Protestant ethics but the Reformation, the Renaissance and the Scientific Revolution---largely the scientific revolution .
In the Reformation he highlights not the Protestants but the concept of rationality: when Martin Luther challenged the Catholic Church, that shows a rational being---someone who can challenge the church rather than simply obey it because it has been done for so long . He was not performing traditional social action; he performed formal rationality, fully conscious: why should you people act as intermediaries between men and God? 
So Tawney says it is the spirit of rationality that gives rise to the spirit of capitalism, and that rationality comes from the Reformation, along with the scientific revolution and the Renaissance .

He may sound like a critic, but he is not---because Max Weber never said Protestant ethics is the reason behind the rise of capitalism . He only drew adequate causal relations: there is a probability that Protestant ethics may have led to it . He is not being determinist, so the question of criticism hardly arises .

\paragraph{2. Kurt Samuelson}
Kurt Samuelson says we cannot really say it was Protestant ethics that led to the rise of capitalism, or that there is even a probable relation between the spirit of capitalism and Protestants---because he can see the spirit of capitalism even in countries which did not have Protestants .
\begin{itemize}
    \item \textbf{Belgium:} the chocolates which are so famous today have been famous for a very long time . Belgium is largely a Catholic country; there are not that many Protestants there . Yet Belgium, despite being a Catholic country, had the spirit of capitalism .
    \item \textbf{Switzerland:} here there were Protestants, but capitalism developed only in very late stages---at least a hundred years later . He is not saying capitalism did not develop; he is saying it took them a hundred more years to develop that spirit . So you cannot really say there is adequate causality .
    \item \textbf{Netherlands:} to Samuelson, capitalism in the Netherlands had developed even before the Protestants . To him the Netherlands was a Catholic country, not a Protestant one---so it is Catholics who had the adequate relation with capitalism .
\end{itemize}
But this depends on what kind of ideal type each has drawn: the ideal type of Catholicism drawn by Max Weber and the one drawn by Kurt Samuelson are very different . And here the role of values comes in---Max Weber was himself a Protestant, his mother being a Calvinist, and Calvinism is a variant of Protestantism .

Samuelson also says the biggest critique of Max Weber should be that he was not objective: when everything is a probability, can his thesis be considered scientific? Where there is no objectivity, is there any scope for falsifiability? 

\paragraph{3. Reinhard Bendix}
Bendix makes a very interesting observation about Japan . Japan lost the Second World War, but that loss evoked such a great spirit of nationalism and patriotism in the Japanese people that they continuously worked hard, relentlessly, unyieldingly and incessantly---and this actually led to the rise of Japanese capitalism . That is why Japanese electronics is the most unique in the world .
And who were the Japanese? They were Buddhists . This shows that Buddhism can be practised in the private sphere: at home you are a Buddhist, but the moment you go out into the industry you become a hard-working individual, a Japanese nationalist who will do anything for productivity or the wealth of the nation .
So capitalism existed without Protestantism . According to Bendix, to say that Protestantism led to capitalism, or even had an adequate relation with it, is not so right .

\begin{ExampleBox}
On Japanese capitalism (to be taken up in the chapter on Work and Economy): Japan practises job rotation---a worker may become a manager, a manager may become a clerk, a manager may for a few days become CEO . Everybody becomes everything .
Unlike traditional capitalist enterprises where you are permanently a manager or permanently a clerk, there is no such strict compartmentalization in Japan and therefore no scope for alienation .
Japan also has a maximum aged population and lacks the demographic dividend that India celebrates . When the elderly do not feel integrated with the young, the suicide rate rises .
\end{ExampleBox}

\paragraph{4. Milton Singer}
Milton Singer conducted a study in Madras, India, and found a process called \textbf{compartmentalization} .
Compartmentalization means the private sphere is different from the public sphere: at home you are completely religious, but the moment you go to the factory you are no longer religious and are concerned only about wealth . That does not mean your religiosity is over---you become religious again when you get back home .

In South India there are more temples than in the whole world---after every hundred metres you will see a temple . This shows that religiosity can exist despite capitalism .
Weber says capitalism has come and religiosity is finished, money has become an end in itself . But Milton Singer says they are concerned with money at the office and become religious again when they come home .
In simple terms, they lead life in two compartments: in private life they are religious, in public life they are businessmen . This is not the opposite of what Weber says; it is simply different from it .

\vspace{6pt}
\noindent
\begin{tabularx}{\textwidth}{>{\hsize=0.25\hsize}Y >{\hsize=0.35\hsize}Y >{\hsize=0.4\hsize}Y}
\toprule
\rowcolor{AccentDark}
\thead{Critic} & \thead{Core Claim} & \thead{Evidence Cited} \\
\midrule
R. H. Tawney & The spirit of capitalism was absorbed by Protestants, not produced by them . & Origin lies in Reformation, Renaissance and above all the Scientific Revolution; Luther's challenge itself displays rationality . \\
\midrule
Kurt Samuelson & No adequate causality between Protestantism and capitalism . & Catholic Belgium had the spirit; Protestant Switzerland developed capitalism a hundred years late; the Netherlands had capitalism before Protestants and was Catholic . \\
\midrule
Reinhard Bendix & Nationalism, not Protestantism, can produce capitalism . & Buddhist Japan's post-war nationalism produced Japanese capitalism . \\
\midrule
Milton Singer & Religiosity and capitalism coexist through compartmentalization . & Study of Madras, India---religious at home, profit-oriented at the factory; the density of temples in South India . \\
\bottomrule
\end{tabularx}
\vspace{6pt}

\begin{CriticalPoint}
A student asked whether the Protestant is similar to the Buddhist who rose against ``Indian Catholicism'' . Both came in reaction to a larger religion---Catholicism and Hinduism respectively---but their values are entirely different .
For Buddhism, the material is the source of suffering . For the Protestant, material wealth is not suffering at all---it is a sign that he is going to go to heaven . For one it is the source of suffering, for the other the source of heaven .
\end{CriticalPoint}

\begin{QuickRevision}
\begin{itemize}
    \item Four critics: Tawney (spirit absorbed, origin in rationality from Reformation/Renaissance/Scientific Revolution), Samuelson (Belgium, Switzerland, Netherlands; objectivity and falsifiability), Bendix (Buddhist Japan's nationalism), Singer (compartmentalization in Madras) .
    \item None of these is strictly a critique, because Weber never asserted a determinist or monocausal link .
    \item Samuelson's deeper point: differing ideal types of Catholicism, and Weber's own Calvinist background, raise the question of value-freedom .
    \item Buddhism and Protestantism are structurally parallel as protests but opposite in values---material as source of suffering versus material as sign of salvation .
\end{itemize}
\end{QuickRevision}

\subsection{Critique of Bureaucracy}
Again, these are not really critiques, because whatever the critics say has already been said by Max Weber himself .

\paragraph{1. Robert Merton---Dysfunctions of Bureaucracy}
\begin{enumerate}
    \item \textbf{Rules become an end in themselves:} Bureaucracy---in government offices and in the private sector alike---exists to ensure efficiency and to help the client . But when the rules made for efficiency become their own end, it leads to the displacement of the goal for which bureaucracy was made . Bureaucracy was made to serve people; you made rules, but now you are more concerned about the rules, and the rules act as obstruction . This is what we call \textbf{red tapeism}---from the red tape used on file holders since colonial times .
    \item \textbf{There is no scope for improvisation:} Once the rules are made, it takes at least an era for reforms in those rules . If you suggest any change to UPSC, it will take fifty or sixty years; committees will be made, members will die, new members will be appointed, and the report will be delayed .
    \item \textbf{Impersonality in procedures may lead to friction between public officials and the public:} Clients expect a certain empathy from bureaucrats, but bureaucrats are engaged in being highly formal, rational and in what Max Weber calls \textbf{cold calculations}---calculations which lack empathy or feeling for human beings . All the feelings are for the rules . The closer you are to the rules, the further you are from people; and this is what you will be trained for .
\end{enumerate}

\begin{CriticalPoint}
Is this really a critique? Weber has already said all of it, which is why Marx, Weber and Durkheim are classical scholars---nobody can criticise them, because they have said everything .
Weber was ambivalent: one opinion was that these features have helped bureaucracy become efficient; the second was that the very same features which made bureaucracy efficient will make it inefficient .
When Marx says we will have a better world in the name of socialism with the dictatorship of the workers, Weber says: the more you advance, the more society gets rationalized and bureaucratized, and the more it is the rule of rules . Humans would not be ruling; rules would rule .
He called this not the dictatorship of the workers but the \textbf{dictatorship of the officials} .
\end{CriticalPoint}

\paragraph{2. Michel Crozier}
\begin{enumerate}
    \item \textbf{Bureaucracy refuses to learn from its own mistakes:} It is rule-bound and does not change . The rules, hierarchy and chain of command delay projects: the plan was to make a highway within five years, and the outcome was a highway that took fifty years with the expenditure rising from one million dollars to one billion .
    \item \textbf{Everybody's responsibility is no one's responsibility:} You go to an office to get your work done, and the person there says he does not know---as satirised in the popular television show \textit{Office Office} .
    \item \textbf{Bureaucracy is engaged in dealing with its own problems and therefore does not get time for public welfare:} ACR reviews remain pending, disputes between officers occupy the system, corruption and appointment issues---the appointment of the wrong person, the overriding of the seniority principle---are its own internal problems .
\end{enumerate}

\begin{GovtPolicy}
The counter offered in class: today we have Digital India, and by being digital we are trying to be more efficient and to cut out all the intermediaries between people and officials .
Middlemen are removed, expenditure is reduced, and we have schemes that transfer the beneficiary amount directly into the beneficiary's account, so the welfare money that used to drain away on the route is saved .
\end{GovtPolicy}

\paragraph{3. Alvin Gouldner}
Alvin Gouldner is the same person who gave the idea of reflexive sociology .
He conducted a field study in America among gypsum plant mine workers . He found that those people were doing their work with full sincerity, discipline and punctuality until one day a manager was appointed to oversee them .
The moment a manager was appointed to monitor them, their efficiency got reduced . So you cannot say that bureaucracy ensures efficiency; rather, bureaucracy can also decrease efficiency .

The parallel drawn: why do some people not study at home and instead come to Delhi? Because they are not comfortable studying at home---they believe the lesser the monitoring, the better . If you appoint managers to monitor someone, they will not be efficient .

\paragraph{4. Jerry Pournelle---The Iron Law of Bureaucracy}
Go to any organization and you will meet two kinds of people: those who are working for the organization with full dedication, and those who are working in the name of the organization .
The second kind boast about the work they do and present themselves as the ones with maximum contacts in the organization; every time they say what they have done for the organization .
Unfortunately these people become visible, and the efficient, hard-working, dedicated people become invisible . And bureaucracy rewards visibility .
That is why you see these people climbing the ladder like anything . So bureaucracy does not reward efficiency---bureaucracy rewards visibility, not efficiency . That is the iron law of bureaucracy .

\paragraph{Mechanical and Organic Models of Bureaucracy}
\vspace{6pt}
\noindent
\begin{tabularx}{\textwidth}{>{\hsize=0.25\hsize}Y >{\hsize=0.38\hsize}Y >{\hsize=0.37\hsize}Y}
\toprule
\rowcolor{AccentDark}
\thead{Dimension} & \thead{Mechanical Model} & \thead{Organic Model} \\
\midrule
Character & Rigid, hierarchical, command-based---the pure Weberian model . & Participative; decisions are not taken only at the top . \\
\midrule
Decision-making & Top-down approach: they give you an order, you follow it . & Participation in decision-making; opinions are taken before deciding, for instance on the timing of a class . \\
\midrule
Example given & Delhi University . & A few (not all) private-sector organisations . \\
\midrule
Outcome & Fails to deliver---degrees still pending years after the batch has passed out . & Brings in more efficiency . \\
\bottomrule
\end{tabularx}
\vspace{6pt}

\begin{QuickRevision}
\begin{itemize}
    \item Merton: rules become ends in themselves (goal displacement, red tapeism), no scope for improvisation, impersonality causing friction and cold calculations .
    \item Crozier: bureaucracy refuses to learn from its mistakes; everybody's responsibility is nobody's responsibility; it is busy with its own problems instead of public welfare .
    \item Gouldner: the gypsum plant study---appointing a monitoring manager reduced worker efficiency, so bureaucracy can decrease efficiency .
    \item Pournelle's iron law: bureaucracy rewards visibility, not efficiency .
    \item Weber pre-empts all of it through his ambivalence---the features that make bureaucracy efficient are the ones that make it inefficient---and his prediction of the dictatorship of the officials .
\end{itemize}
\end{QuickRevision}

\subsection{Critique of Verstehen}
Weber says the method is empathetic understanding---to understand social action in its historical context . Any action that is meaningful, oriented to other actors and performed in a certain historical context is social action . The question is: how are you going to interpret it? 

Weber was a Protestant, and that is why he could interpret that to Protestants work is a religious experience .
A tribal person could not interpret a Protestant's work like this . For a tribal person, spiritual experience is something else---the worship of natural forces: fire, air, earth, the forest . This they call animism and naturism; to them that is not religion in our sense .
Imagine planting a person from the Andaman and Nicobar Islands beside you as you work in a factory . Would he be able to interpret what you are doing? First, he would not even know that this is ``work'', because he is not from an industrial society and has lived in a tribal society forever . Second, that it might be a way of worship is a far more distant thought .

\paragraph{1. Crisis of Representation}
Who is going to understand the social action, and how is he going to maintain value neutrality?  How can he be objective in the way you expect him to be? These are very subjective things .
To put yourself into the shoes of another person, to think about what he is doing while taking into consideration the history of that person, is not an easy task . It is as difficult to do as it is easy to say .

\paragraph{2. Intersubjectivity}
The word was coined by Alfred Schutz, who is highly impressed with Max Weber but also criticises him .
Schutz says he is highly thankful to Max Weber, because it is only because of Weber that we started studying individual actors and the meaning of the action . Weber is therefore indispensable for the rise of phenomenology .

\textbf{Phenomenology:} there is nothing objective in the social world; these are webs of meanings .
What is a pen to you may be a stick or an earbud to someone else . A chair is not just a chair---turned over, its legs look like a defence shield, and a tribal who has never seen one may use it for defence rather than for sitting .
Capitalism is different to Marx, different to Weber, different to Adam Smith, different to Durkheim . The colour red may imply different meanings to different people . A thumbs-up implies different things in different cultures---in Russia it carries a very different and vulgar meaning; in India something else; in America and Japan something else again .

So there is no objective meaning . Since there is intersubjectivity, reality is highly subjective from one individual to another . And if your understanding of social action differs from someone else's, how would you ever understand social action? If nothing is objective, can we ever formulate laws? We can never formulate laws .
At the same time Schutz was very happy that someone tried to study individuals, because scholars before Max Weber were studying only structure . From here comes phenomenology, and Schutz becomes interested in studying these different intersubjective meanings .

\begin{CriticalPoint}
Post-modernism is an extension of phenomenology . Post-modernists say nothing is objective in this world; theories are highly fluid and not so applicable; Marxism has failed; science does not exist as an objective enterprise but is completely subjective, falsifying its own claims and coming up with another claim .
There is no objective truth even in science . Science claims and projects itself as a grand narrative---but there is no grand narrative: no Marxism, no science, no functionalism .
\end{CriticalPoint}

\begin{QuickRevision}
\begin{itemize}
    \item Verstehen = empathetic understanding of social action in its historical context; Weber's own Protestant background made his interpretation possible .
    \item Crisis of representation: who interprets, and how is value neutrality and objectivity maintained when standing in another's shoes? 
    \item Intersubjectivity (Alfred Schutz): meanings differ from person to person and culture to culture (pen, chair, thumbs-up, capitalism itself), so laws cannot be formulated .
    \item Schutz nevertheless credits Weber as indispensable to phenomenology, of which post-modernism is a further extension .
\end{itemize}
\end{QuickRevision}

\subsection{Comparison I: Paradigm, Tradition and Methodology}
First, let us locate the three scholars . Paradigms are sets of beliefs . In natural science we have a dominant paradigm---today what Einstein said is the dominant paradigm, so gravitation must be explained through Einsteinian physics; Newtonian physics is outdated . Tomorrow it may be replaced by a new paradigm .

\begin{CriticalPoint}
But in social science there is no dominant paradigm . No one is superior . You cannot say Marx is superior to Weber, or Weber to Durkheim, or Durkheim to Merton, or Merton to Talcott Parsons or J{\"u}rgen Habermas or Horkheimer .
All of them have looked at the phenomenon through a perspective . When Marx talks about capitalism it is a perspective he is showing, not objective reality---if capitalism really were so exploitative, it would have been over by now .
\end{CriticalPoint}

\vspace{6pt}
\noindent
\begin{tabularx}{\textwidth}{>{\hsize=0.22\hsize}Y >{\hsize=0.26\hsize}Y >{\hsize=0.26\hsize}Y >{\hsize=0.26\hsize}Y}
\toprule
\rowcolor{AccentDark}
\thead{Dimension} & \thead{Durkheim} & \thead{Marx} & \thead{Weber} \\
\midrule
Paradigm & Structural functionalism, which itself falls under positivism . & Conflict paradigm---Marxism . & Interactionism, also known as interpretativism . \\
\midrule
Why & Looks at how institutions help maintain society; does not look at conflict . & Does not see institutions working for society but for a class; sees only two opposing groups---bourgeoisie and proletariat, masters and slaves, lords and serfs, one class oppressing another with the help of institutions . & The subject matter of sociology to him is the study of social action; he studies the action of individuals, their meanings and motives, not conflict (though conflict returns in social stratification) . \\
\midrule
Subject matter & Social fact . & Class conflict / class struggle . & Social action . \\
\bottomrule
\end{tabularx}
\vspace{6pt}

\begin{KeyConcept}
What is positivism? Whether water is in India or in America, water is $\text{H}_2\text{O}$ only---it is the same everywhere, so you can theorise how water is formed . That is a general law .
But water in India may be chemically similar to water in America, while an Indian is not similar to an American . If men are not similar, how can you theorise? Positivists nevertheless believed individuals are similar---every individual is just like matter .
Like apples falling in reaction to the law of gravitation, individuals marry in response to the need of society . That is positivism .
\end{KeyConcept}

\paragraph{Sociological Tradition}
\begin{itemize}
    \item \textbf{Durkheim}---macro tradition . He talks about structures constraining individuals, coercive of individuals: family as social fact, religion as social fact .
    \item \textbf{Marx}---macro . Larger structures again regulating individuals; society following laws .
    \item \textbf{Weber}---micro to macro . He started with micro social action, then studied structures, and then said the macro regulates the micro . The same macro which is the making of the micro now regulates individuals .
    This is known as \textbf{substantive sociology}---the integration of micro and macro . That means Weber is highly comprehensive: he is not talking only macro or only micro, but looking at both .
\end{itemize}

\paragraph{Methodology}
\vspace{6pt}
\noindent
\begin{tabularx}{\textwidth}{>{\hsize=0.25\hsize}Y >{\hsize=0.25\hsize}Y >{\hsize=0.5\hsize}Y}
\toprule
\rowcolor{AccentDark}
\thead{Scholar} & \thead{Methodology} & \thead{What It Involves} \\
\midrule
Durkheim & Positivist methodology & The social world operates just like the natural world; social facts are treated as things; there is cause-effect in the social world; non-material social facts lead to social facts; generalizations can be developed---for example rapid industrialization leading to an increased suicide rate . Visible above all in the study of \textit{Suicide}, where he is purely positivist, follows the scientific method, looks into cause and effect, makes laws, and uses concomitant variation . In the \textit{Division of Labour} he is not positivist but a philosopher writing a dissertation, and it was later read through the functionalist perspective---division of labour creating organic solidarity and integration, a moral force binding people again after individualism weakened social cohesion . \\
\midrule
Marx & Dialectical methodology & Not simply cause and effect . A thesis is rejected by an antithesis; the antithesis is rejected by another; then comes synthesis, which becomes a thesis and is rejected again . ``Due to inherent contradictions within capitalism, socialism would rise'' seems like cause and effect on the surface---but the inherent contradictions themselves are dialectics . He does also believe in causal relations and that the world is governed by laws . \\
\midrule
Weber & Historical method & Note: not ``historical sociology''---historical method . Hermeneutics, the study of historical texts and the intention of the author, was extended to the study of social action, which must be understood in its historical context . That is why \textit{Protestant Ethic} required full world history---Weber had to read history to understand how Protestants came about . It is only in history that Weber could reveal the reasoning behind the meaning they attribute to work . The historical method also implies verstehen, causality (but multi-causal and adequate, not monocausal), and ideal type . \\
\bottomrule
\end{tabularx}
\vspace{6pt}

\begin{PYQLink}
Careful distinction for the examination: the name of Marx's theory is historical materialism, but the method used to explain it is dialectics .
When you see a question on historical method, write only Max Weber---not Karl Marx .
\end{PYQLink}

\begin{QuickRevision}
\begin{itemize}
    \item No dominant paradigm exists in social science, unlike Einstein replacing Newton in physics; all three offer perspectives .
    \item Paradigms: Durkheim structural-functionalist (within positivism), Marx conflict, Weber interactionist/interpretativist .
    \item Traditions: Durkheim macro, Marx macro, Weber micro-to-macro = substantive sociology, hence comprehensive .
    \item Methodologies: Durkheim positivist (\textit{Suicide}, concomitant variation), Marx dialectical (historical materialism is the theory, dialectics the method), Weber historical method (hermeneutics, verstehen, adequate causality, ideal type) .
    \item For a question on historical method, the answer is Weber alone .
\end{itemize}
\end{QuickRevision}

\subsection{Comparison II: Modernity, Division of Labour and Religion}
All three are the result of, and a response to, the period of great transformation---industrialization, the French Revolution, the Reformation, the Renaissance and the Enlightenment .

\paragraph{Subject Matter of Modernity}
\vspace{6pt}
\noindent
\begin{tabularx}{\textwidth}{>{\hsize=0.25\hsize}Y >{\hsize=0.35\hsize}Y >{\hsize=0.4\hsize}Y}
\toprule
\rowcolor{AccentDark}
\thead{Scholar} & \thead{What Modernity Means to Him} & \thead{Principal Concerns} \\
\midrule
Durkheim & Industrialization---he responds to the force of industrialization, which has weakened cohesion and made people individualist, aspiring and ambitious . He was a conservative who believed traditional institutions had broken down; his question was what would bring solidarity again, and the answer was the division of labour---a qualitatively different, weakened, organic solidarity . & Weakened social cohesion / individualism; anomie; abnormal division of labour . \\
\midrule
Marx & Capitalism---not industrialization per se; capitalism is the focus of his work . & Not cohesion but exploitation of the working class; appropriation of surplus value; alienation of the proletariat . \\
\midrule
Weber & Rationalization---not industrialization or capitalism per se, but the larger process of which capitalism is only a part . & Disenchantment; the absence of any space for emotion; societal alienation . \\
\bottomrule
\end{tabularx}
\vspace{6pt}

\begin{CriticalPoint}
Marx is very pessimistic about institutions, because to him they work for the interest of the capitalist class . Bureaucracy is machinery planted to secure the interest of the capitalist class; so is the state; so is religion; so is law .
Education is based on class privilege---whoever has more money gets better education, and the poor are told they could not clear the system because they lacked merit . But he could not clear it because he did not have the class position: whoever has class will have merit . You need money to go to Oxford .
\end{CriticalPoint}

\paragraph{Division of Labour}
\vspace{6pt}
\noindent
\begin{tabularx}{\textwidth}{>{\hsize=0.33\hsize}Y >{\hsize=0.33\hsize}Y >{\hsize=0.34\hsize}Y}
\toprule
\rowcolor{AccentDark}
\thead{Durkheim} & \thead{Marx} & \thead{Weber} \\
\midrule
Division of labour brings solidarity---organic solidarity, interdependence and integration . & Division of labour is an instrument of exploitation . It justifies why you should get a low wage; and in the name of specialization it alienates you from your creativity, so you cannot explore your potential beyond one domain . Marxist scholars such as Braverman call this \textbf{de-skilling}---labour is not skilled but de-skilled . & Division of labour is there to ensure efficiency in the organization---one of the features of bureaucracy . He is not talking about exploitation or solidarity; he looks at it through an instrumental perspective . But given his ambivalence, he also holds that the same division of labour may restrict efficiency, because concern with rules, hierarchy and process prevents you from helping the client---which is the irrationality of rationality . \\
\bottomrule
\end{tabularx}
\vspace{6pt}

\begin{ExampleBox}
The illustration of de-skilling given from personal experience: on first visiting Delhi University and asking teachers the meaning of certain theoretical concepts, the reply was to ask a different professor, ``because they teach theory---my department is family and kinship, I am a specialist in that'' .
Your potential is not polished; rather you are reduced to one domain in the name of specialization . A PhD, in this Marxist reading, is not a good thing, because it means you now know nothing else .
\end{ExampleBox}

\paragraph{Religion}
\vspace{6pt}
\noindent
\begin{tabularx}{\textwidth}{>{\hsize=0.25\hsize}Y >{\hsize=0.75\hsize}Y}
\toprule
\rowcolor{AccentDark}
\thead{Scholar} & \thead{View of Religion} \\
\midrule
Durkheim & Looks at the functions of religion: it is a force of social control constraining individuals; it brings solidarity; it has cognitive functions; it produces collective effervescence, a charge external to and coercive of individuals which unites them---as when people cry on seeing their favourite musicians . \\
\midrule
Marx & The function of religion is to opiate the masses---to nullify their consciousness and to let the capitalist class exploit them forever, in the name of sending them to heaven . It prevents revolutionary consciousness from developing; religion is therefore status-quoist, securing the interest of the capitalist class . \\
\midrule
Weber & He is not asking about the function of religion . He studies the social action of religious people---Hindus, Protestants, Muslims, Buddhists---and concludes that some religions have the power to bring in another force called capitalism . That way religion is a source of social change . But it depends which religion: Protestantism can bring social change; other religions cannot, and Hinduism maintains the status quo . Religion as a system can therefore be a source of both social change and status quoism . \\
\bottomrule
\end{tabularx}
\vspace{6pt}

\begin{CriticalPoint}
The consequence drawn in class: Europe advanced and India could not, remaining underdeveloped with the same system maintained . Caste still rules; editorials are still written on caste; notes are still made on caste; Indian Society Paper 2 is still read as caste .
We know people by their last names; in America they know people by their first names .
\end{CriticalPoint}

\begin{QuickRevision}
\begin{itemize}
    \item Modernity = industrialization for Durkheim, capitalism for Marx, rationalization for Weber .
    \item Concerns: anomie and abnormal division of labour (Durkheim); exploitation, appropriation of surplus value, alienation (Marx); disenchantment and societal alienation (Weber) .
    \item Division of labour: solidarity (Durkheim), instrument of exploitation and de-skilling---Braverman (Marx), efficiency but also its restriction (Weber) .
    \item Religion: functional and solidarity-producing (Durkheim), opiate and status-quoist (Marx), a possible source of social change but only for some religions (Weber) .
\end{itemize}
\end{QuickRevision}

\subsection{Comparison III: The Solutions They Proposed}
\vspace{6pt}
\noindent
\begin{tabularx}{\textwidth}{>{\hsize=0.25\hsize}Y >{\hsize=0.75\hsize}Y}
\toprule
\rowcolor{AccentDark}
\thead{Scholar} & \thead{Solution} \\
\midrule
Durkheim & Reforms: normal division of labour and social reforms---professional associations, a moral state, a code of conduct, moral education and professional ethics . Only then will the normal division of labour come in and society evolve as it was supposed to . \\
\midrule
Marx & Revolution and communism . Nothing else . \\
\midrule
Weber & The hope lies neither with the masses nor with the bureaucrats, because he would see no difference between them---both are manifestations of formal rationality . Whenever a crisis comes, the one who will help bring social change is the charismatic authority: one charismatic man will come one day and bring in a new society . \\
\bottomrule
\end{tabularx}
\vspace{6pt}

\begin{KeyConcept}
Max Weber is also called the \textbf{bourgeois Marx} .
Both Marx and Weber disliked capitalism---Weber clearly stated in his books that he was not happy with the present state---but the reasons are very different .
Marx is unhappy because he sees the exploitation of the working class by the bourgeois class . Weber's analysis is not so reductionist: he is unhappy because in capitalism there is no space for emotion . We have become highly disenchanted, concerned only about wealth and status, no longer concerned with helping people . The rules bring impersonality and separate us from individuals---and this he calls \textbf{societal alienation} .
It is not just alienation; it is societal . You are alienated not because you are doing something that is one per cent of your potential, but because you are not doing something as a human---you are doing something to make a task machine-like efficient .
\end{KeyConcept}

\begin{ExampleBox}
The school example: the teacher is not friendly with or concerned about students . The students to the school teacher are not Prashant, Swapnil, Rupan, Shubhangi, Vijay, Prakash or Raghavan---they are roll number one, roll number two, roll number three .
The teacher has become completely alienated from who the students are . The student to the teacher is nothing but a means to get promotion: the better the result of the class, the better the chance of promotion . Bureaucracy is going on there too .

When a teacher fails a student or gives the lowest marks, she is not going to be emotional about it . She will simply present the mark sheet: see your mark sheet, E grade, this is your performance .
Then there will be a parent-teacher meeting, in which the student sits in the corner as an object of discussion between two conscious, rule-bound, formal, mature beings---one teacher and one parent---both concerned about the efficiency of this ``product'' .
You have become one individual reduced to what is in the grade sheet . This is what Weber was concerned with: the human is no longer seen as human, but as an instrument of efficiency and calculability, someone who should be predictable and successful, someone who works like a machine .
\end{ExampleBox}

\begin{CriticalPoint}
Max Weber did not like Bismarck; it was the times of Bismarck when he was writing these things .
He never had any hope from the masses---the hope that Marx had, that the masses will bring revolution and an egalitarian society will develop . Weber does not share it, although he too is critical of capitalism .
He felt that the more we advance towards socialism or communism, the more we get rationalized and the more bureaucracy gets concentrated . Russia and China are the biggest examples: in the name of socialism, it is the bureaucracy which is ruling there .
Therefore it is not the dictatorship of the proletariat but the dictatorship of the officials that will reign over civilisations in the future---and the hope lies in a creative, honest, charismatic political leader, not in a bureaucrat .
\end{CriticalPoint}

\begin{QuickRevision}
\begin{itemize}
    \item Solutions: reforms and moral institutions (Durkheim), revolution and communism (Marx), a creative and charismatic political leader (Weber) .
    \item Weber is called the bourgeois Marx: also anti-capitalist, but because capitalism leaves no space for emotion, not because of class exploitation .
    \item Societal alienation = being alienated from acting as a human, not merely from one's creative potential .
    \item Weber places no hope in masses or bureaucrats; advance towards socialism only deepens bureaucratization---Russia and China---producing the dictatorship of the officials .
\end{itemize}
\end{QuickRevision}

\subsection{Doubts Answered at the End of Class}
\paragraph{How do we differentiate paradigm from methodology?}
Paradigm means a set of beliefs . Methodology means a set of methods used by a scientist---natural or social---to carry out research . The methods Weber had to use to understand rationalization were ideal type, verstehen and adequate causality .

\paragraph{Is Weber in between Durkheim and Marx when it comes to ideology?}
Weber never gave a functional theory . He is not talking about the function of anything to the system; he was against organicism and against positivism---that was the very first lecture on Max Weber .
He is not looking at laws, so he is not a functionalist . And generally he is not a conflict scholar either; how conflict enters later will be seen in stratification .
For now, what matters: he is neither functionalist nor conflict---he is interpretative or interactionist, an interactionist because he is studying interaction, the social action of individuals .

\paragraph{Subject matter of sociology according to Marx, in one word?}
Class conflict . The history of our society has been the history of class struggle; class struggle is the subject matter of sociology .

\begin{QuickRevision}
\begin{itemize}
    \item Paradigm = set of beliefs; methodology = set of methods (for Weber: ideal type, verstehen, adequate causality) .
    \item Weber is neither functionalist nor conflict theorist---he is interpretative/interactionist, and was against organicism and positivism .
    \item For Marx, the subject matter of sociology in one word is class conflict .
\end{itemize}
\end{QuickRevision}

\sectiondivider

\section{Answer Writing, American Sociology and Talcott Parsons}

\subsection{Answer Practice: Weber's Concepts for Analysing Legitimate Domination}

\begin{PYQLink}
\textbf{Question (2015, 10 marks, 150 words):} Which concept did Weber use to analyse the forms of legitimate domination? Time allowed in class: 7 minutes .
\end{PYQLink}

Let us understand the demand of the question . The examiner asks which concept or concepts; he might have forgotten to write the ``s'' . Concept does not mean you have to write only one concept .

\paragraph{Step 1---Look at the Keywords, Not the Sentence}
We are not supposed to read the question; we are supposed to look at the keywords .
The keywords here are \textbf{legitimate domination}, which is itself a concept, and \textbf{forms}, which is very clear .
So the idea should be to start with the definition of domination .

\paragraph{Step 2---Define Domination}
\begin{KeyConcept}
When it comes to domination and authority, both are one and the same . Authority was the word used by Talcott Parsons when he translated the work of Max Weber from German into English .
The meaning of domination---and remember, it is not power; power is different---is the probability of commands being obeyed, the probability of obeisance to commands .
\end{KeyConcept}

\paragraph{Step 3---Explain How Domination Becomes Legitimate}
What gave legitimacy to the other person to exercise power over you, and what gave legitimacy for his commands to be obeyed by you?  Domination gets legitimacy in three ways :
\begin{enumerate}
    \item \textbf{Traditional domination}---you obey someone's command because he is your father or she is your mother; legitimacy based on tradition, customs, norms and values, or on ancient texts . In the Ramayana the ideal is displayed and everything moves around it---from there we learned that the lady should be like Sita, the father like Dashrath, the man like Ram . So: legitimate domination based on tradition or ancient texts .
    \item \textbf{Charismatic authority}---sometimes the other person is very charismatic and has induced hope in a hopeless situation . You get fed up studying so many books---ancient history, modern history, medieval history, world history---nothing is registering; then you watch a motivational lecture by someone on YouTube and hope returns . You start following that teacher and subscribing to his channel . Why are you obeying him? Because you saw charisma in him---heroic qualities, extraordinary features, oratory skills . Legitimacy given on the basis of extraordinary qualities .
    \item \textbf{Legal-rational authority}---you are stopped on the road by a traffic policeman . Had some other person asked you to stop, you would not have stopped; but a uniformed man asked . You did not stop because you found charisma in him, nor because our ancient texts say we should obey policemen . You obeyed only because it is the command of rules and regulations---constitutionalism; as a citizen of India we have to obey the administration .
\end{enumerate}

\paragraph{Step 4---Answer the Actual Question: Which Concepts?}
\begin{enumerate}
    \item \textbf{Social action}---whether you obey the policeman, listen to the motivational guru, or obey your father, everything you are doing is your social action .
    \item \textbf{Ideal type}---the forms of authority are not pure types . A traffic policeman may also be a charismatic personality; in the real world you cannot cleanly differentiate . There are many IPS and IAS officers, but some are very charismatic . The larger trend is that you obey people based on rules---but traditional, charismatic and legal-rational authority are ideal types, not pure types .
    \item \textbf{Domination / authority} itself .
    \item \textbf{Adequate causality}---why did I stop my vehicle and pay two thousand rupees? A normal person draws a monocausal relation: he stopped you, therefore you stopped, and you had to pay . But in reality multiple reasons might exist---for instance, that person was driving with his father sitting beside him, and therefore he had to stop . It is not only that the policeman waved his hand . In reality you do not have monocausality; it is multi-causality, and at most, as sociologists, we can draw an adequate causal relation: if he stopped, there is a probability that he stopped largely because of this .
\end{enumerate}

\begin{PYQLink}
These are the minimum four concepts that you should have talked about .
A student asked whether we can also write that Weber rejected something---no . When we see the question, the art of answer writing is basically to write what is asked . We cannot talk about the history of the concepts, the evolution or the origin of sociology here .
Criticism is not required in this question---it is a 10-mark question . But you can mention that interpretative sociology was a vision of Max Weber, and that this vision helped sociology depart from the macro perspective to the micro .
How will you be different if you write only what is required? That IS how you are different . Most people do not even address the question; without addressing it they ask how to be different . First address the question .
What matters is structure: start with the definition of domination; then legitimacy and the forms of legitimate domination; then the examples; and finally the concepts Weber used . All of you study similar things from similar notes and lectures, and still your marks will differ .
\end{PYQLink}

\begin{QuickRevision}
\begin{itemize}
    \item Domination = probability of commands being obeyed; not power . Authority is Parsons' translation of Weber's German term .
    \item Three bases of legitimacy: tradition/ancient texts, extraordinary charismatic qualities, and formal rules (constitutionalism) .
    \item Four concepts to name in the answer: social action, ideal type, domination/authority, adequate causality .
    \item For 10 marks: no history, no origins, no criticism---define, classify, illustrate, then name the concepts . Structure is what separates scripts .
\end{itemize}
\end{QuickRevision}

\subsection{Reading the Newspaper Sociologically I: The Cult of the Topper}
Now read a newspaper without any pressure and figure out what the author is trying to say and, as a student of sociology, how you are going to look at it . Notice the difference between how a common-sensical person reads an article and how a student of sociology reads it . Both pieces discussed here are from the \textit{Indian Express}, and the author, revealed at the end, is Abhijit Pathak, a sociologist from JNU .

\paragraph{The Argument of the Article, in Summary}
The writer begins from the observation that our society worships success and adores toppers, and argues that precisely for this reason it needs to recognise there is no necessary link between a person's intellectual ability or academic excellence and their performance in a board examination . One can score full marks in English without any curiosity for poetry, short stories or novels; the topper in physics need not have any scientific temper or urge to understand the laws of nature .
The experience of wonder, the joy of learning and unlearning, and the spirit of inquiry count for nothing, because the prevailing system believes in the fetish of quantification .

The craze for marks converts the learner into a smart consumer---the only skill required is the ability to consume knowledge capsules and facts, through mock tests, endless drilling and an obsession with ranking . Schools sell toppers as perfect products to raise their own market value, and the young learner is reduced to repeating what has been fed in .
In such a world there is no poetic wonder, no criticality of social science, no inquiring spirit of science---one is only an exam warrior .

As fact-centric objective short questions become the new normal and memorisation becomes the key to success, it becomes very difficult not to score at least 80 per cent, while those who score 95 per cent still experience acute distress . The writer names this the \textbf{pathology of success producing the stigma of failure} .

Students also find almost no guidance while choosing subjects and entering college, for two reasons :
\begin{enumerate}
    \item \textbf{Academic disciplines are hierarchized through the pragmatism of market rationality:} Just as in class ten you ``chose'' non-medical, so students, pressured by anxiety-ridden parents and driven by peer culture, prefer prestigious subjects---physics, economics, commerce, English literature, psychology---even when they have no intrinsic inclination toward them .
    \item \textbf{Colleges and universities are themselves hierarchized in an age of ranking:} Institutions like St Stephen's or Lady Shri Ram College are consumed as fancy brands . Students prefer the prestigious college over the subject they like: Sanskrit or education does not matter, Miranda House or Hindu College matters at any cost .
\end{enumerate}

This mass psychology of consumption complicates the politics of the cut-off and destroys human possibilities: someone who could have contributed immensely to the study of Hindi literature ends up doing economics; someone carrying the burden of a famous college's economics course might have found real satisfaction in history at a less reputed institution . The craving for status triumphs over personal liking, and education suffers .

The system, the writer says, works like a soulless machine, more interested in eliminating people than in finding or arousing anyone's real aptitude---if you do not have 99 per cent, you cannot study sociology at Lady Shri Ram College .

The final paragraph asks us to imagine the mental agony these youngsters pass through as social Darwinism is normalised and hyper-competitiveness becomes the philosophy of the age . Schools become highly oppressive institutions: heavily loaded school bags, weekly tests, never-ending home assignments, constant performance anxiety, parental pressure and the urge to be a topper . The world they inhabit, in the writer's words, has no rainbow, no sunset, no music . Even though thinkers like Rabindranath Tagore and Jiddu Krishnamurti felt the pain of this neurosis, the system remains as it is and the society being created is bound to be toxic and violent .

\begin{KeyConcept}
What is social Darwinism here? We take an entrance examination; whoever is intelligent gets selected and the rest are left behind---survival of the fittest, normalised inside universities .
The sociological reading: this entire article is a description of the \textbf{irrationality of rationality} . Formal rationality---the fetish of quantification, marks, cut-offs, rankings, brand value---produces outcomes that defeat the very purpose of education .
\end{KeyConcept}

\begin{QuickRevision}
\begin{itemize}
    \item The article: no necessary correlation between academic excellence and board marks; the system rests on a fetish of quantification .
    \item The learner becomes a consumer of knowledge capsules; schools sell toppers as products; disciplines and colleges are hierarchized by market rationality and brand consciousness .
    \item Social Darwinism is normalised and hyper-competitiveness becomes the philosophy of the age; Tagore and Krishnamurti had already felt this neurosis .
    \item Sociologically, the whole piece illustrates formal rationality producing the irrationality of rationality .
\end{itemize}
\end{QuickRevision}

\subsection{Reading the Newspaper Sociologically II: The Age of Instant Enlightenment}
The second article, by the same author, is titled ``The age of instant enlightenment'' . Its argument, summarised:
We live in an age of instantaneity and quick consumption . We had fast food; today we have fast education and fast motivation .

Everything---a McDonald's burger, a cup of Starbucks coffee, a message on Twitter, even political satire on YouTube---is produced and consumed as fast as one can imagine . There seems to be no end to the rapidly changing flow of commodities, symbolic goods, and Facebook and Instagram images and messages . We are only running, eternally hungry for something exciting and new .

This age of instantaneous consumption transforms even religiosity and spirituality into a product: from a quick look at a book by Osho while waiting at the airport, to a three-day retreat in a fancy ashram to learn a breathing exercise, the aspiring and mobile class turns even the alternative into a consumable product .

Not surprisingly, YouTube is full of innumerable packages of instant enlightenment---Zen Buddhism as the easiest spiritual path on earth, two minutes to heal your upset mind, one daily practice that will bring divine grace . We click on them precisely because the path is advertised as easy .

There are moments when this class too feels a need for peace, tranquillity and self-realization; but it is not easy to be free from attachment to the comforts and privileges it is used to, nor from the mood of instantaneity . Herein lies the sociology of the arrival of all sorts of life coaches, motivational speakers and new-age spiritual messiahs .

They come forward with packages of instant redemption---a feel-good movement for a couple of days before one goes back to the old routine of working like a robotic performer, earning, saving, consuming and buying all sorts of insurance . Come back from the office, have a cup of coffee, open the laptop, spend five minutes with a talk by a well-known spiritual speaker on two steps to think right, or a fifteen-minute morning meditation---and everything is calculated .

\begin{KeyConcept}
Through these two very different pieces the author is describing the \textbf{iron cage of rationality} that we cannot escape .
A caution stated in class: this does not mean the author is a follower of Weber . In Indian sociology no person is a follower of any one paradigm .
\end{KeyConcept}

\begin{QuickRevision}
\begin{itemize}
    \item Instantaneity and quick consumption turn even spirituality into a consumable product---airport Osho, weekend ashram, YouTube enlightenment .
    \item The mobile, privileged class seeks peace but cannot leave either its comforts or its mood of instantaneity .
    \item Hence the sociology of life coaches, motivational speakers and new-age messiahs offering instant redemption in calculated fifteen-minute packages .
    \item Both articles describe the same thing: the inescapable iron cage of rationality .
\end{itemize}
\end{QuickRevision}

\subsection{The Origins of American Sociology}
We now leave European sociology . In American sociology our syllabus has three scholars---Talcott Parsons, Robert Merton and George Mead---and each of the three belongs to a very different tradition, just as Marx, Weber and Durkheim did .

\begin{ComparisonBox}
\begin{itemize}
    \item European sociology had its foundation in historical events: the Renaissance, Reformation, Enlightenment, French Revolution and Industrial Revolution . These produced changes which were analysed, described and critiqued by Marx, Weber and Durkheim .
    \item American sociology has no such historical event or story of its own . The history behind American sociology is European sociology . American sociologists borrowed a great deal from Europe---functionalism, Marxism, feminism, everything was borrowed . But they also did some things European sociology did not do at that point of time .
\end{itemize}
\end{ComparisonBox}

\paragraph{The Reasons for Its Origin}
\begin{enumerate}
    \item \textbf{European theoretical tradition}---this reason is external to America .
    \item \textbf{Immigration}---America is basically a land of immigrants; people made America, through slavery and through groups such as the Quakers . Even today, most of the people working in Facebook and Microsoft are Asians, not Americans; the industries are theirs but many of the workers are from non-American states .
    \item \textbf{Urbanization}---the issues that arose with settling those migrants .
\end{enumerate}

\begin{ComparisonBox}
Recall Durkheim: migration from rural areas to urban pockets led to a rise in material density and population volume . The increased population could not be adjusted in limited land, so a specialised division of labour was allotted---and instead of competition or conflict among people, you now have cooperation . That is how specialisation came into existence .
\end{ComparisonBox}

In America, immigration came from multiple pockets---from Europe, and from Africa through slavery; after the abolitionist movement most of these people migrated within America .
With immigration came the problems of urbanization: how to settle these migrants . And with migration, anomie also built up in American society---crime rates increased, along with poverty and homelessness . These are precisely the things Durkheim calls anomie, and Durkheim knew they were the result of migration, which he studied while looking into the impact of industrialization .

\paragraph{The Golden Triangle}
\begin{DataFact}
Urbanization, immigration and industrialization were happening at a very rapid pace in the \textbf{golden triangle}: Chicago, Washington DC and Boston .
Rapid change means anomie---sudden change .
\end{DataFact}

With rapid change, society was moving from simplicity to complexity, from pathology to diversity, from absolutism to relativism . The scholars who discuss simple-to-complex are Durkheim and, in fact first, T{\"o}nnies .
Diversity here means many people coming from different backgrounds---black slaves, Europeans, Quakers .

\paragraph{The Immigration Timeline Shown in Class}
\vspace{6pt}
\noindent
\begin{tabularx}{\textwidth}{>{\hsize=0.35\hsize}Y >{\hsize=0.65\hsize}Y}
\toprule
\rowcolor{AccentDark}
\thead{Period} & \thead{Who Came} \\
\midrule
1600s & Puritans; Jamestown; Quakers; slavery . \\
\midrule
1700 onwards, around 1750 & Scots-Irish . \\
\midrule
1800 & Irish, following the potato famine---Irish Catholics . \\
\midrule
1850 & Germans, Danish, Icelandic, Norwegian, Swedish . \\
\midrule
Later nineteenth century & Many people from Eastern Europe . \\
\midrule
1890s & The peak of immigration into America . \\
\midrule
1900s & When anti-Jewish sentiment was at its peak in Europe, most Jews went to America---the figure given in class was 80 per cent . Members of the Frankfurt School also fled there during the war . \\
\bottomrule
\end{tabularx}
\vspace{6pt}

\begin{CriticalPoint}
A consequence noted in class: you would not see Israel defying America; there is a bonhomie, a relation of harmony, between America and Israel which has sustained for so long, because America is influenced by a Jewish lobby .
Although these people fled to America, they were such intellectuals that they contributed immensely in scientific advancement, defence equipment and agriculture .
\end{CriticalPoint}

\begin{QuickRevision}
\begin{itemize}
    \item American sociology arises from three sources: the European theoretical tradition, immigration, and urbanization .
    \item Its ``history'' is European sociology itself---theories were borrowed wholesale---unlike Europe, whose sociology grew out of Renaissance, Reformation, Enlightenment and the two revolutions .
    \item Immigration produced urbanization and anomie---crime, poverty, homelessness---concentrated in the golden triangle of Chicago, Washington DC and Boston .
    \item The immigration timeline runs from Puritans and Quakers in the 1600s to the 1890s peak and the Jewish and Frankfurt School exodus of the 1900s .
\end{itemize}
\end{QuickRevision}

\subsection{The Chicago School and Its Founders}

\begin{DataFact}
The first university department of sociology was at the University of Chicago, founded in 1892 on the appointment of Albion Small as Head Professor of Sociology .
Himself a researcher on the histories of institutions and schools of thought, he played a central role in creating sociology as an academic discipline .
He founded the \textit{American Journal of Sociology} and edited it for thirty years . Emile Durkheim was on the editorial board of that journal---described in class as its chief editor . \textit{[as stated in lecture]}
At the same time Albion Small helped found and manage the American Sociological Association, the source of the maximum literature in sociology .
Small was also deeply involved with social reform---a tradition that has continued throughout the department's history .
\end{DataFact}

\paragraph{The Founders and Their Debate}
The founders of American sociology---of the Chicago School, the first school of sociology in America---were Albion Small, Lester Ward and William Sumner . All three had read European sociology . The school was established in the 1890s .

\begin{CriticalPoint}
Ward and Sumner were engaged in a debate about what the approach should be: how should we solve the issues of poverty and homelessness produced by immigration, urbanization and industrialization? 
One position was: do not do anything---rather let us study it theoretically and understand why these things happened, uncover the laws behind the changes and then discover when order will come . This is speaking like Auguste Comte, and this position was that of social Darwinism and evolutionism . In the lecture this position was attributed to Ward, who was described as completely against intervention . \textit{[as stated in lecture]}
The other position was: no, we have to bring reforms; we have to intervene . In the lecture this was attributed to Sumner . \textit{[as stated in lecture]}
From here come two different schools in American sociology: one interested in reforms, the second against reforms .
\end{CriticalPoint}

The analogy given: the Congress in India also had two groups, moderates and extremists, and initially the moderates dominated . Similarly, in the founding of American sociology, the group that dominated was the social reformists . Those who said ``we will only theorize'' did not dominate---Europe had already done that; something different was needed .
So social reformism and the Chicago School are one and the same, because the Chicago School proposed and was dominated by social reformism .

\paragraph{The Four Challenges}
\begin{enumerate}
    \item Poverty 
    \item Homelessness 
    \item Crime rate 
    \item Racism 
\end{enumerate}

Most of the migrants by the end of the 1900s were black, coming from Africa; today we label them African-Americans . In contemporary politics they are regarded as African-Americans, not simply as Americans .
In India we do not say Sikh Indians or Muslim Indians; we do not label people as Kashmiri Indians or North-East Indians as a separate category---we simply say Kashmiri, North-Eastern, South Indian, Delhi or Mumbai .
But in America the labels are African-American, Hispanic-American (again, Latinos), Asian-American . Most Indian-Americans live in the suburbs of America---that is, at the periphery, not at the core . And there are ghettos as well .

\begin{QuickRevision}
\begin{itemize}
    \item The department at the University of Chicago was founded in 1892 under Albion Small, who founded and edited the \textit{American Journal of Sociology} for thirty years and helped found the American Sociological Association .
    \item Founders: Albion Small, Lester Ward, William Sumner; their debate was between non-intervention/social Darwinism and social reform/intervention .
    \item The reformists dominated, so Chicago School = social reformism .
    \item The four challenges before American society: poverty, homelessness, crime rate and racism---with hyphenated labels and suburban/ghetto settlement patterns as the racial backdrop .
\end{itemize}
\end{QuickRevision}

\subsection{The Reformers: Jane Addams and Du Bois}

\paragraph{Jane Addams---Poverty and Homelessness}
Jane Addams is described as the first ever sociologist to get the Nobel Prize .
She established the \textbf{Hull House}---like the charity organisations working to ameliorate the conditions of poor people . Through it she helped marginalized sections by giving them education, talking about their rights, giving them free food and a healthy lifestyle: the minimum desires that should be met . She was a reformist---acting in the world, not theorizing about the world .

From here on came social activism, which is not sociology and is a very different thing . You may be a social activist without being a sociologist, and you may be a sociologist without being a social activist . And there is a third category, social work---anybody can be a social worker . These three categories are very different from one another .

\paragraph{W. E. B. Du Bois---Racism}
Du Bois was himself an African-American, and he thought about the reasons behind racism in America . He gave the concept of the \textbf{colour line}---whites and blacks .
This colour line is a creation of Europe . It was in Europe that the age of discovery began; while exploring different continents they came across people who did not look like them and called them blacks or browns . The labels we use today---brown, black---are actually from the European perspective; we are using European labels .

Du Bois says the white-black colour line is not just a colour line: it is oppressive, and it is imbued with the historical exploitation of blacks by whites . The races may be different, of course, but there is also inequality in this difference .
Even when an African becomes an American by citizenship, he is still not purely an independent man with reasoning who can exercise fundamental rights just like any other white American . He is still an African-American---which means that man has dual consciousness .

\begin{KeyConcept}
\textbf{Dual consciousness:} I am African as well and American as well . In the private sphere I am African at home; the moment I go outside, I have to act like an American .
Imagine sitting with white Americans as a student at Harvard or Cambridge, thinking of equal opportunity, expecting to graduate just like them---and then after class, in the canteen, finding that whites have different tables and blacks have different tables . If a black student tries to eat with the whites, the whites get up and move elsewhere .
So Du Bois could see a clear-cut inequality in democracy . America is considered the oldest democracy---and yes, it is the oldest democracy---but to Du Bois it was not really a democracy .
\end{KeyConcept}

\paragraph{Chicago as a Laboratory}
Chicago at that point of time was not just a city; Chicago was a laboratory, because so many issues were cropping up---homelessness, poverty, crime rate, racism .
To these reformers it was a laboratory where they could freely work: go out into the world instead of theorizing from an armchair, have discussion and interaction with people, and get to know their everyday experience .

And here Max Weber returns . When these reformers went out into the public sphere and tried to interact with the victims of immigration---Africans, the poor, the homeless, and women as well---they got to know their concerns . This is why we call it an application of verstehen at the end of American sociology .
They started with reformism; they had challenges before them; and gradually, in order to understand those challenges better, they moved towards verstehen---discussion with women, with criminals, with blacks, with homeless people, chatting with them . They were not like Durkheim, Marx and Weber .

\begin{QuickRevision}
\begin{itemize}
    \item Jane Addams: Hull House, education, rights, food and a healthy lifestyle for the marginalized; the first sociologist to receive the Nobel Prize; a reformist acting in the world .
    \item Social activism, social work and sociology are three distinct categories .
    \item Du Bois: the colour line as a European creation, oppressive and historically loaded; dual consciousness; inequality inside the oldest democracy .
    \item Chicago became a laboratory, and everyday interaction with the oppressed amounted to an application of verstehen .
\end{itemize}
\end{QuickRevision}

\subsection{How American Sociology Differs from European Sociology}
\begin{enumerate}
    \item American sociology, unlike European sociology, was more pragmatic (pragmatism) and less theoretical .
    \item Unlike European sociology, which had historical antecedents in its emergence, American sociology developed as a reaction to issues related to immigration and the exploitation of the immigrants---homelessness, racism, poverty, discrimination and the resultant consequence, that was increasing crime rate .
    \item American sociology developed as an elixir to heal the wounds in the soil of Chicago (the golden triangle) .
    \item Unlike European sociology, which was largely structural and macro, American sociology was mainly micro and engaged in everyday interaction with the oppressed .
\end{enumerate}

For example, Jane Addams going to the homes of poor people, sitting with them, understanding their concerns; going to the homes of homeless or single women whose husbands had left them with their children .

The contemporary parallel drawn: the flood of immigrants from the Middle East to Germany and France . Most European states closed their doors, and those who were absorbed were taken in as people required to perform menial and dirty tasks . Imagine American scholars dealing with such people through everyday interaction, trying to understand their concerns deeply rather than theorizing about them---verstehen .

\begin{QuickRevision}
\begin{itemize}
    \item Four dictated distinctions: pragmatic rather than theoretical; a reaction to immigration and its exploitation rather than to historical antecedents; an elixir for the wounds of Chicago; micro and interactional rather than structural and macro .
    \item The method that emerges out of reformist fieldwork is verstehen .
\end{itemize}
\end{QuickRevision}

\subsection{From Chicago to Harvard: Sorokin and Parsons}
The Chicago School developed in 1890 and dominated American sociology until 1930; Chicago School was synonymous with American sociology in this period .
In 1930 it lost its relevance, because Pitirim Sorokin, a Russian-American, came to America and established a department at Harvard University---the Harvard School of sociology .

Sorokin was himself a victim of the Russian revolution, whose timeline runs from 1917 through the following decades; it was a series of revolutions . From 1922 to 1930 he was neither Russian nor American---he was stateless, because the state was still to be formed after the revolution .
He was a supporter of socialism, a crude supporter of the Marxian idea of socialism, and had joined the Socialist Revolutionary Party in Russia . In 1930 he fled to America and became an American citizen .

In America he was writing on why the revolution failed, and why the idea of Marxism did not turn out as it was supposed to---why it turned out brutal and violent . Marx never said revolution would entail horrors or murders; to Marx revolution simply entails a change in the mode of production that is inevitable . If it is inevitable, why are so many murders taking place? Is murder also part of the inevitability principle? 

\begin{KeyConcept}
Sorokin is also known for the \textbf{cyclical theory} (to be taken up in later chapters) . Instead of moving linearly---ancient to feudal to capitalist to socialist to communist---he said that in reality it is a cycle: where you start, you come back there again .
Even Marx's movement from primitive communism to pure communism can be read as a cycle, though it seems linear---as always, interpretation varies with the academician .
For Marx the transition entails inherent contradictions within the mode of production and the domination of one group over another . According to Sorokin there are no two groups opposed to each other in that way .
\end{KeyConcept}

\paragraph{Parsons versus Sorokin at Harvard}
One day Talcott Parsons came to study at Harvard and became a student of Sorokin---and their views turned out to be very opposite . A student not agreeing with the teacher: from there on two different schools within Harvard got established, with some agreeing with Parsons and some with Sorokin, just as in Germany there were supporters and non-supporters of Hegel .
Then came Robert Merton, who as a student had the choice to do his PhD under either Sorokin or Talcott Parsons, Parsons having by then gained professorship .

From 1930 to 1960---thirty years---Parsons dominated over Sorokin, and structural functionalism dominated .
Why? Parsons was a great admirer of Durkheim and Max Weber; Sorokin was a great admirer of Marxism . America was a democracy---indeed a capitalist democracy---and a democracy will never support violent revolutions . Sorokin, following his cyclical theory, believed capitalist democracy was about to decline . So at the end of the day Harvard University promoted Talcott Parsons and demoted Sorokin, because Parsons spoke like an American democratic person who would never talk about radical or revolutionary ideas .
That paradigm dominated for about twenty years and declined in 1960 .

\paragraph{How Parsons Met Weber}
When Max Weber died of pneumonia in 1920, his wife Marianne Weber conducted conferences, speeches and discussions on Max Weber's thought, over which she used to preside .
Talcott Parsons, visiting Germany, entered such a discussion hall in 1925 and listened to Marianne Weber . On the thoughts of Max Weber he got enlightened; that was his first encounter with Weber's thought .

\begin{KeyConcept}
\textbf{Talcott Parsons = Weber + Durkheim.}
Weber gave the theory of social action, which is micro-sociology . Durkheim gave the theory of social fact, which is macro-sociology . Parsons had before him two people with two opposite ideas---and he supported both .
Instead of choosing micro or macro, he decided to come up with an integrated approach, integrating both .
And this is not an easy task: he is integrating not just micro with macro, but positivism---a theory about laws of the social world---with a web of meanings in which there is no theory . One says there is objective reality, the second says there is no objective reality . Weber's social action theory was itself indispensable for the construction of phenomenology and ethnomethodology, which are micro tools .
\end{KeyConcept}

\begin{QuickRevision}
\begin{itemize}
    \item Chicago School dominated 1890--1930; Sorokin's arrival and the founding of the Harvard department in 1930 ended that dominance .
    \item Sorokin: stateless victim of the Russian revolution, socialist, author of the cyclical theory against linear evolution .
    \item Parsons, his student, disagreed; Harvard, in a capitalist democracy, promoted Parsons and structural functionalism dominated 1930--1960 .
    \item Parsons met Weber's thought through Marianne Weber's discussions in 1925; Parsons = Weber + Durkheim, and his project is the integration of micro and macro, of meanings and of positivism .
\end{itemize}
\end{QuickRevision}

\subsection{Talcott Parsons: The Structure of Social Action}
In the pursuit of integration, the first theory Talcott Parsons gave is the \textit{Structure of Social Action} . The integration is visible in the name itself: positivism is about structures, and social action is not about structure but about the meanings of individuals---and he says structure of social action .

\paragraph{The Unit Act}
\begin{KeyConcept}
Just as we have units in chemistry---atoms, molecules---and units in physics, there must be units in the social world too .
The unit of the body is the cell or tissue; tissues build up into organs, organs into the human body . Similarly there must be a unit of society: the unit of society is action---the actor and the action .
He called it the \textbf{unit act} . Many unit acts form society, and then society dominates the unit acts .
Notice how much like Weber this is: social action comes, then structure, then structure dominates the social action of individuals .
\end{KeyConcept}

\paragraph{The Marriage Example}
There is a man who wants to date a girl, because he wants to get married . The goal is marriage---and goals are always momentary, never permanent . Today your goal is to become an IAS; tomorrow it will be to become chief secretary; then a politician; then a motivational guru . Goals keep changing .

What would Weber say? If I am looking at photographs of girls my parents have shown me, my looking at those photographs is a social action . It could be \textit{zweckrational} or affectual---you would not be able to figure out which, and there is no monocausal explanation . But a researcher can empathetically understand the meanings behind it .
And yet, be honest: if I am going through pictures of girls on my mobile phone, is there any certainty about what someone in my shoes would conclude? He might say I am on a dating app; it could equally be timepass . Both are very different things .

So Parsons says: it is not an easy task . We should understand social action---but how? And while understanding social action, you first need to understand why that person has such goals and why he adopts particular means .

\paragraph{Five Choices, and How Four Were Eliminated}
\vspace{6pt}
\noindent
\begin{tabularx}{\textwidth}{>{\hsize=0.3\hsize}Y >{\hsize=0.7\hsize}Y}
\toprule
\rowcolor{AccentDark}
\thead{Choice} & \thead{What Happened, and Why} \\
\midrule
Parents' choice & Went with it and married---and after three months discovered the man was a cheat with other relationships, and the marriage ended in divorce . Next time the parents offer again, but the actor says he will look for himself . \\
\midrule
Matrimonial.com & Rejected . The site presents people like a menu of food; and in any case the profiles there are made by parents for their sons and daughters---the very source of the earlier deception . ``This is not me'' . \\
\midrule
School or college & Rejected . That girl had already rejected you once; there is no courage to go back, and in any case she is now in Bangalore . \\
\midrule
Dating app & Rejected on two grounds: normative---``these are not my values, dating apps are for losers, all the people there are fake, I believe in meeting people face to face''; and situational---the app charges ten thousand rupees to maximise your chances, and you checked your pocket . \\
\midrule
Friend's suggestion & Finally chosen . Your friend suggests someone according to your taste, preference, family background and status---according to your caste, class, religion, family background, and your requirement that the person be educated . \\
\bottomrule
\end{tabularx}
\vspace{6pt}

\begin{QuickRevision}
\begin{itemize}
    \item \textit{The Structure of Social Action} integrates structure (positivism) with social action (meaning) in its very title .
    \item The unit of society is the unit act---actor plus action; many unit acts make society, and society then dominates them .
    \item Goals are momentary, not permanent; Weber can interpret a social action but cannot fix its meaning with certainty .
    \item Parsons therefore asks how goals are formed and how means are selected---illustrated by the elimination of four of five marriage options .
\end{itemize}
\end{QuickRevision}

\subsection{Normative and Situational Constraints}
Before arriving at the decision and cancelling all the other choices, many things were going on in the actor's head . Those things, in the language of Talcott Parsons, are \textbf{normative constraints} acting from above on the actor's orientation and \textbf{situational constraints}, acting from below .

\begin{KeyConcept}
Situational constraints have three criteria: \textbf{income} or money; \textbf{geography}, that is your location; and the \textbf{capability} of the individual .
Normative constraints are the \textbf{norms and values} which decide what you consider available to you at all .
The dating app was refused both because the norms tell you that you are not the kind of person who is available on dating apps, and because you could not pay the fee .
\end{KeyConcept}

\begin{KeyConcept}
Write down: Max Weber, in his theory of social action, discussed how an actor acts in a situation and what meanings he attributes to the situation .
However, Weber did not take into account the constraints faced by the actor in a situation---constraints which not only influence the availability of the means before the actor, but also influence the goals the actor decides for himself .
\end{KeyConcept}

\paragraph{Even the Goal Is Not Yours}
Why are you thinking of getting married? Because marriage is an institution and a norm---so even in this so-called individual decision, norms are entering .
Your goal to become an IAS is not culturally undefined . If you were American or European, would you have given a thought to becoming an IAS? Perhaps something else entirely . The goals you set are influenced by the norms and culture of society .

\paragraph{The IAS Example: Five Means and the Constraints on Them}
\vspace{6pt}
\noindent
\begin{tabularx}{\textwidth}{>{\hsize=0.35\hsize}Y >{\hsize=0.65\hsize}Y}
\toprule
\rowcolor{AccentDark}
\thead{Means Available} & \thead{How the Constraints Operate} \\
\midrule
Coaching & A student from Nashik, Maharashtra cannot come to Delhi because Delhi is expensive---an income constraint; there is doubt about capability; and geography is against it . For a girl child there is a further normative constraint: parents often do not send a girl child to Delhi even when they would send a boy, because spending on the boy feels like investment and spending on the girl feels like expenditure---and because of anxieties about safety in a ``crime city'' . Covid cases and the norm of social distancing added a further situational constraint . \\
\midrule
Self-study standard books and YouTube videos & You are not free to decide on your own which book is best for you; you refer to many toppers before buying a book, and you also refer to teachers, friends and seniors . Some books you cannot purchase---the thick, expensive Haralambos, the ``bible'', is set aside for the cheap condensed version with small pages and poor paper quality . So even self-study is chosen under situational and normative constraints . \\
\midrule
Library & You see friends coming back from the library every day, and anxiety builds---everyone studies there for ten hours, so what will become of me? You google the best library; you meet the owner, who offers breakfast, lunch, dinner, internet, books for every optional, water, coffee, air-conditioning---for ten thousand rupees a month, over and above your rent, when your parents have said twenty thousand is all they will pay monthly . So your choice of library is also influenced by situational constraints, and you end up in a library that has none of that but does have an atmosphere for studying . \\
\midrule
Toppers' notes & Ultimately used inside the library in place of standard books . \\
\midrule
Corruption and contacts & It does not work for UPSC, though the claim in class was that it works in BPSC and UPPCS . It is rejected because it is against your values and against what your family would allow . \\
\bottomrule
\end{tabularx}
\vspace{6pt}

\begin{CriticalPoint}
Were these three decisions---library, toppers' notes, no corruption---independent of external factors? No . Budget, family, gender, even religion and region shaped them .
On region: although UPSC is open to everyone, largely it is North Indians and South Indians who come to Delhi to prepare . You do not see many people from the North East, or from Andaman, or Jammu and Kashmir, or parts of Gujarat and Punjab preparing for civil services . No statistics were claimed for this---it was offered as a general observation .
There are many pockets where people do not share the ambition to become IAS; they have their own ambitions to go to Australia, Canada or the Middle East---from Gujarat and Kerala to the Middle East, for instance .
So even your goal of becoming an IAS is highly situational: where you were born, what kind of family you grew up in, what kind of education you got .
\end{CriticalPoint}

\begin{QuickRevision}
\begin{itemize}
    \item Two constraints: normative (norms and values, from above) and situational (income, geography, capability, from below) .
    \item Dictated critique: Weber discussed how an actor acts and what meanings he attributes, but did not take into account the constraints that shape both the available means and the goals themselves .
    \item Even the goal---marriage, or becoming an IAS---is culturally supplied, not individually invented .
    \item Illustrated across coaching, self-study, library, toppers' notes and corruption, with gender, budget, region and family operating throughout .
\end{itemize}
\end{QuickRevision}

\subsection{The Three Influences on Parsons}
Write down: Talcott Parsons was highly influenced by three schools of thought---utilitarianism, idealism and positivism .

\paragraph{Utilitarianism}
Utilitarian scholars believe that individuals are rational actors: they have multiple choices before them and will make their choice rationally .
Parsons agrees that actors make rational choices---but says rational choices are not made in ignorance of the external constraints . Those constraints actually go into the making of that so-called final rational choice .

\paragraph{The Dinner Example}
The goal is to have dinner---in fact, to have pizza . The choices: go to a restaurant; order on Swiggy or Zomato; go to a dhaba; order from Domino's; or make it at home .
There is very little money left in the pocket, and the auto or Uber fare to the place that is far from home adds to it---situational constraints .
Why did you not go to the dhaba? Because it feels ``low class'' to you---that is a norm entering, along with your class position, your status and your family upbringing . Even if a lavishly named dhaba charged a thousand rupees you would go there, but not to the ordinary one .
So normative and situational constraints finally made you eat at home .

\begin{ComparisonBox}
What would Weber see? A person eating a sandwich at home . Would Weber have the means to know that this is a pizza to me---a pizza made by my own hard work? To me it is pizza; to someone else it is simply bread and sauce . Reality is subjective .
What does Parsons ask? Tell me first why you are eating pizza, why in this manner, why at home---and also tell me the constraints which made you perform this action . That is the larger idea of Parsons over Weber .
\end{ComparisonBox}

\paragraph{The Gutter Cleaner and the Limits of Historical Context}
If I am working as a gutter cleaner, there is a historical context which has to be understood to know why I am performing this menial task . A researcher would find that lower castes have been engaged in such menial tasks for historical reasons, and that this has continued among some communities till date .
So my action of cleaning a gutter is not just my own individual action---there have been historical constraints, situational constraints, an inability to afford education, and many other things .

But when you talk about historical context, that is a very large picture, and it may not capture the individual situational essence . This is what Parsons is saying to Max Weber .

\begin{ComparisonBox}
The analogy used: utilitarians study very short, minute-by-minute actions---like YouTube shorts .
Max Weber is not studying shorts; Max Weber is studying the full movie, asking why this actor plays only this much of a role in the whole film .
But Talcott Parsons says nothing can be traced simply by looking at history: situational and normative constraints have also to be looked into to understand the essence of reality .
\end{ComparisonBox}

\begin{QuickRevision}
\begin{itemize}
    \item Parsons draws on utilitarianism, idealism and positivism .
    \item The ``rational'' choice is itself produced by external constraints---the dinner and dhaba example .
    \item Weber gives the meaning of the act; Parsons demands the meaning plus the constraints behind it .
    \item History gives the large picture (gutter cleaner, caste) but misses the situational essence .
\end{itemize}
\end{QuickRevision}

\sectiondivider

\section{Objectivity in Weber, and Parsons on Pattern Variables and the Social System}

\subsection{Answer Practice: Weber's Method of Maintaining Objectivity}

\begin{PYQLink}
\textbf{Question (2016):} Examine Max Weber's method of maintaining objectivity in social research . Reference given in class: Vishal, 14th edition, page 13, question 4 .
The question is of 20 marks, but let us assume it is a 10-mark, 150-word question . Writing 10 marks is more challenging than writing 20 .
\end{PYQLink}

The process is the same: first get uncomfortable by looking at the question, then get comfortable by figuring out something that gives you comfort .

\paragraph{Step 1---Identify the Keyword}
``Examine'' is not a very important keyword . As always, we have to write both sides . It does not matter whether the direction is ``elucidate''---even then the other side must be written .
The single keyword here is \textbf{objectivity} .

\paragraph{Step 2---Where Does Objectivity Come From?}
The notion of objectivity comes from the natural sciences . Positivists are the people who look for objectivity in the social world; they believe the social world is as objective as the natural world, and therefore that the scientific method can be applied in the study of the social world .

Max Weber believed the social world is not objective, at least not on the lines of the natural sciences . But although he believed there is no objectivity in the social world, he supported the idea of considering the meanings and motivations of individuals objectively .
In the social world there is no objectivity; but we can study it objectively .

\begin{ExampleBox}
The cockfight illustration . Two people are engaged in a cockfight . I am a European watching them . Just by looking at the crowd, I may call it a source of entertainment---a show, some theatre .
But when I start looking at the cockfight by putting myself in the shoes of those people, and after going through the history of the whole community, I realise it is not just a cockfight . It is a symbol of power struggle---a struggle for power and for status between two people .
The meaning comes out very differently, and this is objective . So although the social world is highly subjective, the researcher can study that subjective world objectively if the researcher does not allow his values to enter the interpretation of social action .
\end{ExampleBox}

\paragraph{Step 3---The Lines to Write}
\begin{enumerate}
    \item Objectivity has been a notion derived from the natural sciences that claims the research to be value free .
    \item Positivists believed value-free analysis can be done with the application of the scientific method in the social world, which would help to discover laws of society . Example: Durkheim arrived at the law that with industrialization the suicide rate rises .
    \item Max Weber, on the other hand, believed that the social world consists of meaningful social actions which are motivated and subjective . Therefore objectivity can be attained only if the researcher studies the web of meanings in a value-neutral fashion .
    \item How can we study something value-neutrally? Through the concept of the \textbf{ideal type}, which helps the researcher study the social world objectively and helps maintain value neutrality .
\end{enumerate}

\paragraph{Step 4---The Ideal Type as the Device of Value Neutrality}
\vspace{6pt}
\noindent
\begin{tabularx}{\textwidth}{>{\hsize=0.25\hsize}Y >{\hsize=0.45\hsize}Y >{\hsize=0.3\hsize}Y}
\toprule
\rowcolor{AccentDark}
\thead{Scholar} & \thead{Treatment of Capitalism} & \thead{Is It a Value Judgment?} \\
\midrule
Adam Smith & A neoliberal economy gives rise to the wealth of nations . & Yes---a value judgment . \\
\midrule
Karl Marx & Capitalism is highly exploitative and alienates individuals . & Yes---a value judgment . \\
\midrule
Max Weber & Capitalism is the accumulation of wealth---accentuated features only . & No---he does not say whether wealth accumulation is good or bad, who exploits whom, or whether the economy grows or solidarity results . \\
\bottomrule
\end{tabularx}
\vspace{6pt}

In history there have been different methods of accumulating wealth . In tribal society it was booty capitalism---accumulating wealth by capturing or defeating other people . This is objective: it simply depicts what happened in history, with no judgment .
Contrast Marx: ``the history of all society has been the history of class struggle'' . Weber simply portrays that wealth was accumulated; he does not talk about exploitation, expropriation, appropriation, surplus value or alienation .
This is why it is a very unique method of objectivity in social science . And remember, ideal type also has types: historical, universal and historical-particular .

\paragraph{Step 5---Adequate Causality and the Claim to Science}
These ideal types helped Weber not only to maintain objectivity but also to draw adequate causal relations---to establish cause and effect in the social world . Cause and effect is also a feature of natural science, so there is science here .
Hence: sociology is a science---but a science of social action with interpretative understanding of the meanings and motives, along with the causes and the consequences .

\paragraph{Step 6---The Critics}
\begin{CriticalPoint}
Introduce scholars---advisable in a 20-mark answer, ignorable in a 10-mark one .
Gunnar Myrdal and Alvin Gouldner clearly say objectivity is an illusion . David Marsland may also be introduced .
Gouldner's concept of \textbf{domain assumptions}: if I am a feminist, I study domestic violence, gender inequality and sexual harassment at the workplace . My personal life, my values and my reading of feminism influenced me to study these things---those are my domain assumptions .
Similarly, as a Marxist scholar I would study class inequality, poverty, homelessness, misery, alienation, monotonous work, the platform economy, cyber and digital capitalism, migration, reverse migration and forced labour . Being a Marxist, you deal with those topics---so how can you say you are highly objective? 

Max Weber's answer: as far as the selection of the topic is concerned, your values are very much welcome; after the topic is chosen, no . The critics reply that even after the selection of the topic values will still be there .
\end{CriticalPoint}

\begin{KeyConcept}
So Weber's idea of objectivity is itself very ``ideal'': the ideal type gives us a direction, a target, a goal that everyone should try to attain . It is not something you can really live up to, but something you can always look up to while doing research .
\end{KeyConcept}

\begin{QuickRevision}
\begin{itemize}
    \item Keyword: objectivity . Objectivity is a natural-science notion of value-free research; positivists import it wholesale, Weber does not .
    \item For Weber the social world is subjective but can be studied objectively---the cockfight read as a struggle for power and status .
    \item The device is the ideal type, which secures value neutrality: capitalism as accumulation of wealth, with no judgment, unlike Adam Smith or Marx .
    \item Ideal types also permit adequate causal relations, so sociology is a science of social action with interpretative understanding of causes and consequences .
    \item Critics: Myrdal, Gouldner (objectivity is an illusion; domain assumptions) and Marsland; Weber allows values only at the stage of choosing the topic .
\end{itemize}
\end{QuickRevision}

\subsection{Weber's Objectivity versus Positivist Objectivity}
\vspace{6pt}
\noindent
\begin{tabularx}{\textwidth}{>{\hsize=0.22\hsize}Y >{\hsize=0.38\hsize}Y >{\hsize=0.4\hsize}Y}
\toprule
\rowcolor{AccentDark}
\thead{Element} & \thead{Weber} & \thead{Positivists} \\
\midrule
Type & Ideal type---not the actual type, nor the average type; it is an approximation of reality, not reality . & Actual type---what they study is taken to be the thing itself . \\
\midrule
Understanding & Verstehen---empathetic understanding; Weberians understand us . & They do not understand people, they explain people . \\
\midrule
Values & Value neutrality . & Value freedom, with strict separation between researcher and researched; prediction of reality . \\
\midrule
Causation & Adequate causality---probability, multi-causality . & Monocausal: A led to B . \\
\midrule
Assumption about people & A web of meanings; reality is highly subjective . & Just as they have the atom, we have people who are like atoms; behaviour is static, people in similar situations behave similarly everywhere, and therefore laws of the social world can be formulated . \\
\bottomrule
\end{tabularx}
\vspace{6pt}

\begin{KeyConcept}
Capitalism does not exist; it is a model . Democracy is a model, not a reality . Justice is a model---equal opportunity to everyone, free and fair elections as political justice, plus economic and social justice . In reality you do not have a hundred per cent of economic or social justice; you get one feature here and another there .
Every scholar uses concepts and holds models in mind to explain reality .
If nothing is certain, how can you come up with a law? That is Weber's position in one line .
\end{KeyConcept}

\begin{QuickRevision}
\begin{itemize}
    \item Weber: ideal type, verstehen, value neutrality, adequate causality, web of meanings .
    \item Positivists: actual type, explanation, value freedom and researcher/researched separation, monocausality, prediction, and laws .
    \item Ideal type is neither the actual nor the average type---only an approximation .
    \item Democracy, justice and capitalism are all models rather than realities, which is why no law can be formulated .
\end{itemize}
\end{QuickRevision}

\subsection{Parsons on Orientation, and His Five Questions}
Recall the diagram: an actor, arrows for the possible means, situational constraints below and normative constraints above . Parsons is studying the orientation of the actor---a word that comes from Weber's writings, since social action is meaningful action oriented to other actors .

\begin{ComparisonBox}
Weber: the actor is oriented to other actors .
Parsons: the actor is oriented to other actors, of course---but this orientation is not so neat and clean . There are constraints, normative and situational, which Weber did not take into consideration .
Weber tried to understand the social action of large numbers of people to figure out why large-scale social structures came into existence---his structures are the formation of individuals . That way structures are not objective reality, since they are our making, and all of us may attach different meanings to the same structure . To the concept of work a Protestant attributes one meaning and a Hindu another; so the very notion of capitalism is not objective---it is a reflection of the values of Protestants .
\end{ComparisonBox}

\paragraph{The Five Questions in Parsons' Mind}
\begin{enumerate}
    \item How does the actor finalize the means? 
    \item How are the goals or ultimate ends determined? 
    \item Why is the actor, despite having a different set of goals and means, able to interact with someone who does not share those goals and means? 
    \item How do the normative and situational constraints coerce the individual into certain means and certain goals? 
    \item \textbf{How come we have social order in society, with so many people having different sets of goals and different sets of means?} 
\end{enumerate}

\begin{ExampleBox}
Parsons is defamiliarizing what is very familiar . Two people talking is nothing extraordinary---someone asks the way to the temple, the other answers . But is this the kind of interaction Parsons is interested in? Yes: he is unearthing everything we take for granted .
UPSC has approximately 1000 seats and about 3 lakh students appear . Finally 1000 make it and 2,99,000 are out . Why did they simply accept that they are not fit? Why did they not revolt or protest? What is it in society that stops them? 
Or: suppose you marry someone completely opposite to you . After a year both of you know that this is not the person each thought the other to be . Why then is there still no divorce? Why is the marriage still there? Social order is still maintained despite tension, anxiety and strain .
\end{ExampleBox}

\begin{QuickRevision}
\begin{itemize}
    \item Parsons studies orientation, adding normative and situational constraints to Weber's orientation to other actors .
    \item His five questions: how means are finalized; how goals are determined; why actors with different goals still interact; how constraints coerce; and above all how social order is possible .
    \item The method is defamiliarization---treating the ordinary (an exam result accepted without protest, a marriage that does not end) as a problem requiring explanation .
\end{itemize}
\end{QuickRevision}

\subsection{Pattern Variables}
Whatever choice you make can be studied easily through a pattern, because the choices individuals make constitute patterns . Large numbers of people go through similar means to achieve similar goals . There is complexity, but there is a pattern---and Parsons calls this pattern the \textbf{pattern variable} .

\begin{KeyConcept}
\textbf{Definition:} a pattern variable is a model used by the researcher to understand how the actor resolves the dilemma that he confronts in an action situation .
Multiple choice means dilemma . Parents' choice or your own; UPSC or higher education or a foreign university; pizza or burger; hot coffee or cold coffee---every actor in any action situation confronts a dilemma, and finally makes one choice .
Since it is a model, it is very much influenced by Weber's notion of the ideal type .
\end{KeyConcept}

After analysing all the modes of social action, orientation, and normative and situational constraints, Parsons found there are two modes---two choices the actor has: to go with tradition, or to go with his own choice, that is individualism .
Corruption is a manifestation of individualism; to write the UPSC examination is considered a traditional social action, or it may be considered value-rational---whatever it may be, it is not individualism .

\paragraph{The Five Pattern Variables}
\vspace{6pt}
\noindent
\begin{tabularx}{\textwidth}{>{\hsize=0.23\hsize}Y >{\hsize=0.25\hsize}Y >{\hsize=0.52\hsize}Y}
\toprule
\rowcolor{AccentDark}
\thead{Pattern Variable A\newline(Traditional, Expressive)} & \thead{Pattern Variable B\newline(Modern, Instrumental)} & \thead{What It Means} \\
\midrule
Ascription & Achievement & Ascribed status---you are Punjabi, from MP, of a specific caste, of a specific religion . Achieved status---you are an MBA, an MA, an M.Tech . If you choose a partner by caste, region or religion, that is ascription; if by educational background and achievement, that is achievement orientation . In reality both are looked at, which is why it is only a model . \\
\midrule
Affectivity & Affective neutrality & At home your mother is emotional about you and you may get annoyed with her . The same son would not express that annoyance to his boss . In the office you cannot be emotional; if you have an issue you write it in a letter and email it, and the boss replies by the same email . Family = affectivity; work or organization = affective neutrality . \\
\midrule
Particularism & Universalism & At home my mother and family give me special treatment . In school no student gets special treatment---all students are equal to the teacher, we are roll numbers one to ten, not Sonu or Monu . In the public sphere we aim at universal standards, and universal standards are the feature of modern society: this is what democracy is about---equal opportunity to everyone . \\
\midrule
Diffuseness & Specificity & Your role at home is highly diffused: we do not talk to our mother and father only because we have some specific business; we talk to them generally . In the office people talk because they have work, and productive discussion happens through conferences and seminars . You cannot talk to the police the way you talk to your father; with a doctor the conversation is specific about treatment and medication; with the police about a specific law, violation or accident . \\
\midrule
Collective orientation & Self-orientation & In traditional society you always take the collectivity into consideration---your family, your nation, your religion, your caste . In modern society you are highly self-oriented . Examples: first-generation IFS officers who are not concerned about lineage or traditional occupation and simply want to serve India from outside it; a person with a family business who does not continue it and prepares for UPSC instead; a Dalit no longer engaged in a menial task and preparing for UPSC . When self-orientation spreads you see a change in the occupational structure and a breakdown in the structure of caste hierarchy---which shows modernity is coming . \\
\bottomrule
\end{tabularx}
\vspace{6pt}

\begin{DataFact}
On intercaste marriage: till last year, 93 per cent of marriages in India were intra-caste---within the same caste even after 75 years of independence . Hardly 5 or 6 per cent are intercaste .
Intra-caste marriage is not individualism; it is traditional action . Going out of your caste means making your own choice---individualism .
A definitional problem also arises: if a Mishra marries a Pandey, is that intercaste or intra-caste?  Generally we think both are traditionally Brahmins, so it is intra-caste endogamy---but the castes are in fact different, so it could be counted as intercaste . It depends on the assumptions the researcher works with .
\end{DataFact}

\begin{QuickRevision}
\begin{itemize}
    \item A pattern variable is a model for understanding how an actor resolves the dilemma of an action situation---dilemma meaning multiple choice .
    \item Two modes of choice: tradition (A) and individualism (B) .
    \item The five pairs: ascription/achievement, affectivity/affective neutrality, particularism/universalism, diffuseness/specificity, collective orientation/self-orientation .
    \item A describes traditional society, B modern society; India's 93\% intra-caste marriage rate is the class's test case, complicated by how ``intercaste'' is defined .
\end{itemize}
\end{QuickRevision}

\subsection{Significance of Pattern Variables and the Classification of Societies}
\begin{enumerate}
    \item Pattern variable as a model can help us to understand the social action of the individual inclusively, holistically and comprehensively .
    \item Pattern variable A and pattern variable B are two different modes of orientation---in brackets, expressive and instrumental . A is highly expressive; B is instrumental .
    \item Pattern variable B helps in understanding the larger social system: if the social action of individuals falls under pattern variable A, the social system is traditional; if under B, it is modern .
\end{enumerate}

Recall that Max Weber talked about instrumental social action---means and ends rationally calculated, everything done by yourself, all the means at your disposal, your own decision . In every one of the five pairs, the right-hand side is all about you as an individual; when you have so much of ``you'', you have instrumental orientation, and when you have so much of ``us'', it is expressive .

If out of one crore people I find eighty lakh doing what falls under pattern variable B, I will say society is moving towards modernity; if eighty per cent fall in the traditional domain, I would say India is still not modern---highly traditional with very few elements of modernity . But before any research one must have such models in mind before drawing conclusions .

\begin{CriticalPoint}
Whenever Parsons talks about social action he is considering only and only formal rationality, because his first assumption was that man is utilitarian and as a utilitarian man he is rational, has choices, and makes them accordingly . The other forms of rationality and the other forms of social action he does not talk about .
And even when people opt for what is expressive or falls under pattern variable A, he does not give importance to it---because to him being modern means pattern variable B, and A is traditional and non-modern . He looks down on A and gives all the importance to B .
\end{CriticalPoint}

\paragraph{How Each Scholar Classified Society}
\vspace{6pt}
\noindent
\begin{tabularx}{\textwidth}{>{\hsize=0.3\hsize}Y >{\hsize=0.7\hsize}Y}
\toprule
\rowcolor{AccentDark}
\thead{Scholar} & \thead{Basis of Classification} \\
\midrule
Durkheim & Division of labour and collective consciousness and law as well . \\
\midrule
Marx & Modes of production and classes . \\
\midrule
Weber & Rationality which includes authority and the forms of capitalism; essentially the way we think . \\
\midrule
Parsons & Pattern variables, that is, modes of orientation . \\
\bottomrule
\end{tabularx}
\vspace{6pt}

\begin{ExampleBox}
On Weber's ``way we think'': today you opt for a college because you consider it a brand, not because you like the subject . You are not following your preference or your passion . Everything is over---we are in a disenchanted world, highly rational . That is the irrationality of rationality .
The change in the way of thinking---from traditional social action to formal, highly rational social action---is what the rationalization process is all about . Today, before getting married, you would not be thinking about extravagant expenses; you would ask why not simply invest the money instead .
\end{ExampleBox}

\paragraph{Traditional and Modern Society by Scholar}
\vspace{6pt}
\noindent
\begin{tabularx}{\textwidth}{>{\hsize=0.2\hsize}Y >{\hsize=0.4\hsize}Y >{\hsize=0.4\hsize}Y}
\toprule
\rowcolor{AccentDark}
\thead{Scholar} & \thead{Traditional Society} & \thead{Modern Society} \\
\midrule
Weber & Traditional social action, traditional authority, practical rationality . & Zweckrational social action, legal-rational authority, formal rationality . \\
\midrule
Marx & Limited technological advancement, limited needs, community with no class structure . & Highest form of technological development; relations of production shaped by technological advancement, resulting in the formation of classes aligned asymmetrically, one class oppressed by another; false needs, oppression, exploitation, no more community . \\
\midrule
Parsons & Orientation falling under pattern variable A . & Orientation falling under pattern variable B . \\
\bottomrule
\end{tabularx}
\vspace{6pt}

\begin{QuickRevision}
\begin{itemize}
    \item Three significances: holistic understanding of individual social action; A and B as expressive and instrumental modes of orientation; and B as the key to classifying the larger social system .
    \item Parsons assumes a utilitarian, formally rational actor and treats A as merely traditional and non-modern .
    \item Classification bases: division of labour and collective consciousness (Durkheim), modes of production and classes (Marx), rationality (Weber), pattern variables (Parsons) .
\end{itemize}
\end{QuickRevision}

\subsection{Modes of Control and the Problem of Social Order}
Take one actor whose goal is to make money . Three legitimate choices---a startup, a government job, a private job---and a fourth, gambling . Parsons asks: why did you not go with gambling? 
You may answer: because if I gamble I can be sent to jail, so I have a fear of punishment, or my status will be defamed . But how did you develop this consciousness? Why are you opting for the means seen as legitimate in society? 

It was developed through three modes of control :
\begin{enumerate}
    \item \textbf{Sanctions}---that is, punishment . Deviance is not encouraged; it is prohibited by the fear of law . Law is one institution that maintains solidarity and prohibits deviance; you follow the consensus because of sanctions . In UPSC or any examination you will see the instruction that anyone using unfair means will be debarred---the sanction is given right at the page of instructions .
    \item \textbf{Values given by school and family}---the values that you have to be ambitious and aspirational . Right at school you start getting punished for something that is not seen as right for the discipline of the classroom . The classroom is a mini society: all the students have different aspirations and come from different families, yet all of them are disciplined, sitting silently and listening to the teacher, because of the values families have given them and because of the sanction represented by the monitor standing in the middle .
    \item \textbf{The larger social system}---societal values . These are values you do not really learn from school or family . For instance you only read about democracy in school; you do not practise or internalize it there . After school you go to college and then to a job, and in the job itself you learn how democracy actually operates---because your appointment, recruitment and selection are democratic, with everybody chosen on the basis of skill, opportunity and education, so you see equal opportunity irrespective of sex or gender .
\end{enumerate}

\begin{KeyConcept}
These are the three modes of control, and these are the things which maintain social order . Without disciplinary authority and without values, social order would not be possible .
\end{KeyConcept}

\paragraph{A Second Application: Entering Politics}
\vspace{6pt}
\noindent
\begin{tabularx}{\textwidth}{>{\hsize=0.45\hsize}Y >{\hsize=0.55\hsize}Y}
\toprule
\rowcolor{AccentDark}
\thead{Means to Enter Politics} & \thead{What Constrains It} \\
\midrule
Wealth---you get a ticket from a political party because you are a millionaire or billionaire, or because you have been sponsoring political parties & Available only to those with the wealth . \\
\midrule
Contacts---as an IAS or IPS officer you may have built up many contacts & Requires the prior position . \\
\midrule
Marrying the son or daughter of a President, Chief Minister or Prime Minister & The choice finally taken in the example---but itself shaped by normative constraints, since as a bureaucrat you cannot afford to offend the minister . \\
\midrule
Criminal behaviour---what we call the criminalization of politics, since politics is about majority, not about moral or immoral & Blocked by sanctions, values and the larger social system: the Supreme Court and the Election Commission act as watchdogs on the background of candidates and may not even allow you to contest . \\
\bottomrule
\end{tabularx}
\vspace{6pt}

\begin{CriticalPoint}
The larger point: in the theory of the structure of social action, Talcott Parsons is always and only talking about formal rationality---you have multiple means before you, and you opt for certain means by eliminating others, considering these right and those wrong .
\end{CriticalPoint}

\begin{QuickRevision}
\begin{itemize}
    \item Three modes of control: sanctions (fear of law and punishment), values instilled by school and family, and the values of the larger social system .
    \item Together they explain why deviant means---gambling, criminality---are eliminated, and thus how social order is possible .
    \item The classroom is a mini society in which sanction and value operate simultaneously .
    \item Every such elimination of means is an exercise of formal rationality, which is the only rationality Parsons works with .
\end{itemize}
\end{QuickRevision}

\subsection{The Social System and the Four Needs}
The human body has different needs---to eat, to drink, to meet and greet someone, to sleep, to rest . Just as the human body has different needs, society, if considered as an organism---which is what functionalists have always maintained---also has different needs .

\begin{KeyConcept}
The needs of society are four only:
\begin{itemize}
    \item Adaptation
    \item Goal attainment
    \item Integration
    \item Latency
\end{itemize}
What makes Talcott Parsons difficult is only one thing---the words he uses . Otherwise he is a very simple scholar .
\end{KeyConcept}

\paragraph{The Shimla Illustration}
\vspace{6pt}
\noindent
\begin{tabularx}{\textwidth}{>{\hsize=0.25\hsize}Y >{\hsize=0.75\hsize}Y}
\toprule
\rowcolor{AccentDark}
\thead{Need} & \thead{In the Shimla Trip} \\
\midrule
Adaptation & You land in Shimla and first take rest---you are at high altitude and feeling breathless . You adapt to the external environment . \\
\midrule
Goal attainment & Then you decide to have a nice dinner, or to go out for a long walk . After adapting, you go on to fulfil your goals . \\
\midrule
Integration & You go out and meet people; you eat the ordinary local food comfortably; you talk to friends and to the local people of Shimla and get photographs with them; you sit with the taxi driver and talk to him; you integrate with the hotel staff . Unless you integrate with people, people would not cooperate with you . \\
\midrule
Latency & In the whole process of discussion, integration and fulfilling goals, tension also builds up---you come home and find the balance in your account is nothing like what you assumed . Latency is tension management . \\
\bottomrule
\end{tabularx}
\vspace{6pt}

\paragraph{The Same Applied to Indian Society after 1990}
\begin{itemize}
    \item \textbf{Adaptation}---when liberalization, privatization and globalization happened in the 1990s, Indian society had to adapt to these external changes . There is an institution in society which adapts to the external environment .
    \item \textbf{Goal attainment}---after LPG the goals of our society also changed .
    \item \textbf{Integration}---changes came in the organs of our society: changes in laws, changes in the economy, changes in the way law, economy, state, school and family were integrated with each other as organs . After the 1990s, the school curriculum changed and NCERT textbooks changed: chapters on liberalization, privatization, capitalism and liberalism started appearing . The education structure changed . The family as an institution changed---where there were joint families, aspirations arose and people moved to Delhi and Mumbai, so families went from joint to nuclear . Laws changed, the state changed, everything changed---and yet they remained in integration with each other .
    \item \textbf{Latency}---with so many changes, tension builds up . Strains develop in the system, just as strains develop in the body or the mental system . Therefore we also have a tension management mechanism .
\end{itemize}

\begin{KeyConcept}
Unless we have four institutions catering to the four different needs of society, we will never be able to maintain social order .
For social order to maintain itself, the four different needs have to be fulfilled by four different institutions of society . Which institutions serve which need is the next topic .
\end{KeyConcept}

\begin{QuickRevision}
\begin{itemize}
    \item Society, treated as an organism, has four needs: adaptation, goal attainment, integration and latency .
    \item The Shimla trip illustrates all four: rest, dinner and walks, meeting people, and managing the tension that builds up .
    \item Post-1990 India illustrates them structurally: adapting to LPG, changing goals, reintegrating law, economy, state, school and family, and absorbing the resulting strains .
    \item Social order requires four institutions, each catering to one need .
\end{itemize}
\end{QuickRevision}

\subsection{Doubts Answered at the End of Class}
\paragraph{What is the difference between situational and normative constraints and the modes of control?}
The modes of control include these constraints; they are sanctions .
\begin{itemize}
    \item \textbf{Normative constraint}---you want to marry someone and are not able to, because he or she belongs to a different caste or religion . Or the normative constraint may be absent on caste and religion and appear elsewhere: she is from your caste and your religion but poor, and dowry cannot be paid . (Dowry is of course a bad institution and has been prohibited; the point is only what happens in practice.) 
    \item \textbf{Situational constraint}---suppose you do not have the money to arrange the marriage, or you are not well enough situated economically to take care of the next generation . That way situational constraint is not about norms; it is about your income, geography and capability . Or: your wife is an IAS officer and so are you, but you are in Rajasthan and she is in UP, and ultimately it ends in divorce . That is not about norms; it is about geographical location, income and capability .
\end{itemize}

\paragraph{Parsons was influenced by the idealism school of thought---is it the same idealism?}
Idealism has multiple variants . Idealism simply means ideas are primary and matter is secondary . Remember this .

Parsons is influenced by three schools . From utilitarianism he takes that individuals have choice and are rational . From idealism he takes that norms will constrain the individual, that ideas are primary: the means you choose to fulfil your goals will be influenced by ideas, and the goal you set in your life is also influenced by ideas .

So the actor may be rational, in the sense that he can make all the choices by himself; but the means available before him would be governed by norms, and the goal he sets for himself would also be governed by norms . Both are there: the individual actor making his choice, and the norms constraining his action . And the third influence is positivism .

\begin{QuickRevision}
\begin{itemize}
    \item The modes of control include the constraints; sanctions are their core .
    \item Normative constraints concern caste, religion, dowry and other norms; situational constraints concern income, geography and capability---including two spouses posted in different states .
    \item Idealism = ideas are primary, matter secondary; through it norms shape both the means and the goals .
    \item Parsons' three influences remain utilitarianism, idealism and positivism .
\end{itemize}
\end{QuickRevision}

\sectiondivider

\section{Answer-Writing Toolkit}
The following collects only the answer-writing guidance actually given across these lectures .

\subsection{How to Approach the Question}
\begin{itemize}
    \item Do not read the question as a sentence---look at the keywords . Everything in the answer is built from them .
    \item First get uncomfortable by looking at the question, then get comfortable by finding the part you can begin from .
    \item ``Examine'' is not an important direction word---both sides have to be written . Even for ``elucidate'' the other side must appear .
    \item Write what is asked . Do not write the history of the concepts, the evolution of the idea, or the origin of sociology, when the question does not demand it .
    \item The way to be different is to address the question . Most candidates never address it and then ask how to be different .
    \item Everyone studies from similar notes and hears similar lectures; what separates the marks is structure .
\end{itemize}

\subsection{Length, Time and the Use of Scholars}
\begin{itemize}
    \item A 10-mark question is 150 words, and in class 7 minutes were given for it . Writing 10 marks is harder than writing 20 .
    \item In a 20-mark answer you should mention scholars for the critique; in a 10-mark answer you may ignore them .
    \item Criticism is not required in a 10-mark question on a definitional demand .
\end{itemize}

\subsection{Dimensions and Structure for the Two Questions Practised}

\paragraph{Which concept did Weber use to analyse the forms of legitimate domination? (2015, 10 marks)}
\begin{enumerate}
    \item Define domination---the probability of commands being obeyed; note that it is not power, and that authority is Parsons' translation of the term .
    \item Explain legitimacy and the three bases: tradition or ancient texts, extraordinary/charismatic qualities, and formal rules and regulations .
    \item Give the forms with quick examples---obeying the father, following a motivational speaker, stopping for a uniformed traffic policeman .
    \item Name the concepts used: social action, ideal type, domination/authority, adequate causality .
    \item Optional closing line: interpretative sociology was Weber's vision, and it helped sociology depart from the macro perspective to the micro .
\end{enumerate}

\paragraph{Examine Max Weber's method of maintaining objectivity in social research (2016)}
\begin{enumerate}
    \item Define objectivity as a notion derived from the natural sciences claiming research to be value free .
    \item State the positivist position: value-free analysis through the scientific method, to discover laws of society---with Durkheim's law relating industrialization to suicide rates as the example .
    \item State Weber's position: the social world consists of meaningful social actions which are motivated and subjective, so objectivity is attained only if the researcher studies the web of meanings in a value-neutral fashion .
    \item Show how the ideal type secures value neutrality---capitalism as accumulation of wealth, against the value judgments of Adam Smith and Marx; mention that ideal type has types: historical, universal, historical-particular .
    \item Add adequate causality, and conclude that sociology is a science of social action with interpretative understanding of meanings, motives, causes and consequences .
    \item Close with critics: Gunnar Myrdal and Alvin Gouldner (objectivity is an illusion; domain assumptions), and David Marsland .
\end{enumerate}

\subsection{Keywords to Carry Into the Answer}
\begin{itemize}
    \item \textbf{Weber on objectivity:} ideal type, value neutrality, verstehen, adequate causality .
    \item \textbf{Weber on authority:} domination, legitimacy, ideal type, social action, adequate causality .
    \item \textbf{Traditional authority:} gerontocracy, primary patriarchalism, patrimonialism (prebendalism), sultanism, patrimonial bureaucracy .
    \item \textbf{Charismatic authority:} routinization of charisma, historic and transitory phenomenon, hero or tyrant .
    \item \textbf{Bureaucracy:} machine-like efficiency, iron cage of rationality, irrationality of rationality, dictatorship of the officials .
    \item \textbf{Protestant Ethics:} predestination, salvation anxiety, calling, ascetic materialism, this-worldly orientation, Puritan work ethic, elective affinity, adequate causality, unintended consequence .
    \item \textbf{Parsons:} unit act, normative constraint, situational constraint, pattern variable, expressive and instrumental orientation, modes of control, adaptation-goal attainment-integration-latency .
\end{itemize}

\subsection{Specific Instructions Given for Particular Topics}
\begin{itemize}
    \item \textbf{On the Protestant Ethic question:} explain the ideal type of Protestantism (the Protestant work ethic) and the ideal type of rational capitalism, and draw the relation between the two . Do not narrate the history of Martin Luther, John Calvin or Benjamin Franklin in the answer---that is only for your understanding, and it will help in other units .
    \item \textbf{On historical method:} when you see a question on historical method, write only Max Weber, not Karl Marx . Marx's theory is named historical materialism, but his method is dialectics .
    \item The history-versus-sociology distinction developed while teaching the Reformation can itself be used in a question on history versus sociology .
\end{itemize}

\subsection{How to Read for the Examination (Stated in the Second Lecture)}
\begin{itemize}
    \item Read instrumentally: previous year questions first, then the weightage of areas, then the minimum book that gives maximum result in minimum time .
    \item \textbf{Ancient history:} Harappa/IVC, Buddhism, Jainism, Mauryan Empire, Guptas, scientific achievements, and the philosophies (Shad Darshanas) . Skip the Greeks, Shakas, Parthians, Kushans and the Satavahanas . Three days maximum .
    \item \textbf{Medieval history:} only the Mughals, the Cholas, the Pallavas, and Bhakti and Sufi . No Sultanate .
    \item \textbf{Modern history:} skip the early modern period before 1857; read the major events along with the major acts, always alongside previous year questions . Do not try to become a scholar of modern history .
    \item \textbf{Art and culture:} NCERT only---Buddhist art, paintings, Indo-Islamic architecture, folk dances and classical dances .
    \item Whatever you find boring should be read first, because otherwise you will never read it .
\end{itemize}

\subsection{Applying the Syllabus to the Newspaper}
A newspaper article should be read twice: once as a common-sensical reader would, and once as a student of sociology . The class demonstration used two \textit{Indian Express} pieces on examination culture and on instant spirituality, both read through formal rationality, the fetish of quantification, social Darwinism and the iron cage of rationality .